\documentclass[fs85,seriesRus]{korfarm-base}
\usepackage{korfarm-russell}

% ============================================================
% passage-note.tex (v11)
% 지문 박스 + 첫 페이지 우측 메모란 (동적 높이)
% 정책: PAGE_BREAK_POLICY.md §3, §4 참조
%
% 변경 (v11):
%  · 지문 박스 폭 확대: enlarge right by=-3.7cm  (메모 3.4cm + gap 3mm)
%  · 메모란 동적 높이: frame.north east ~ frame.south east 사이 자동 채움
%  · 메모 줄선: clip 영역으로 박스 내부만 채움 (박스 높이 모르고도 작동)
%  · 문단 들여쓰기 활성화: tcolorbox 내부 \parindent=1em, \noindent 제거
%
% 사용:
%   \kfPassageNote{제목}{지문 본문}      → 메모란 포함
%   \kfPassageNoteNT{지문 본문}          → 메모란 포함, 제목 없음
%   \kfPassageFull{제목}{본문}           → 메모란 없음 (활동 내 보기)
% ============================================================

\makeatletter

% 메모란 규격 (cm 고정)
\newlength{\kf@memoW}
\setlength{\kf@memoW}{3.4cm}
\newlength{\kf@memoGap}
\setlength{\kf@memoGap}{3mm}

% --- 메인: 제목 있는 버전 ---
\newcommand{\kf@passagenote@with}[2]{%
  \par\ifkfPassageNewPage\ifdim\pagetotal>0.4\textheight\clearpage\fi\fi\smallskip\needspace{6\baselineskip}%
  \begin{tcolorbox}[
    enhanced, breakable,
    width=\dimexpr\linewidth-4cm\relax,
    colback=kfPassageBg, colframe=kfPrimary!70,
    boxrule=0.6pt, arc=3pt,
    left=\kfBoxPad, right=\kfBoxPad, top=\dimexpr\kfBoxPad+8pt\relax, bottom=\kfBoxPad,
    title={\small\bfseries\color{kfPrimary} #1},
    coltitle=kfPrimary,
    colbacktitle=kfPrimaryLight!60,
    attach boxed title to top left={yshift=-2.4mm, xshift=8pt},
    boxed title style={colframe=kfPrimary!50, boxrule=0.5pt, arc=2pt},
    parbox=false,
    overlay unbroken and first={%
      \begin{scope}
        \draw[kfRule, line width=0.4pt, fill=kfNoteBg, rounded corners=2pt]
          ([xshift=\kf@memoGap]frame.north east) rectangle
          ([xshift=\kf@memoGap+\kf@memoW]frame.south east);
        \node[anchor=north east, font=\footnotesize\bfseries\color{kfMute}, inner sep=3pt]
          at ([xshift=\kf@memoGap+\kf@memoW, yshift=-2pt]frame.north east) {메모};
        \node[anchor=north east, fill=kfPrimary!15, text=kfPrimary,
              font=\scriptsize\bfseries, rounded corners=2pt,
              inner xsep=4pt, inner ysep=2pt]
          at ([xshift=\kf@memoGap+\kf@memoW-3pt, yshift=-19pt]frame.north east) {\kf@memocat};
        \begin{scope}
          \path[clip]
            ([xshift=\kf@memoGap]frame.north east) rectangle
            ([xshift=\kf@memoGap+\kf@memoW]frame.south east);
          \foreach \i in {1,...,40}{%
            \draw[kfRule, line width=0.25pt]
              ([xshift=\kf@memoGap+2mm, yshift=-7mm - \i*5mm]frame.north east) --
              ([xshift=\kf@memoGap+\kf@memoW-2mm, yshift=-7mm - \i*5mm]frame.north east);%
          }
        \end{scope}
      \end{scope}
    }
  ]
    \setlength{\parindent}{1em}%
    \hspace*{1em}\ignorespaces#2%
  \end{tcolorbox}\par\smallskip
}

% --- 제목 없는 버전 ---
\newcommand{\kf@passagenote@notitle}[1]{%
  \par\ifkfPassageNewPage\ifdim\pagetotal>0.4\textheight\clearpage\fi\fi\smallskip\needspace{6\baselineskip}%
  \begin{tcolorbox}[
    enhanced, breakable,
    width=\dimexpr\linewidth-4cm\relax,
    colback=kfPassageBg, colframe=kfPrimary!70,
    boxrule=0.6pt, arc=3pt,
    left=\kfBoxPad, right=\kfBoxPad, top=\dimexpr\kfBoxPad+8pt\relax, bottom=\kfBoxPad,
    parbox=false,
    overlay unbroken and first={%
      \begin{scope}
        \draw[kfRule, line width=0.4pt, fill=kfNoteBg, rounded corners=2pt]
          ([xshift=\kf@memoGap]frame.north east) rectangle
          ([xshift=\kf@memoGap+\kf@memoW]frame.south east);
        \node[anchor=north east, font=\footnotesize\bfseries\color{kfMute}, inner sep=3pt]
          at ([xshift=\kf@memoGap+\kf@memoW, yshift=-2pt]frame.north east) {메모};
        \node[anchor=north east, fill=kfPrimary!15, text=kfPrimary,
              font=\scriptsize\bfseries, rounded corners=2pt,
              inner xsep=4pt, inner ysep=2pt]
          at ([xshift=\kf@memoGap+\kf@memoW-3pt, yshift=-19pt]frame.north east) {\kf@memocat};
        \begin{scope}
          \path[clip]
            ([xshift=\kf@memoGap]frame.north east) rectangle
            ([xshift=\kf@memoGap+\kf@memoW]frame.south east);
          \foreach \i in {1,...,40}{%
            \draw[kfRule, line width=0.25pt]
              ([xshift=\kf@memoGap+2mm, yshift=-7mm - \i*5mm]frame.north east) --
              ([xshift=\kf@memoGap+\kf@memoW-2mm, yshift=-7mm - \i*5mm]frame.north east);%
          }
        \end{scope}
      \end{scope}
    }
  ]
    \setlength{\parindent}{1em}%
    \hspace*{1em}\ignorespaces#1%
  \end{tcolorbox}\par\smallskip
}

\newcommand{\kfPassageNote}[2]{%
  \def\kf@tmptitle{#1}%
  \ifx\kf@tmptitle\@empty
    \kf@passagenote@notitle{#2}%
  \else
    \kf@passagenote@with{#1}{#2}%
  \fi
}
\newcommand{\kfPassageNoteNT}[1]{\kf@passagenote@notitle{#1}}
\makeatother

% ============================================================
% 풀폭 지문 박스 (메모란 없음) — 활동 내 보기 박스
% PAGE_BREAK_POLICY.md §3
% ============================================================
\makeatletter
\newcommand{\kf@passagefull@with}[2]{%
  \par\ifkfPassageNewPage\ifdim\pagetotal>0.4\textheight\clearpage\fi\fi\smallskip\needspace{6\baselineskip}%
  \begin{tcolorbox}[
    enhanced, breakable,
    colback=kfPassageBg, colframe=kfPrimary!70,
    boxrule=0.6pt, arc=3pt,
    left=\kfBoxPad, right=\kfBoxPad, top=\dimexpr\kfBoxPad+8pt\relax, bottom=\kfBoxPad,
    title={\small\bfseries\color{kfPrimary} #1},
    coltitle=kfPrimary,
    colbacktitle=kfPrimaryLight!60,
    attach boxed title to top left={yshift=-2.4mm, xshift=8pt},
    boxed title style={colframe=kfPrimary!50, boxrule=0.5pt, arc=2pt},
    parbox=false
  ]
    \setlength{\parindent}{1em}%
    \hspace*{1em}\ignorespaces#2%
  \end{tcolorbox}\par\smallskip
}
\newcommand{\kf@passagefull@notitle}[1]{%
  \par\ifkfPassageNewPage\ifdim\pagetotal>0.4\textheight\clearpage\fi\fi\smallskip\needspace{6\baselineskip}%
  \begin{tcolorbox}[
    enhanced, breakable,
    colback=kfPassageBg, colframe=kfPrimary!70,
    boxrule=0.6pt, arc=3pt,
    left=\kfBoxPad, right=\kfBoxPad, top=\dimexpr\kfBoxPad+8pt\relax, bottom=\kfBoxPad,
    parbox=false
  ]
    \setlength{\parindent}{1em}%
    \hspace*{1em}\ignorespaces#1%
  \end{tcolorbox}\par\smallskip
}
\newcommand{\kfPassageFull}[2]{%
  \def\kf@tmptitle{#1}%
  \ifx\kf@tmptitle\@empty
    \kf@passagefull@notitle{#2}%
  \else
    \kf@passagefull@with{#1}{#2}%
  \fi
}
\newcommand{\kfPassageFullNT}[1]{\kf@passagefull@notitle{#1}}
\makeatother

% ============================================================
% 작은 문장 박스 (문장 독해용) — PAGE_BREAK_POLICY.md §2.5
% ============================================================
\newcommand{\kfSentenceBox}[2]{%
  \par\smallskip\needspace{3\baselineskip}%
  \noindent\begin{tcolorbox}[
    enhanced, breakable=false,
    colback=kfPassageBg, colframe=kfMute,
    boxrule=0.4pt, arc=2pt,
    left=8pt, right=6pt, top=4pt, bottom=4pt,
    title={\footnotesize\bfseries\color{kfPrimary} #1},
    coltitle=kfPrimary, colbacktitle=kfPaper,
    attach boxed title to top left={yshift=-1.6mm, xshift=4pt},
    boxed title style={colframe=kfMute, boxrule=0.3pt, arc=1pt}
  ]%
    \setstretch{1.3}#2%
  \end{tcolorbox}\par\smallskip%
}

% ============================================================
% 지문 본문 아래 안내/정리 박스 (v16 신규)
% 사용: \kfPassageHint{한 줄 정리}{본문 안내 텍스트}
% ============================================================
\newcommand{\kfPassageHint}[2]{%
  \par\smallskip\needspace{4\baselineskip}%
  \begin{tcolorbox}[
    enhanced, breakable=false,
    colback=kfPrimary!5, colframe=kfPrimary!40,
    boxrule=0.4pt, arc=2pt,
    left=8pt, right=6pt, top=4pt, bottom=4pt,
    title={\footnotesize\bfseries\color{kfPrimary} #1},
    coltitle=kfPrimary, colbacktitle=kfPaper,
    attach boxed title to top left={yshift=-1.5mm, xshift=4pt},
    boxed title style={colframe=kfPrimary!40, boxrule=0.3pt, arc=1pt}
  ]%
    \small\setstretch{1.3}#2%
  \end{tcolorbox}\par\smallskip%
}

% ============================================================
% 환경 버전 (호환성) — 메모란 포함
% ============================================================
\makeatletter
\newenvironment{kfPassageNoteEnv}[1]%
  {\par\needspace{6\baselineskip}%
   \def\kf@psrc{#1}%
   \begin{tcolorbox}[
     enhanced, breakable,
     width=\dimexpr\linewidth-4cm\relax,
     colback=kfPassageBg, colframe=kfPrimary!70,
     boxrule=0.6pt, arc=3pt,
     left=\kfBoxPad, right=\kfBoxPad, top=\dimexpr\kfBoxPad+8pt\relax, bottom=\kfBoxPad,
     title={\small\bfseries\color{kfPrimary} \kf@psrc},
     coltitle=kfPrimary, colbacktitle=kfPrimaryLight!60,
     attach boxed title to top left={yshift=-2.4mm, xshift=8pt},
     boxed title style={colframe=kfPrimary!50, boxrule=0.5pt, arc=2pt},
     parbox=false,
     overlay unbroken and first={%
       \begin{scope}
         \draw[kfRule, line width=0.4pt, fill=kfNoteBg, rounded corners=2pt]
           ([xshift=\kf@memoGap]frame.north east) rectangle
           ([xshift=\kf@memoGap+\kf@memoW]frame.south east);
         \node[anchor=north east, font=\footnotesize\bfseries\color{kfMute}, inner sep=3pt]
           at ([xshift=\kf@memoGap+\kf@memoW, yshift=-2pt]frame.north east) {메모};
         \begin{scope}
           \path[clip]
             ([xshift=\kf@memoGap]frame.north east) rectangle
             ([xshift=\kf@memoGap+\kf@memoW]frame.south east);
           \foreach \i in {1,...,40}{%
             \draw[kfRule, line width=0.25pt]
               ([xshift=\kf@memoGap+2mm, yshift=-7mm - \i*5mm]frame.north east) --
               ([xshift=\kf@memoGap+\kf@memoW-2mm, yshift=-7mm - \i*5mm]frame.north east);%
           }
         \end{scope}
       \end{scope}
     }
   ]%
   \setlength{\parindent}{1em}%
  }%
  {\end{tcolorbox}\par\smallskip}
\makeatother

% ============================================================
% condition-box.tex
% <조건> 박스 컴포넌트 (tcolorbox)
% 사용:
%   \begin{kfCondition}
%     \item 첫째 조건
%     \item 둘째 조건
%   \end{kfCondition}
% ============================================================

\newenvironment{kfCondition}%
  {\par\smallskip\needspace{4\baselineskip}%
   \begin{tcolorbox}[
     enhanced, breakable=false,
     colback=kfPrimaryLight!40, colframe=kfPrimary,
     boxrule=0.5pt, arc=2pt,
     left=\kfBoxPad, right=\kfBoxPad, top=\kfBoxPad, bottom=\kfBoxPad,
     title={\small\bfseries\color{kfPrimary} \,조건\,},
     coltitle=kfPrimary, colbacktitle=kfPaper,
     attach boxed title to top left={yshift=-2mm, xshift=6pt},
     boxed title style={colframe=kfPrimary, boxrule=0.4pt, arc=1pt}
   ]%
   \begin{enumerate}[leftmargin=16pt, itemsep=1pt, topsep=1pt,
                    label=\textbf{\arabic*.}]}%
  {\end{enumerate}\end{tcolorbox}\par\smallskip}

% ============================================================
% evidence-box.tex
% <보기> 박스 컴포넌트
% 사용:
%   \begin{kfEvidence}
%     보기 본문 ...
%   \end{kfEvidence}
% ============================================================

\newenvironment{kfEvidence}%
  {\par\smallskip\needspace{4\baselineskip}%
   \begin{tcolorbox}[
     enhanced, breakable=false,
     colback=kfPaper, colframe=kfMute,
     boxrule=0.5pt, arc=2pt,
     left=\kfBoxPad, right=\kfBoxPad, top=\kfBoxPad, bottom=\kfBoxPad,
     title={\small\bfseries\color{kfInk} \,보기\,},
     coltitle=kfInk, colbacktitle=kfPrimaryLight!50,
     attach boxed title to top left={yshift=-2mm, xshift=6pt},
     boxed title style={colframe=kfMute, boxrule=0.4pt, arc=1pt}
   ]%
   \leavevmode\mbox{}\ignorespaces%
  }%
  {\end{tcolorbox}\par\smallskip}

% ============================================================
% choice-list.tex
% 5선지 (① ② ③ ④ ⑤) 자동 들여쓰기 컴포넌트
% 사용:
%   \begin{kfChoices}
%     \kfChoice 첫째 선택지
%     \kfChoice 둘째 선택지
%     \kfChoice 셋째 선택지
%     \kfChoice 넷째 선택지
%     \kfChoice 다섯째 선택지
%   \end{kfChoices}
% ============================================================

\newcounter{kfChoiceCnt}
\newcommand{\kfChoiceMark}[1]{%
  \ifcase#1\relax%
  \or ①\or ②\or ③\or ④\or ⑤\or ⑥\or ⑦\or ⑧\or ⑨%
  \fi
}

% 환경: 자동 번호
\newenvironment{kfChoices}%
  {\setcounter{kfChoiceCnt}{0}%
   \par\smallskip%
   \begin{list}{}{%
     \setlength{\leftmargin}{20pt}%
     \setlength{\itemindent}{0pt}%
     \setlength{\labelwidth}{18pt}%
     \setlength{\labelsep}{6pt}%
     \setlength{\itemsep}{2pt}%
     \setlength{\topsep}{1pt}%
     \setlength{\parsep}{0pt}%
     \renewcommand{\makelabel}[1]{##1\hss}%
   }}%
  {\end{list}\par\smallskip}

\newcommand{\kfChoice}{%
  \stepcounter{kfChoiceCnt}%
  \item[\textbf{\kfChoiceMark{\value{kfChoiceCnt}}}]\ignorespaces%
}

% --- 한 줄 5선지 (짧은 선택지용) ---
\newcommand{\kfChoicesInline}[5]{%
  \par\smallskip%
  \noindent\textbf{①}~#1\quad\textbf{②}~#2\quad\textbf{③}~#3\quad\textbf{④}~#4\quad\textbf{⑤}~#5%
  \par\smallskip%
}

% ============================================================
% answer-note.tex
% 답안 작성란 (노트 스타일, 줄친 영역)
% 사용:
%   \kfAnswerLines{5}        % 5줄 답안란
%   \kfAnswerNote{서술형 답안란}{8} % 라벨+8줄
% ============================================================

% 단일 줄 (필요 시 인라인)
\newcommand{\kfAnswerLine}{%
  \par\noindent\textcolor{kfRule}{\rule{\linewidth}{0.4pt}}\par
}

% N줄 답안란 (간단)
\newcommand{\kfAnswerLines}[1]{%
  \par\smallskip%
  \begin{minipage}{\linewidth}%
  \foreach \i in {1,...,#1}{%
    \textcolor{kfRule}{\rule{\linewidth}{0.4pt}}\\[6pt]%
  }%
  \end{minipage}%
  \par\smallskip%
}

% 라벨이 있는 노트형 답안란
\newcommand{\kfAnswerNote}[2]{%
  \par\smallskip\needspace{#2\baselineskip}%
  \begin{tcolorbox}[
    enhanced, breakable=false,
    colback=kfPaper, colframe=kfRule,
    boxrule=0.4pt, arc=2pt,
    left=\kfBoxPad, right=\kfBoxPad, top=\kfBoxPad, bottom=\kfBoxPad,
    title={\footnotesize\bfseries\color{kfMute} #1},
    coltitle=kfMute, colbacktitle=kfPaper,
    attach boxed title to top left={yshift=-1.5mm, xshift=6pt},
    boxed title style={colframe=kfRule, boxrule=0.3pt, arc=1pt}
  ]%
  \foreach \i in {1,...,#2}{%
    \textcolor{kfRule}{\rule{\linewidth}{0.3pt}}\\[6pt]%
  }%
  \end{tcolorbox}\par\smallskip%
}

% 짧은 빈칸 (단답형)
\newcommand{\kfBlank}[1][3cm]{\underline{\hspace{#1}}}
% 넓은 빈칸 (서술 한 줄)
\newcommand{\kfBlankWide}[1][6cm]{\underline{\hspace{#1}}}
% 괄호 빈칸
\newcommand{\kfBlankParen}{(\,\hspace{1cm}\,)}
% O/X 빈칸
\newcommand{\kfBlankOX}{(\,\hspace{1.5cm}\,)}

% ============================================================
% question-block.tex
% 문제 블록 (번호 배지 + 발문 + 보기/조건 + 선택지 + 답안란)
% 모든 구성을 하나의 minipage로 묶어 단/페이지 단절 금지
%
% 사용:
%   \begin{kfQBlock}{1}{객관식}
%     \kfQStem{다음 글에 대한 설명으로 가장 적절한 것은?}
%     \begin{kfEvidence} ... \end{kfEvidence}    % 선택
%     \begin{kfCondition} ... \end{kfCondition}  % 선택
%     \begin{kfChoices}
%       \kfChoice ...
%     \end{kfChoices}
%     \kfAnswerLines{2}                            % 선택
%   \end{kfQBlock}
% ============================================================

% 무단절 minipage로 묶고, 발문/보기/조건/선택지/답안란을 자유롭게 배치
% 사용자 피드백 #4: [객관식]/[서술형] 등 텍스트 라벨 제거 — 번호 배지만 표시
% #2(유형 라벨)는 호환성 인자로만 받고 출력하지 않음.
\newenvironment{kfQBlock}[2]%
  {\par\smallskip\needspace{6\baselineskip}%
   \noindent\begin{minipage}{\linewidth}%
   \samepage%
   \par\noindent
   \kfQNum{#1}\hspace{4pt}%
   \begin{minipage}[t]{\dimexpr\linewidth-30pt\relax}%
  }%
  {\end{minipage}\end{minipage}\par\smallskip}

% 문제 발문
\newcommand{\kfQStem}[1]{%
  \par\noindent#1\par\vspace{2pt}%
}

% 짧은 소문항 라벨 (1), (2), (3)
\newcounter{kfSubQ}
\newcommand{\kfSubQReset}{\setcounter{kfSubQ}{0}}
\newcommand{\kfSubQ}{%
  \stepcounter{kfSubQ}%
  \par\noindent\textbf{(\arabic{kfSubQ})}~\ignorespaces%
}

% ============================================================
% activity-table.tex
% 활동: 정리표 / 내용 확인표 (마크다운 표 → tabularx 변환)
% 사용:
%   % 일반 tabular (직접 컬럼 명시)
%   \begin{kfActTable}{|l|X|X|}
%     \kfTHead{항목} & \kfTHead{설명} & \kfTHead{학생 작성} \\\hline
%     ① & ... & \kfTBlank \\\hline
%   \end{kfActTable}
%
%   % adjustbox 자동 축소 (정해진 폭으로 강제)
%   \kfActTableWrap{
%     ... (\begin{tabular}...\end{tabular} 코드) ...
%   }
% ============================================================

% 사용 시 본문 마지막 행은 \\\hline 으로 끝나야 함.
% v17.1: 흰색 mask로 Canva 배경 본문 침범 차단
\newenvironment{kfActTable}[1]%
  {\par\smallskip\needspace{4\baselineskip}%
   \renewcommand{\arraystretch}{1.3}%
   \arrayrulecolor{kfRule}%
   \begin{tcolorbox}[blanker, colback=white, left=2pt, right=2pt, top=2pt, bottom=2pt]%
   \noindent\begin{adjustbox}{max width=\linewidth}%
   \begin{tabular}{#1}%
   \hline}%
  {\end{tabular}%
   \end{adjustbox}%
   \end{tcolorbox}%
   \arrayrulecolor{black}%
   \par\smallskip}

% 헤더 행 헬퍼
\newcommand{\kfTHead}[1]{\textbf{\color{kfPrimary} #1}}
\newcommand{\kfTHeadRow}[1]{\rowcolor{kfPrimaryLight!50}\textbf{\color{kfPrimary} #1}}

% 빈 셀 (학생 작성용)
\newcommand{\kfTBlank}{\rule[-1.6em]{0pt}{3.4em}}
\newcommand{\kfTBlankSm}{\rule[-1.0em]{0pt}{2.4em}}

% adjustbox 강제 축소 래퍼 — Python에서 변환된 raw tabular 블록을 감쌀 때 사용
\newcommand{\kfActTableWrap}[1]{%
  \par\smallskip\needspace{4\baselineskip}%
  \begin{tcolorbox}[blanker, colback=white, left=2pt, right=2pt, top=2pt, bottom=2pt]%
  \noindent\begin{adjustbox}{max width=\linewidth, center}%
  #1%
  \end{adjustbox}%
  \end{tcolorbox}\par\smallskip%
}

% ============================================================
% activity-vocab.tex
% 어휘 학습 표 (3열: 번호 - 뜻 - 빈칸)
% 사용:
%   \begin{kfVocab}
%     \kfVocabItem{1}{눈으로 보고 마음으로 깨달음}
%     \kfVocabItem{2}{사물의 이치를 따져 헤아림}
%   \end{kfVocab}
% ============================================================

% adjustbox로 폭 강제 + 일반 tabular + p{} 열 사용
\newenvironment{kfVocab}%
  {\par\smallskip\needspace{4\baselineskip}%
   \renewcommand{\arraystretch}{1.35}%
   \arrayrulecolor{kfRule}%
   \noindent\begin{adjustbox}{max width=\linewidth, center}%
   \begin{tabular}{|>{\centering\arraybackslash}m{1.0cm}|m{8.5cm}|>{\centering\arraybackslash}m{4.0cm}|}%
   \hline
   \rowcolor{kfPrimaryLight!50}%
   \multicolumn{1}{|c|}{\textbf{\color{kfPrimary}번호}} &
   \multicolumn{1}{c|}{\textbf{\color{kfPrimary}뜻풀이}} &
   \multicolumn{1}{c|}{\textbf{\color{kfPrimary}어휘 (정답 작성)}} \\\hline
   \ignorespaces}%
  {\end{tabular}%
   \end{adjustbox}%
   \arrayrulecolor{black}%
   \par\smallskip}

\newcommand{\kfVocabItem}[2]{%
  \textbf{#1} & #2 & \rule[-1.0em]{0pt}{2.4em}\\\hline%
}

% ============================================================
% activity-blank.tex
% 빈칸 채우기 활동
% 사용:
%   \begin{kfBlankList}
%     \kfBlankItem 첫 번째 문장에서 (\kfBlank)은(는) ...
%     \kfBlankItem 둘째 문장에서 ...
%   \end{kfBlankList}
% ============================================================

\newenvironment{kfBlankList}%
  {\par\smallskip%
   \begin{enumerate}[leftmargin=20pt, itemsep=4pt, topsep=2pt,
                    label=\textbf{\arabic*.}, labelsep=6pt]}%
  {\end{enumerate}\par\smallskip}

\newcommand{\kfBlankItem}{\item}

% 텍스트 안에 박스형 빈칸 (단어 1개 채우기)
\newcommand{\kfBlankBox}[1][2.4cm]{%
  \tikz[baseline=-0.5ex]{%
    \draw[kfRule, line width=0.4pt] (0,0) -- ++(#1,0);%
  }%
}

% ============================================================
% activity-ox.tex (v15)
% O/X 표기 활동 — 그래픽 디자인
% 사용:
%   \begin{kfOXList}
%     \kfOXItem 글쓴이는 자연을 예찬하고 있다.\kfOX
%     \kfOXItem 1연과 4연은 수미상관 구조이다.\kfOX
%   \end{kfOXList}
%
% \kfOX = 우측에 [O] [X] 두 박스 (학생이 하나 동그라미)
% ============================================================

\newenvironment{kfOXList}%
  {\par\smallskip%
   \begin{enumerate}[leftmargin=20pt, itemsep=5pt, topsep=2pt,
                    label=\textbf{\arabic*.}, labelsep=6pt]}%
  {\end{enumerate}\par\smallskip}

\newcommand{\kfOXItem}{\item}

% v16: O/X 그래픽 박스 키움 (18×14 → 24×20pt) + 영역색 강조
\newcommand{\kfOX}{%
  \hfill\nolinebreak
  \tikz[baseline=-0.9ex]{%
    \node[draw=kfPrimary, fill=kfPrimary!5, line width=0.6pt, rounded corners=3pt,
          minimum width=24pt, minimum height=20pt, inner sep=0pt,
          font=\normalsize\bfseries\color{kfPrimary}] (ob) {O};%
    \node[draw=kfMute, fill=kfMute!5, line width=0.6pt, rounded corners=3pt,
          minimum width=24pt, minimum height=20pt, inner sep=0pt,
          font=\normalsize\bfseries\color{kfMute},
          right=4pt of ob] (xb) {X};%
  }%
}

% ============================================================
% activity-connect.tex
% 좌우 선 긋기 활동 (2단 양식 + 중앙 여백)
% 사용:
%   \begin{kfConnect}
%     \kfPair{사르트르}{실존은 본질에 앞선다}
%     \kfPair{니체}{초인 사상}
%     \kfPair{하이데거}{현존재}
%   \end{kfConnect}
% ============================================================

\newenvironment{kfConnect}%
  {\par\smallskip%
   \renewcommand{\arraystretch}{1.6}%
   \begin{tabular}{@{}>{\raggedleft\arraybackslash}m{4.5cm}
                     @{\hspace{2.5cm}}
                     >{\raggedright\arraybackslash}m{6.5cm}@{}}}%
  {\end{tabular}\par\smallskip}

\newcommand{\kfPair}[2]{%
  \tikz[baseline]{\node[draw=kfRule, fill=kfPrimaryLight!30, rounded corners=2pt,
        inner sep=4pt, font=\small]{#1};} \tikz[baseline]{\node[circle, draw=kfMute, fill=kfPaper, inner sep=1.5pt]{};}
  &
  \tikz[baseline]{\node[circle, draw=kfMute, fill=kfPaper, inner sep=1.5pt]{};} \tikz[baseline]{\node[draw=kfRule, fill=kfNoteBg, rounded corners=2pt,
        inner sep=4pt, font=\small]{#2};} \\
}

% ============================================================
% activity-writing.tex
% 글쓰기 활동 (큰 노트 영역)
% 사용:
%   \kfWriting{독후감}{12}     % 라벨 + 12줄 큰 노트
%   \kfWritingPrompt{프롬프트 텍스트}
% ============================================================

\newcommand{\kfWritingPrompt}[1]{%
  \par\smallskip%
  \begin{tcolorbox}[
    enhanced, breakable=false,
    colback=kfPrimaryLight!30, colframe=kfPrimary,
    boxrule=0.4pt, arc=3pt,
    left=\kfBoxPad, right=\kfBoxPad, top=\kfBoxPad, bottom=\kfBoxPad,
  ]%
  \small\textbf{\color{kfPrimary}글쓰기 안내}\\[2pt]%
  #1
  \end{tcolorbox}\par\smallskip%
}

\newcommand{\kfWriting}[2]{%
  \par\smallskip\needspace{\numexpr#2+2\relax\baselineskip}%
  \begin{tcolorbox}[
    enhanced, breakable=false,
    colback=kfPaper, colframe=kfRule,
    boxrule=0.4pt, arc=2pt,
    left=\kfBoxPad, right=\kfBoxPad, top=\kfBoxPad, bottom=\kfBoxPad,
    title={\footnotesize\bfseries\color{kfMute} #1},
    coltitle=kfMute, colbacktitle=kfPaper,
    attach boxed title to top left={yshift=-1.5mm, xshift=6pt},
    boxed title style={colframe=kfRule, boxrule=0.3pt, arc=1pt}
  ]%
  \foreach \i in {1,...,#2}{%
    \textcolor{kfRule}{\rule{\linewidth}{0.3pt}}\\[7pt]%
  }%
  \end{tcolorbox}\par\smallskip%
}

% ============================================================
% activity-structure.tex
% 구조도 (트리 + 빈칸)
% 사용:
%   \begin{kfStructure}
%     \kfNode{0}{서론}{문제 제기}
%     \kfNode{1}{본론}{(\,\quad\,)}
%     \kfNode{1}{본론}{근거 2}
%     \kfNode{0}{결론}{주장 강화}
%   \end{kfStructure}
% ============================================================

\newenvironment{kfStructure}%
  {\par\smallskip\needspace{6\baselineskip}%
   \begin{tcolorbox}[
     enhanced, breakable=false,
     colback=kfPaper, colframe=kfRule,
     boxrule=0.4pt, arc=2pt,
     left=\kfBoxPad, right=\kfBoxPad, top=\kfBoxPad, bottom=\kfBoxPad,
     title={\footnotesize\bfseries\color{kfMute} 글의 구조},
     coltitle=kfMute, colbacktitle=kfPrimaryLight!40,
     attach boxed title to top left={yshift=-1.5mm, xshift=6pt},
     boxed title style={colframe=kfRule, boxrule=0.3pt, arc=1pt}
   ]%
   \begin{tabbing}
     \hspace{1.4cm}\=\hspace{1.4cm}\=\hspace{8cm}\kill}%
  {\end{tabbing}\end{tcolorbox}\par\smallskip}

% 노드: #1=들여쓰기 단계 (0~3), #2=라벨, #3=내용
\newcommand{\kfNode}[3]{%
  \ifcase#1\relax%
    \>\>\textbf{\color{kfPrimary} ▶ #2}\quad #3\\
  \or
    \>\>\quad\textbf{\color{kfMute} ◦ #2}\quad #3\\
  \or
    \>\>\quad\quad\textbf{\color{kfMute} - #2}\quad #3\\
  \or
    \>\>\quad\quad\quad\textbf{\color{kfMute} \textperiodcentered\ #2}\quad #3\\
  \fi
}

% 빈칸 노드
\newcommand{\kfNodeBlank}[2]{\kfNode{#1}{#2}{(\hspace{2.5cm})}}

% ============================================================
% activity-sentence-reading.tex
% 문장 독해 활동 — 문장 박스 1개 + 그에 묶인 문제 N개를 한 minipage로
% 사용:
%   \begin{kfSentenceUnit}{1}{"바람은 내 귀에 속삭이며..."}
%     \kfSentQ{(1)}{서술어 '속삭이며'의 주어는?}
%     \begin{kfChoices}
%       \kfChoice 바람은
%       ...
%     \end{kfChoices}
%   \end{kfSentenceUnit}
% ============================================================

% kfSentenceBox 매크로는 passage-note.tex에 정의됨

\newenvironment{kfSentenceUnit}[2]%
  {\par\smallskip\needspace{6\baselineskip}%
   \noindent\begin{minipage}{\linewidth}%
   \samepage%
   \kfSentenceBox{문장 #1}{#2}%
   \par\smallskip%
  }%
  {\end{minipage}\par\smallskip}

% 같은 문장에 묶인 소문항 (문장 안내 + 발문)
\newcommand{\kfSentQ}[2]{%
  \par\noindent\textbf{#1}~#2\par\smallskip%
}

% 헤더 없이 문장만 있는 묶음 (sentence_ref가 없는 일반 문항)
\newenvironment{kfSentencelessUnit}%
  {\par\smallskip\noindent\begin{minipage}{\linewidth}\samepage}%
  {\end{minipage}\par\smallskip}

% ============================================================
% chapter-cover.tex (v6)
% 챕터 진입 페이지 (표지) — 하단에 영역 목차 추가
% 사용:
%   \kfChapterCover{1}{근대성과 자아}{Chapter 1}{이번 챕터에서는 ...}
%   \begin{kfChapterAreaList}
%     \kfChapterAreaItem{어휘}{kfAreaVocab}{핵심 어휘 12개}
%     ...
%   \end{kfChapterAreaList}
%
%   영역 목차가 없을 때는 \kfChapterCoverEnd 호출
% ============================================================

\newcommand{\kfChapterCover}[4]{%
  \clearpage%
  \thispagestyle{kfBlank}%
  \kfMarkChapter{Chapter #1. #2}%
  \kfChapterCoverBg%
  % v18.5: 챕터 표지 우상단에 로고
  \begin{tikzpicture}[remember picture, overlay]
    \node[anchor=north east, inner sep=0pt] at ($(current page.north east)+(-18mm,-18mm)$) {\kfLogoMd};
  \end{tikzpicture}%
  \vspace*{20mm}%
  \noindent
  \begin{tikzpicture}[remember picture, overlay]
    % 좌측 액센트 바
    \fill[kfPrimary] (current page.west) ++(12mm,25mm) rectangle ++(5mm, -50mm);
    % 챕터 번호 배지 (이중 원, 약간 작게)
    \node[anchor=north west, inner sep=0pt] at ($(current page.north west)+(24mm,-45mm)$) {%
      \begin{tikzpicture}
        \draw[kfPrimary, line width=1pt] (0,0) circle (10mm);
        \draw[kfPrimary, line width=0.4pt] (0,0) circle (8mm);
        \node[font=\normalsize\bfseries\color{kfPrimary}] at (0,2mm) {CHAPTER};
        \node[font=\LARGE\bfseries\color{kfPrimary}] at (0,-3mm) {#1};
      \end{tikzpicture}%
    };
  \end{tikzpicture}%
  \vspace{42mm}%
  \par\noindent\hspace{32mm}%
  \begin{minipage}{0.65\linewidth}%
    {\footnotesize\color{kfMute} #3}\\[4pt]
    {\LARGE\bfseries\color{kfPrimary} #2}\\[8pt]
    {\color{kfRule}\rule{50mm}{0.5pt}}\\[6pt]
    {\small\color{kfInk} #4}
  \end{minipage}%
  \par\vspace{14mm}%
}

% v6 / v18.5 / v18.6: 챕터 표지의 영역 목차 — 흰색 반투명 박스
% v18.6: 배경 가로선/도형과 안 겹치도록 TikZ overlay로 우하단 절대 위치
\makeatletter
% 영역 항목들을 모아 둘 박스
\newsavebox{\kf@arealistbox}
\newenvironment{kfChapterAreaList}{%
  \begin{lrbox}{\kf@arealistbox}%
  \begin{minipage}{0.62\linewidth}%
    \begin{tcolorbox}[%
      enhanced, colback=white, colframe=white,
      opacityback=0.92, opacityframe=0,
      boxrule=0pt, arc=4pt,
      width=\linewidth,
      top=8pt, bottom=8pt, left=10pt, right=10pt,
    ]%
      {\footnotesize\bfseries\color{kfMute} 이 챕터의 영역}\par\vspace{4pt}%
      {\color{kfRule}\hrule height 0.3pt}\par\vspace{4pt}%
      \begin{itemize}[leftmargin=14pt, itemsep=2pt, topsep=0pt, label={\color{kfPrimary!70}\textbullet}]%
}{%
      \end{itemize}%
    \end{tcolorbox}%
  \end{minipage}%
  \end{lrbox}%
  % v18.7: 배경 상단 가로선 ~ 하단 가로선 사이의 중간 영역에 배치
  % 가로선이 내용을 가리지 않도록 박스 중심이 두 선 사이 중점에 위치
  \begin{tikzpicture}[remember picture, overlay]%
    \node[anchor=center, inner sep=0pt]
      at ($(current page.south)+(0mm,90mm)$)
      {\usebox{\kf@arealistbox}};%
  \end{tikzpicture}%
  \vfill%
  \noindent\hfill\kfSeriesBadge\hspace{12mm}%
  \clearpage%
}
\makeatother

% 영역 항목: \kfChapterAreaItem{영역명}{kfArea색}{부제}
\makeatletter
\newcommand{\kfChapterAreaItem}[3]{%
  \item {\small\bfseries\color{#2} #1}%
  \def\kf@cci@sub{#3}%
  \ifx\kf@cci@sub\@empty\else
    {\small\color{kfMute}\, \textemdash\, #3}%
  \fi
}
\makeatother

% 영역 목록 없이 cover만 (legacy)
\newcommand{\kfChapterCoverEnd}{%
  \vfill%
  \noindent\hfill\kfSeriesBadge\hspace{12mm}%
  \clearpage%
}

% ============================================================
% area-header.tex (v16)
% 영역 시작 표제 — 시리즈별 차별화 라벨 + 큰 영역명 + 부제
% PAGE_BREAK_POLICY.md §1.4
%
% v16 (디자인 고도화 Phase 1):
%   - 시리즈별 라벨 박스 추가:
%     소쉬르 = AREA START (영역색 채움 박스)
%     프레게 = SECTION OPEN (영역색 테두리)
%     러셀   = RUSSELL (영역색 라이트 박스)
%     비트   = WITTGENSTEIN (회색 라이트 박스)
%   - 라벨 위치: 영역명 위 좌측
% ============================================================

\makeatletter

% 시리즈 라벨 텍스트
\newcommand{\kf@serieslabel}{%
  \ifcase\kf@series\relax
    AREA START%
  \or
    AREA START%
  \or
    SECTION OPEN%
  \or
    RUSSELL%
  \or
    WITTGENSTEIN%
  \fi
}

% 시리즈 라벨 박스 스타일 (#1 = 영역색)
\newcommand{\kf@labelbox}[1]{%
  \ifcase\kf@series\relax
    % 0/1 = 소쉬르: 영역색 채움
    \tikz[baseline=-2pt]{\node[fill=#1, text=white,
      rounded corners=3pt, inner xsep=6pt, inner ysep=2pt,
      font=\scriptsize\bfseries]{\kf@serieslabel};}%
  \or
    % 1 = 소쉬르: 영역색 채움
    \tikz[baseline=-2pt]{\node[fill=#1, text=white,
      rounded corners=3pt, inner xsep=6pt, inner ysep=2pt,
      font=\scriptsize\bfseries]{\kf@serieslabel};}%
  \or
    % 2 = 프레게: 영역색 테두리
    \tikz[baseline=-2pt]{\node[draw=#1, line width=0.6pt, text=#1,
      rounded corners=3pt, inner xsep=6pt, inner ysep=2pt,
      font=\scriptsize\bfseries]{\kf@serieslabel};}%
  \or
    % 3 = 러셀: 영역색 라이트 박스
    \tikz[baseline=-2pt]{\node[fill=#1!12, text=#1,
      rounded corners=2pt, inner xsep=5pt, inner ysep=1.5pt,
      font=\scriptsize\bfseries]{\kf@serieslabel};}%
  \or
    % 4 = 비트: 회색 미니 박스
    \tikz[baseline=-2pt]{\node[fill=kfMute!10, text=kfMute,
      rounded corners=2pt, inner xsep=5pt, inner ysep=1.5pt,
      font=\scriptsize\bfseries]{\kf@serieslabel};}%
  \fi
}

% v17: 시리즈+영역별 Canva 배경 자동 매핑
% \kf@hasbg=1로 set되면 \kfAreaStart에서 라벨 박스 출력 생략
\def\kf@areaVocab{어휘}\def\kf@areaGrammar{문법}\def\kf@areaConcept{개념}
\def\kf@areaLit{문학}\def\kf@areaNonlit{비문학}
\newif\ifkf@bgon
% 파일 있으면 배경 적용 + boolean true. 없으면 nothing.
\newcommand{\kfBgSet}[1]{%
  \IfFileExists{#1.pdf}{\kfPageBg{#1}\kf@bgontrue}{}%
}
\newcommand{\kf@trybg}[1]{%
  % v18.4: 영역 배경 PDF 비활성화 — 큰 도형이 본문 영역을 침범해
  % "영역 헤더 후 빈 페이지" 문제를 일으킴. 라벨 박스로만 영역 표시.
  \kf@bgonfalse%
}

% Note: 러셀(rus_*)/비트(wit_*) 배경 PDF 미존재 시 \kfPageBg가 에러 → \IfFileExists로 가드
% (현재는 파일이 모두 존재한다고 가정. 부분 제공 시 cls 수정 필요)

\newcommand{\kfAreaStart}[3]{%
  \clearpage%
  \thispagestyle{fancy}%
  \kfMarkArea{#1}%
  \kf@trybg{#1}%
  \vspace*{1mm}%
  % 배경 없으면 라텍스 라벨 박스 출력 (배경 있으면 일러스트가 라벨 역할)
  \ifkf@bgon
    \vspace*{4mm}%
  \else
    \noindent\kf@labelbox{#2}\par\vspace{4pt}%
  \fi
  % 영역명 + 부제
  \noindent
  \begin{tikzpicture}
    \node[anchor=west, inner sep=0pt] (bar) at (0,0) {\color{#2}\rule{3pt}{22pt}};
    \node[anchor=base west, inner sep=0pt, xshift=8pt, yshift=2pt] (lbl) at (bar.east |- 0,0) {%
      {\LARGE\bfseries\color{#2} #1}%
    };
    \def\kf@areasub{#3}%
    \ifx\kf@areasub\@empty\else
      \node[anchor=base west, inner sep=0pt, xshift=8pt, font=\normalsize\color{kfMute}]
        at (lbl.base east) {\,\textperiodcentered\, #3};%
    \fi
  \end{tikzpicture}%
  \par\vspace{2pt}
  {\color{#2!70}\hrule height 1pt}%
  \par\vspace{1pt}%
  {\color{kfRule}\hrule height 0.3pt}%
  \par\vspace{4pt}%
  \needspace{6\baselineskip}\nopagebreak[4]%
  \global\kf@areaJustOpenedtrue% v18.5: 첫 컨텐츠 needspace 무시
}
\makeatother

% 시리즈/영역 단축 매크로
\newcommand{\kfStartVocab}[1]{\kfAreaStart{어휘}{kfAreaVocab}{#1}}
\newcommand{\kfStartGrammar}[1]{\kfAreaStart{문법}{kfAreaGrammar}{#1}}
\newcommand{\kfStartConcept}[1]{\kfAreaStart{개념}{kfAreaConcept}{#1}}
\newcommand{\kfStartLit}[1]{\kfAreaStart{문학}{kfAreaLiterature}{#1}}
\newcommand{\kfStartNonlit}[1]{\kfAreaStart{비문학}{kfAreaNonlit}{#1}}
\newcommand{\kfStartWeekly}[1]{\kfAreaStart{실력 확인}{kfAreaWeekly}{#1}}

% ============================================================
% section-header.tex
% 섹션 시작 라벨 헤더 — [지문] / [활동] / [문제] 라벨
% 좌측 액센트 바 + 라벨 + 부제 (영역 액센트 색)
%
% 사용:
%   \kfSecHeader{kfAreaLiterature}{지문}{빼앗긴 들에도 봄은 오는가 / 이상화}
%   \kfSecHeader{kfAreaConcept}{활동}{내용 확인}
%   \kfSecHeader{kfAreaNonlit}{문제}{}
%
% 헤더 직후 페이지 분할 금지: \needspace{4\baselineskip} + \nopagebreak
% ============================================================

\providecommand{\kfSecHeader}[3]{%
  \par\needspace{10\baselineskip}\filbreak%
  \vspace{3pt}%
  \noindent
  \begin{tikzpicture}
    \node[anchor=west, inner sep=0pt] (bar) at (0,0) {\color{#1}\rule{3pt}{14pt}};
    \node[anchor=west, inner sep=0pt, xshift=6pt, font=\normalsize\bfseries\color{#1}]
       (lbl) at (bar.east) {[\,#2\,]};
    \node[anchor=west, inner sep=0pt, xshift=4pt, font=\small\color{kfMute}]
       at (lbl.east) {#3};
  \end{tikzpicture}%
  \par\vspace{1pt}%
  {\color{#1!50}\hrule height 0.4pt}%
  \par\vspace{2pt}\nopagebreak%
}

% 영역별 단축 매크로
\newcommand{\kfSecPassageVocab}[1]{\kfSecHeader{kfAreaVocab}{지문}{#1}}
\newcommand{\kfSecPassageGrammar}[1]{\kfSecHeader{kfAreaGrammar}{지문}{#1}}
\newcommand{\kfSecPassageConcept}[1]{\kfSecHeader{kfAreaConcept}{지문}{#1}}
\newcommand{\kfSecPassageLit}[1]{\kfSecHeader{kfAreaLiterature}{지문}{#1}}
\newcommand{\kfSecPassageNonlit}[1]{\kfSecHeader{kfAreaNonlit}{지문}{#1}}
\newcommand{\kfSecPassageWeekly}[1]{\kfSecHeader{kfAreaWeekly}{지문}{#1}}

\newcommand{\kfSecActVocab}[1]{\kfSecHeader{kfAreaVocab}{활동}{#1}}
\newcommand{\kfSecActGrammar}[1]{\kfSecHeader{kfAreaGrammar}{활동}{#1}}
\newcommand{\kfSecActConcept}[1]{\kfSecHeader{kfAreaConcept}{활동}{#1}}
\newcommand{\kfSecActLit}[1]{\kfSecHeader{kfAreaLiterature}{활동}{#1}}
\newcommand{\kfSecActNonlit}[1]{\kfSecHeader{kfAreaNonlit}{활동}{#1}}
\newcommand{\kfSecActWeekly}[1]{\kfSecHeader{kfAreaWeekly}{활동}{#1}}

\newcommand{\kfSecQVocab}[1]{\kfSecHeader{kfAreaVocab}{문제}{#1}}
\newcommand{\kfSecQGrammar}[1]{\kfSecHeader{kfAreaGrammar}{문제}{#1}}
\newcommand{\kfSecQConcept}[1]{\kfSecHeader{kfAreaConcept}{문제}{#1}}
\newcommand{\kfSecQLit}[1]{\kfSecHeader{kfAreaLiterature}{문제}{#1}}
\newcommand{\kfSecQNonlit}[1]{\kfSecHeader{kfAreaNonlit}{문제}{#1}}
\newcommand{\kfSecQWeekly}[1]{\kfSecHeader{kfAreaWeekly}{문제}{#1}}


\begin{document}

\thispagestyle{empty}
\vspace*{40mm}
\begin{center}
{\Huge\bfseries\color{kfPrimary} 러셀1}\\[10pt]
{\Large\color{kfMute} 학생용 교재 — 제 1 권}\\[20pt]
{\small\color{kfMute} (챕터 1 \textendash{} 4)}
\end{center}
\vfill
\hfill\kfSeriesBadge\hspace{15mm}
\clearpage
\kfChapterCover{1}{러셀1 (챕터1)}{Chapter 1}{이번 챕터의 학습 내용을 시작합니다.}

\kfStartGrammar{문법 영역 — 문장 속 단어들의 역할 (품사)}


\kfSecHeader{kfAreaGrammar}{개념}{문장 속 단어들의 역할 (품사)}

우리말에는 헤아릴 수 없이 많은 단어가 있다. 이 수많은 단어를 공통된 성질에 따라 분류해 놓은 갈래를 ‘품사’라고 한다. 품사는 단어가 지닌 고유한 성격을 알려 주는 이름표와도 같다.

품사를 분류하는 기준은 크게 세 가지이다. 첫 번째 기준은 ‘형태’로, 단어가 문장 속에서 형태가 변하는지 여부에 따라 나눈다. ‘나무’, ‘헌’, ‘아주’와 같은 단어들은 문장에서 쓰일 때 그 형태가 변하지 않는데, 이를 ‘불변어’라고 한다. 반면에 ‘먹다’, ‘예쁘다’와 같은 단어들은 문장에서 쓰일 때 형태가 다양하게 변한다. ‘먹다’는 ‘먹고, 먹으니, 먹어서’로, ‘예쁘다’는 ‘예쁘고, 예쁘니, 예뻐서’로 활용된다. 이처럼 문장 속에서 형태가 바뀌는 단어를 ‘가변어’라고 한다.

두 번째 기준은 ‘기능’으로, 단어가 문장 안에서 어떤 역할을 하는지에 따라 분류한다. 문장의 몸통과 같이 중심이 되어 주어나 목적어 등으로 쓰이는 단어를 ‘체언’이라고 한다. 주로 체언 뒤에 붙어 다른 말과의 문법적 관계를 나타내는 단어는 ‘관계언’이라고 한다. 문장에서 주어의 움직임이나 상태를 풀이하는 단어는 ‘용언’이라고 하며, 다른 말을 꾸며 주는 역할을 하는 단어는 ‘수식언’이라고 한다. 마지막으로 문장의 다른 성분과 직접적인 관계를 맺지 않고 독립적으로 쓰이는 단어를 ‘독립언’이라고 한다.

세 번째 기준은 ‘의미’로, 단어가 지닌 공통된 뜻에 따라 품사를 분류하는 것이다. 체언에는 사람이나 사물의 이름을 나타내는 ‘명사’, 명사를 대신하는 ‘대명사’, 수량이나 순서를 나타내는 ‘수사’가 있다. 관계언에는 ‘조사’가 속한다. 용언은 움직임을 나타내는 ‘동사’와 성질이나 상태를 나타내는 ‘형용사’로 나뉜다. 수식언에는 체언을 꾸미는 ‘관형사’와 주로 용언을 꾸미는 ‘부사’가 있다. 마지막으로 독립언에는 놀람이나 부름, 대답을 나타내는 ‘감탄사’가 속한다.

이처럼 우리말은 형태 기준에 따라 불변어와 가변어로 나누고, 기능 기준에 따라 체언, 관계언, 용언, 수식언, 독립언으로 나누며, 의미 기준에 따라 명사, 대명사, 수사, 조사, 동사, 형용사, 관형사, 부사, 감탄사의 9품사로 나눈다.
\par\medskip\begin{itemize}[leftmargin=18pt, itemsep=2pt, topsep=2pt, label=\textbullet]
\item 형태(불변/가변) — '나무·헌·아주'(불변) / '먹다·예쁘다'(가변)
\item 기능 — '체언(학교)·용언(뛰다)·수식언(예쁜)·관계언(이/가)·독립언(아!)'
\item 의미 9품사 — '명·대·수 / 조 / 동·형 / 관·부 / 감'
\end{itemize}

\kfSecHeader{kfAreaGrammar}{활동}{내용 확인}
\noindent\textit{\small 아래 표의 빈칸에 알맞은 말을 채워 문법 내용을 확인해 보자.}\par\medskip
* 품사 분류의 세 가지 기준 *

우리말 단어를 나누는 세 가지 기준과 9품사의 관계를 표로 정리해 보자. 

\par\medskip\needspace{4\baselineskip}
\noindent\begin{adjustbox}{max width=\linewidth, center}
\renewcommand{\arraystretch}{1.45}
\arrayrulecolor{kfRule}
\begin{tabular}{|>{\raggedright\arraybackslash}p{\dimexpr 0.4306\linewidth-12pt\relax}|>{\raggedright\arraybackslash}p{\dimexpr 0.4167\linewidth-12pt\relax}|>{\raggedright\arraybackslash}p{\dimexpr 0.1528\linewidth-12pt\relax}|}
\hline
\rowcolor{kfPrimaryLight!50}\textbf{\color{kfPrimary} 형태} & \textbf{\color{kfPrimary} 기능} & \textbf{\color{kfPrimary} 의미} \\\hline
[ ① \_\_\_\_\_\_ ] (형태가 변하지 않음) & [ ② \_\_\_\_\_\_ ] (문장의 몸통) & 명사, 대명사, 수사 \\\hline
\rule[-1.0em]{0pt}{2.4em} & [ ③ \_\_\_\_\_\_ ] (관계 표시) & 조사 \\\hline
\rule[-1.0em]{0pt}{2.4em} & [ ④ \_\_\_\_\_\_ ] (꾸며주는 말) & 관형사, 부사 \\\hline
\rule[-1.0em]{0pt}{2.4em} & [ ⑤ \_\_\_\_\_\_ ] (독립적인 말) & 감탄사 \\\hline
[ ⑥ \_\_\_\_\_\_ ] (형태가 변함) & [ ⑦ \_\_\_\_\_\_ ] (움직임·상태 서술) & 동사, 형용사 \\\hline
\end{tabular}
\arrayrulecolor{black}
\end{adjustbox}\par\medskip

\kfSecHeader{kfAreaGrammar}{활동}{분석 훈련}
\noindent\textit{\small  9품사 분류 훈련! 다음 단어의 품사를 9품사 중 하나로 쓰시오.}\par\medskip
\begin{enumerate}[leftmargin=22pt, itemsep=6pt, topsep=4pt, label=\textbf{\arabic*.}]
\item 하늘
\item 달린다
\item 아름답다
\item 우리
\item 셋
\item 에게
\item 온갖
\item 아주
\item 어머나
\item 먹고
\end{enumerate}

\kfSecHeader{kfAreaGrammar}{문제}{}

\begin{kfQBlock}{1}{}
\kfQStem{품사를 나누는 세 가지 기준에 대한 설명으로 적절하지 \uline{않은} 것은?}
\begin{kfChoices}
\kfChoice 단어의 ‘소리(음운)’가 가장 우선적인 분류 기준이 된다.
\kfChoice 문장 안에서 하는 ‘기능’에 따라 체언·용언 등으로 나뉜다.
\kfChoice 단어가 가지는 ‘의미’에 따라 명사·동사 등으로 세분화된다.
\kfChoice 형태·기능·의미의 세 기준은 함께 작용해 품사를 결정한다.
\kfChoice 단어의 ‘형태’가 변하는지에 따라 가변어와 불변어로 나뉜다.
\end{kfChoices}
\end{kfQBlock}
\begin{kfQBlock}{2}{}
\kfQStem{형태가 변하는 '가변어'에 속하는 품사는?}
\begin{kfChoices}
\kfChoice 명사
\kfChoice 동사
\kfChoice 부사
\kfChoice 관형사
\kfChoice 감탄사
\end{kfChoices}
\end{kfQBlock}
\begin{kfQBlock}{3}{}
\kfQStem{기능에 따른 묶음이 나머지와 \uline{다른} 하나는?}
\begin{kfChoices}
\kfChoice 학생
\kfChoice 우리
\kfChoice 첫째
\kfChoice 아주
\kfChoice 나무
\end{kfChoices}
\end{kfQBlock}
\begin{kfQBlock}{4}{}
\kfQStem{밑줄 친 단어의 품사가 나머지와 \uline{다른} 하나는?}
\begin{kfChoices}
\kfChoice 새가 \uline{달린다}.
\kfChoice 밥을 모두 \uline{먹다}.
\kfChoice 하늘이 \uline{아름답다}.
\kfChoice 동생이 \uline{잔다}.
\kfChoice 기차가 \uline{오다}.
\end{kfChoices}
\end{kfQBlock}
\begin{kfQBlock}{5}{}
\kfQStem{'수식언'에 대한 설명으로 옳지 \uline{않은} 것은?}
\begin{kfChoices}
\kfChoice 다른 말을 꾸며 주는 기능을 한다.
\kfChoice 사람이나 사물의 이름을 가리키는 기능을 한다.
\kfChoice 관형사는 체언(명사 등)을 꾸며 준다.
\kfChoice 부사는 주로 용언(동사·형용사)을 꾸며 준다.
\kfChoice 관형사와 부사가 수식언에 속한다.
\end{kfChoices}
\end{kfQBlock}
\begin{kfQBlock}{6}{}
\kfQStem{<보기>의 밑줄 친 단어와 품사가 같은 것은?}
\begin{kfEvidence}나는 \uline{새} 옷을 입었다.\end{kfEvidence}
\begin{kfChoices}
\kfChoice \uline{하늘}이 무척 푸르고 맑다.
\kfChoice 오늘 나는 \uline{매우} 기쁘다.
\kfChoice \uline{저} 사람이 내 친구다.
\kfChoice \uline{어머나}! 정말 멋지구나.
\kfChoice 동생이 \uline{밥}을 다 먹는다.
\end{kfChoices}
\end{kfQBlock}
\begin{kfQBlock}{7}{}
\kfQStem{밑줄 친 단어의 품사가 '부사'인 것은?}
\begin{kfChoices}
\kfChoice 토끼가 \uline{빨리} 달린다.
\kfChoice 나는 \uline{책}을 읽는다.
\kfChoice \uline{셋}이서 놀러 갔다.
\kfChoice \uline{와}! 정말 멋지다.
\kfChoice \uline{모든} 학생이 모였다.
\end{kfChoices}
\end{kfQBlock}
\begin{kfQBlock}{8}{}
\kfQStem{문장 속에서 형태가 변하지 않는 '불변어'가 \uline{아닌} 것은?}
\begin{kfChoices}
\kfChoice 하늘
\kfChoice 그리고
\kfChoice 깨끗하다
\kfChoice 에게
\kfChoice 헌
\end{kfChoices}
\end{kfQBlock}
\begin{kfQBlock}{9}{}
\kfQStem{'용언'에 대한 설명으로 옳지 \uline{않은} 것은?}
\begin{kfChoices}
\kfChoice 사물의 움직임이나 상태를 서술하는 기능을 한다.
\kfChoice 형태가 고정된 불변어로, 다른 말을 꾸며 주는 기능을 한다.
\kfChoice 활용을 통해 형태가 변하는 가변어에 해당한다.
\kfChoice 문장에서 주로 서술어 자리에 쓰여 의미를 매듭지어 준다.
\kfChoice 동사와 형용사가 모두 용언의 범주에 포함된다.
\end{kfChoices}
\end{kfQBlock}
\begin{kfQBlock}{1}{}
\kfQStem{일정한 기준에 따라 단어의 종류를 나누어 놓은 체계는 무엇인가?}
\kfAnswerLines{1}
\end{kfQBlock}
\begin{kfQBlock}{2}{}
\kfQStem{문장 속에서 형태가 바뀌는 단어들을 무엇이라고 하는가?}
\kfAnswerLines{1}
\end{kfQBlock}
\begin{kfQBlock}{3}{}
\kfQStem{문장 속에서 형태가 고정되어 변하지 않는 단어들은 무엇인가?}
\kfAnswerLines{1}
\end{kfQBlock}
\begin{kfQBlock}{4}{}
\kfQStem{문장에서 몸통처럼 중심이 되어 주어나 목적어 등으로 쓰이는 단어들의 묶음은 무엇인가?}
\kfAnswerLines{1}
\end{kfQBlock}
\begin{kfQBlock}{5}{}
\kfQStem{체언을 꾸미거나 용언을 꾸미는 기능을 하는 단어들의 묶음은 무엇인가?}
\kfAnswerLines{1}
\end{kfQBlock}
\begin{kfQBlock}{6}{}
\kfQStem{주어의 움직임이나 상태를 풀이하는 기능을 하는 단어들의 묶음은 무엇인가?}
\kfAnswerLines{1}
\end{kfQBlock}
\begin{kfQBlock}{7}{}
\kfQStem{체언에 속하는 3가지 품사를 쓰시오.}
\kfAnswerLines{1}
\end{kfQBlock}
\begin{kfQBlock}{8}{}
\kfQStem{용언에 속하는 2가지 품사를 쓰시오.}
\kfAnswerLines{1}
\end{kfQBlock}
\begin{kfQBlock}{9}{}
\kfQStem{'새', '헌', '모든'처럼 체언을 꾸며주는 품사는 무엇인가?}
\kfAnswerLines{1}
\end{kfQBlock}
\begin{kfQBlock}{10}{}
\kfQStem{'와', '네', '여보세요'처럼 독립적으로 쓰이는 품사는 무엇인가?}
\kfAnswerLines{1}
\end{kfQBlock}
\begin{kfQBlock}{1}{}
\kfQStem{다음 문장에서 '헌'과 '아주'의 품사를 각각 쓰고, 이 단어들이 문장에서 어떤 기능을 하는지 서술하시오.}
\begin{kfEvidence}나는 헌 책을 아주 재미있게 읽었다.\end{kfEvidence}
\kfAnswerNote{답안}{4}
\end{kfQBlock}


\kfStartConcept{개념 영역 — 표현을 생생하게 만드는 비유법}


\kfSecHeader{kfAreaConcept}{개념}{표현을 생생하게 만드는 비유법}

비유법은 표현하고 싶은 대상을 그와 비슷한 다른 대상에 빗대어 더 생생하고 인상 깊게 나타내는 표현 방법이다. 이때 진짜 말하고 싶은 대상을 '원관념'이라 하고, 빗대어 표현하는 대상을 '보조관념'이라고 한다. 원관념과 보조관념은 서로 공통점이 있어야 자연스러운 비유가 이루어진다.

직유법은 '같이', '처럼', '듯이', '-인 양' 등 비유를 나타내는 말을 사용하여 원관념과 보조관념을 직접 연결하는 방법이다. 「빼앗긴 들에도 봄은 오는가」에서 "가르마 같은 논길"은 '논길(원관념)'을 '가르마(보조관념)'에 비유하여 일직선으로 뻗은 모습을 표현한 것이다. 또한 "살찐 가슴처럼 부드러운 이 흙"은 '흙(원관념)'을 '살찐 가슴(보조관념)'에 빗대어 흙의 부드럽고 풍요로운 느낌을 강조한다. 직유법은 비교하는 말이 직접 드러나 있어 독자가 그 의미를 이해하기 쉽다.

은유법은 'A는 B이다' 또는 'A의 B' 같은 형태로 원관념과 보조관념을 간접적으로 연결하는 방법이다. "내 마음은 호수요"는 '마음(A)은 호수(B)이다'의 형태로, 마음이 호수처럼 넓고 고요하다는 의미를 함축적으로 드러낸다. 직유법보다 더 깊은 생각을 하게 만든다.

의인법은 사람이 아닌 사물이나 동식물에 인격을 부여하여 사람처럼 생각하고 행동하는 것으로 표현하는 방법이다. 「빼앗긴 들에도 봄은 오는가」에서는 "입을 다물고 있는 하늘아, 들판아"라고 하여 하늘과 들판을 사람처럼 표현했다. 또한 "바람은 내 귀에 속삭이며"나, "착한 개울이 ... 아기를 달래는 노래를 하며 ... 어깨춤을 추며 흘러가네"처럼 자연물이 사람처럼 행동하게 하여 대상에 친근감을 부여하고 생동감을 불어넣는다.

대유법도 비유법의 한 종류로, 한 부분으로 전체를 나타내거나 사물의 특징으로 그 사물 자체를 나타내는 방법이다. 「빼앗긴 들에도 봄은 오는가」의 '빼앗긴 들'은 '들'이라는 부분(보조관념)을 통해 '빼앗긴 우리나라 국토'라는 전체(원관념)를 상징적으로 보여준다.

이러한 비유법은 평범한 내용을 더 새롭고 감동적으로 전달하며, 독자가 작품의 의미를 더 깊이 있게 감상하도록 돕는 중요한 표현 기법이다.
\par\medskip\begin{itemize}[leftmargin=18pt, itemsep=2pt, topsep=2pt, label=\textbullet]
\item 직유법 — '가르마 같은 논길'(논길↔가르마, '같은' 연결)
\item 은유법 — '내 마음은 호수요'(A는 B이다)
\item 대유법 — '빼앗긴 들'→'들'로 국토 전체를 상징
\end{itemize}

\kfSecHeader{kfAreaConcept}{문제}{}

\begin{kfQBlock}{1}{}
\kfQStem{비유법을 사용하는 이유로 적절하지 \uline{않은} 것은 무엇인가?}
\begin{kfChoices}
\kfChoice 어려운 한자어를 쉽게 풀어쓰기 위한 맞춤법 규정이기 때문이다.
\kfChoice 추상적 개념을 구체적인 대상에 빗대어 이해하기 쉽게 만들기 위해서이다.
\kfChoice 익숙한 대상으로 낯선 대상의 특징을 효과적으로 떠올리게 하기 위해서이다.
\kfChoice 독자의 흥미를 끌고 정서적 공감을 자아내기 위해서이다.
\kfChoice 표현을 더 생생하고 인상 깊게 만들기 위해서이다.
\end{kfChoices}
\end{kfQBlock}
\begin{kfQBlock}{2}{}
\kfQStem{"내 마음은 호수요"라는 표현에 대한 설명으로 알맞지 \uline{않은} 것은 무엇인가?}
\begin{kfChoices}
\kfChoice '처럼·같이' 같은 연결어를 사용한 직유법이다.
\kfChoice 'A는 B이다' 형식으로 두 대상을 직접 빗댄 은유법이다.
\kfChoice 추상적인 마음을 구체적인 호수에 빗대어 표현하였다.
\kfChoice 호수가 지닌 고요함·잔잔함의 이미지가 마음에 전이된다.
\kfChoice 원관념은 '마음', 보조관념은 '호수'이다.
\end{kfChoices}
\end{kfQBlock}
\begin{kfQBlock}{3}{}
\kfQStem{다음 중 직유법을 만들기 위해 사용되는 말이 \uline{아닌} 것은 무엇인가?}
\begin{kfChoices}
\kfChoice 같이
\kfChoice 이다
\kfChoice 듯이
\kfChoice -인 양
\kfChoice 처럼
\end{kfChoices}
\end{kfQBlock}
\begin{kfQBlock}{4}{}
\kfQStem{"가르마 같은 논길"이라는 표현에 대한 설명으로 알맞지 \uline{않은} 것은 무엇인가?}
\begin{kfChoices}
\kfChoice 사람이 아닌 것을 사람처럼 표현했다.
\kfChoice '가르마'가 보조관념이다.
\kfChoice 직유법이 사용되었다.
\kfChoice 논길의 일직선으로 뻗은 모습을 표현했다.
\kfChoice '논길'이 원관념이다.
\end{kfChoices}
\end{kfQBlock}
\begin{kfQBlock}{5}{}
\kfQStem{은유법에 대한 설명으로 적절하지 \uline{않은} 것은 무엇인가?}
\begin{kfChoices}
\kfChoice 'A는 B이다'의 형태로 두 대상을 간접적으로 연결한다.
\kfChoice '같이·처럼·같은' 같은 연결어를 사용해 두 대상을 직접 잇는다.
\kfChoice 추상적 대상을 구체적 대상에 빗대어 형상화하기에 좋다.
\kfChoice '내 마음은 호수요'처럼 원관념과 보조관념의 공통점을 부각한다.
\kfChoice 원관념과 보조관념을 직접 등치시켜 의미를 강조한다.
\end{kfChoices}
\end{kfQBlock}
\begin{kfQBlock}{6}{}
\kfQStem{다음 중, 사람처럼 표현한 대상이 \uline{아닌} 것은 무엇인가?}
\begin{kfChoices}
\kfChoice 입을 다물고 있는 하늘아
\kfChoice 바람은 내 귀에 속삭이며
\kfChoice 착한 개울이 신나게 어깨춤을 추며
\kfChoice 살찐 가슴처럼 부드러운 이 흙
\kfChoice 아기를 달래는 노래
\end{kfChoices}
\end{kfQBlock}
\begin{kfQBlock}{7}{}
\kfQStem{의인법을 사용했을 때 얻을 수 있는 효과로 가장 알맞은 것은 무엇인가?}
\begin{kfChoices}
\kfChoice 대상의 의미를 반대로 바꾸어 강조한다.
\kfChoice 대상을 더 객관적으로 보이게 한다.
\kfChoice 대상에 친근감을 부여하고 생동감을 준다.
\kfChoice 표현이 간결해지고 딱딱한 느낌을 준다.
\kfChoice 숨겨진 의미를 찾아야 해서 이해하기 어렵다.
\end{kfChoices}
\end{kfQBlock}
\begin{kfQBlock}{8}{}
\kfQStem{다음 중, 지문에서 설명하는 비유법의 종류가 \uline{다른} 하나는 무엇인가?}
\begin{kfChoices}
\kfChoice 바다가 푸른 숨을 쉰다
\kfChoice 노란 은행잎이 우리에게 손짓한다
\kfChoice 내 마음은 호수요
\kfChoice 바람은 내 귀에 속삭이며
\kfChoice 착한 개울이 아기를 달래는 노래를 하며
\end{kfChoices}
\end{kfQBlock}
\begin{kfQBlock}{9}{}
\kfQStem{'빼앗긴 들'이라는 표현에 대한 설명으로 알맞은 것은 무엇인가?}
\begin{kfChoices}
\kfChoice 직유법 - '들'이라는 대상을 '빼앗긴 것'에 직접적으로 비유한 표현이다.
\kfChoice 은유법 - 들은 결국 적에게 모두 빼앗긴 우리의 땅이라는 뜻을 담고 있다.
\kfChoice 의인법 - '들'이 사람처럼 슬퍼하며 눈물을 흘리고 있다는 뜻이다.
\kfChoice 대유법 - 국토의 일부인 '들'을 통해 우리나라 국토 전체를 나타낸다.
\kfChoice 반어법 - 실제로는 빼앗기지 않았는데 빼앗겼다고 일부러 반대로 표현했다.
\end{kfChoices}
\end{kfQBlock}
\begin{kfQBlock}{10}{}
\kfQStem{비유법에 대한 설명으로 알맞지 \uline{않은} 것은 무엇인가?}
\begin{kfChoices}
\kfChoice 비유법은 표현을 단순하고 건조하게 만든다.
\kfChoice 은유법은 'A는 B이다'의 형태를 띤다.
\kfChoice 의인법은 사물을 사람처럼 표현하여 생동감을 준다.
\kfChoice 대유법은 부분으로 전체를 나타낸다.
\kfChoice 직유법은 비교하는 말이 직접 드러나 이해하기 쉽다.
\end{kfChoices}
\end{kfQBlock}
\begin{kfQBlock}{1}{}
\kfQStem{표현하고 싶은 대상을 그와 비슷한 다른 대상에 빗대어 생생하게 나타내는 표현 방법을 무엇이라고 하는가?}
\kfAnswerLines{1}
\end{kfQBlock}
\begin{kfQBlock}{2}{}
\kfQStem{비유법에서 진짜 말하고 싶은 대상, 즉 '논길'이나 '흙'을 무엇이라고 하는가?}
\kfAnswerLines{1}
\end{kfQBlock}
\begin{kfQBlock}{3}{}
\kfQStem{비유법에서 빗대어 표현하는 대상, 즉 '가르마'나 '살찐 가슴'을 무엇이라고 하는가?}
\kfAnswerLines{1}
\end{kfQBlock}
\begin{kfQBlock}{4}{}
\kfQStem{'같이', '처럼', '듯이' 등을 사용하여 원관념과 보조관념을 직접 연결하는 비유법은 무엇인가?}
\kfAnswerLines{1}
\end{kfQBlock}
\begin{kfQBlock}{5}{}
\kfQStem{"살찐 가슴처럼 부드러운 이 흙"에서 원관념은 무엇인가?}
\kfAnswerLines{1}
\end{kfQBlock}
\begin{kfQBlock}{6}{}
\kfQStem{'A는 B이다'의 형태로 원관념과 보조관념을 간접적으로 연결하는 비유법은 무엇인가?}
\kfAnswerLines{1}
\end{kfQBlock}
\begin{kfQBlock}{7}{}
\kfQStem{"내 마음은 호수요"는 '마음'을 '호수'처럼 넓고 고요하다는 의미로 표현한 것입니다. 이러한 비유법을 무엇이라고 하는가?}
\kfAnswerLines{1}
\end{kfQBlock}
\begin{kfQBlock}{1}{}
\kfQStem{'직유법'과 '은유법'의 차이점을 <조건>에 맞게 서술하시오.}
\kfAnswerNote{답안}{4}
\end{kfQBlock}


\kfStartLit{문학 영역 — 작품}


\kfSecHeader{kfAreaLiterature}{지문}{작품}
\kfPassageNote{작품}{%
지금은 남의 땅 \textbackslash{}- 빼앗긴 들에도 봄은 오는가? \\[14pt]
나는 온몸에 따뜻한 햇살을 받고 \\[4pt]
푸른 하늘과 푸른 들판이 만나는 곳으로 \\[4pt]
가르마 같은 논길을 따라 꿈속을 걷듯 걸어만 간다. \\[14pt]
입을 다물고 있는 하늘아, 들판아 \\[4pt]
내 마음에는 내가 혼자 온 것 같지가 않구나. \\[4pt]
네가 나를 이끌었느냐, 누가 나를 불렀느냐? 답답하구나, 대답을 해 다오. \\[14pt]
바람은 내 귀에 속삭이며 \\[4pt]
한 발짝도 멈추지 마라, 옷자락을 흔들고 \\[4pt]
종달새는 울타리 너머에서 아가씨처럼 구름 뒤에서 반갑게 웃네. \\[14pt]
고맙게 잘 자란 보리밭아 \\[4pt]
간밤 자정이 넘어서 내린 고운 비로 \\[4pt]
너는 삼단 머리처럼 깨끗이 씻었구나. 내 머리까지 가벼워진다. \\[14pt]
혼자라도 힘차게 가자 \\[4pt]
마른 논을 끼고 도는 착한 개울이 \\[4pt]
아기를 달래는 노래를 하며 혼자서 어깨춤을 추며 흘러가네. \\[14pt]
나비야, 제비야, 장난치지 마라 \\[4pt]
맨드라미와 들꽃에게도 인사를 해야지 \\[4pt]
피마자기름을 바른 사람이 농사짓던 그 들판을 다 보고 싶다. \\[14pt]
내 손에 호미를 쥐어 다오 \\[4pt]
살찐 가슴처럼 부드러운 이 흙을 \\[4pt]
발목이 아프도록 밟아도 보고 좋은 땀도 흘리고 싶다. \\[14pt]
강가에 나온 아이처럼 \\[4pt]
시간도 모르고 끝도 없이 떠도는 내 마음아 \\[4pt]
무엇을 찾느냐, 어디로 가느냐? 우스워라, 대답을 해 보거라. \\[14pt]
나는 온몸에 풀냄새를 맡으며 \\[4pt]
푸른 웃음과 푸른 슬픔이 섞인 사이로 \\[4pt]
다리를 절며 하루를 걷는다. 아마도 봄 신령이 내게 내려왔나 보다. \\[14pt]
그러나 지금은 \textbackslash{}- 들을 빼앗겨 봄까지 빼앗기겠네.
}


\kfSecHeader{kfAreaLiterature}{활동}{문학 문장 독해}
\noindent\textit{\small 작품 속 문장을 자세히 읽고, 문장별로 주어진 물음에 답해 보자.}\par\medskip
\begin{kfSentenceUnit}{1}{"바람은 내 귀에 속삭이며 한 발짝도 멈추지 마라, 옷자락을 흔들고 종달새는 울타리 너머에서 아가씨처럼 구름 뒤에서 반갑게 웃네."}
\kfSentQ{1.}{다음 문장에서 서술어를 기준으로 각각의 주어는 무엇인가? (1) 서술어 '속삭이며'의 주어는 무엇인가?}
\begin{kfChoices}
\kfChoice 바람은
\kfChoice 종달새는
\kfChoice 구름
\kfChoice 아가씨
\end{kfChoices}
\kfSentQ{1.}{(2) 서술어 '흔들고'의 주어는 무엇인가?}
\begin{kfChoices}
\kfChoice 바람은
\kfChoice 종달새는
\kfChoice 내 귀
\kfChoice 울타리
\end{kfChoices}
\kfSentQ{1.}{(3) 서술어 '웃네'의 주어는 무엇인가?}
\begin{kfChoices}
\kfChoice 바람은
\kfChoice 종달새는
\kfChoice 울타리
\kfChoice 구름
\end{kfChoices}
\end{kfSentenceUnit}

\begin{kfSentenceUnit}{2}{"나는 온몸에 따뜻한 햇살을 받고 푸른 하늘과 푸른 들판이 만나는 곳으로 가르마 같은 논길을 따라 꿈속을 걷듯 걸어만 간다."}
\kfSentQ{2.}{다음 문장에서 서술어를 기준으로 각각의 주어는 무엇인가? (1) 서술어 '받고'의 주어는 무엇인가?}
\begin{kfChoices}
\kfChoice 나는
\kfChoice 햇살을
\kfChoice 하늘
\kfChoice 들판
\end{kfChoices}
\kfSentQ{2.}{(2) 서술어 '만나는'의 주어는 무엇인가?}
\begin{kfChoices}
\kfChoice 나는
\kfChoice 햇살
\kfChoice 푸른 하늘과 푸른 들판이
\kfChoice 논길
\end{kfChoices}
\kfSentQ{2.}{(3) 서술어 '걸어만 간다'의 주어는 무엇인가?}
\begin{kfChoices}
\kfChoice 나는
\kfChoice 햇살
\kfChoice 하늘
\kfChoice 들판
\end{kfChoices}
\end{kfSentenceUnit}

\begin{kfSentenceUnit}{3}{"고맙게 잘 자란 보리밭아 간밤 자정이 넘어서 내린 고운 비로 너는 삼단 머리처럼 깨끗이 씻었구나."}
\kfSentQ{3.}{'너는'이 가리키는 대상은 무엇인가?}
\begin{kfChoices}
\kfChoice 보리밭
\kfChoice 고운 비
\kfChoice 삼단 머리
\kfChoice 자정
\end{kfChoices}
\end{kfSentenceUnit}

\begin{kfSentenceUnit}{4}{"입을 다물고 있는 하늘아, 들판아 내 마음에는 내가 혼자 온 것 같지가 않구나. 네가 나를 이끌었느냐, 누가 나를 불렀느냐?"}
\kfSentQ{4.}{'네가'가 가리키는 대상은 무엇인가?}
\begin{kfChoices}
\kfChoice 나
\kfChoice 마음
\kfChoice 하늘과 들판
\kfChoice 누
\end{kfChoices}
\end{kfSentenceUnit}

\kfSecHeader{kfAreaLiterature}{활동}{해설 학습}
\noindent\textit{\small 작품 해설을 단원별로 읽고, 각 단원에 딸린 물음에 답해 보자.}\par\medskip
\par\noindent\textbf{\color{kfPrimary} 해설 단원 1}\par\smallskip
\kfPassageNote{해설 1}{%
「빼앗긴 들에도 봄은 오는가」는 일제강점기에 쓰인 대표적인 저항시이다. 이 시는 직접적으로 일제를 비판하지는 않지만, 상징적인 방법으로 나라를 잃은 슬픔과 되찾고 싶은 의지를 드러낸다.

화자는 "내 손에 호미를 쥐어 다오"라며 이 땅에서 직접 농사짓고 싶다는 소망을 표현한다. 이는 단순히 농사를 짓고 싶다는 뜻이 아니라, 이 땅의 진정한 주인이 되고 싶다는 의미이다. "좋은 땀도 흘리고 싶다"는 표현도 조국을 위해 기꺼이 노력하겠다는 의지를 담고 있다.

하지만 시의 마지막에서 화자는 "봄까지 빼앗기겠네"라며 절망적인 현실을 인정한다. 자연의 봄마저도 제대로 누릴 수 없게 될 것이라는 두려움을 표현한 것으로, 일제강점기의 암울한 현실을 보여준다. 그럼에도 불구하고 이 시 자체가 우리말로 우리의 정서를 표현했다는 것 자체가 저항 정신의 표현이라고 할 수 있다.
}

\begin{kfQBlock}{1}{문제}
\kfQStem{이 시는 일제강점기에 쓰인 대표적인 어떤 종류의 시인가?}
\kfAnswerLines{1}
\end{kfQBlock}

\begin{kfQBlock}{2}{문제}
\kfQStem{화자가 "좋은 땀도 흘리고 싶다"고 한 말에 담긴 의지는 무엇인가?}
\kfAnswerLines{1}
\end{kfQBlock}

\begin{kfQBlock}{3}{문제}
\kfQStem{시의 마지막 구절 "봄까지 빼앗기겠네"에서 드러나는 화자의 현실에 대한 인식은 무엇인가?}
\kfAnswerLines{1}
\end{kfQBlock}

\begin{kfQBlock}{4}{문제}
\kfQStem{화자가 이 땅의 주인이 되고 싶다는 소망을 표현하기 위해 손에 쥐어달라고 한 도구는 무엇인가?}
\kfAnswerLines{1}
\end{kfQBlock}

\begin{kfQBlock}{5}{문제}
\kfQStem{화자는 들판에 이어 무엇까지 빼앗길 것이라고 두려워하고 있는가?}
\kfAnswerLines{1}
\end{kfQBlock}

\begin{kfQBlock}{6}{문제}
\kfQStem{이 시의 화자가 절망 속에서도 우리말로 우리의 정서를 표현한 것 자체가 무엇의 표현이라고 할 수 있는가?}
\kfAnswerLines{1}
\end{kfQBlock}

\par\noindent\textbf{\color{kfPrimary} 해설 단원 2}\par\smallskip
\kfPassageNote{해설 2}{%
「빼앗긴 들에도 봄은 오는가」는 첫 연과 마지막 연이 질문과 대답으로 호응하는 수미상관 구조가 특징적인 작품이다. 수미상관이란 시나 소설에서 시작과 끝을 유사하게 구성해 주제를 강조하거나 리듬을 형성하는 기법이다. 보통은 첫 부분과 마지막 부분이 똑같은 표현을 반복하는데, 이 작품에서는 조금 다르게 질문과 대답의 형식으로 수미상관을 이루고 있다.

첫 연에서 "빼앗긴 들에도 봄은 오는가?"라고 물었다면, 마지막 연에서는 "들을 빼앗겨 봄까지 빼앗기겠네"라고 절망적으로 답하고 있다. 똑같은 말은 아니지만 '빼앗긴 들'과 '봄'이라는 핵심 단어가 반복되면서 질문에 대한 대답 형식으로 연결되어 있어서 수미상관의 효과를 내고 있다.

시 전체의 구조를 보면 화자의 감정이 변화하는 과정이 잘 드러난다. 처음에는 빼앗긴 현실을 인식하고 절망하다가, 중간 부분에서는 봄의 아름다운 풍경에 빠져들며 기쁨을 느낀다. 하지만 마지막에는 다시 현실을 깨닫고 더 깊은 절망에 빠지게 된다.

"그러나 지금은"이라는 표현은 시상이 전환되는 중요한 부분이다. 아름다운 봄에 도취되어 있던 화자가 갑자기 차가운 현실로 돌아오는 순간을 보여준다. 이런 구조를 통해 일제강점기 지식인의 고뇌와 내적 갈등을 효과적으로 드러내고 있다.
}

\begin{kfQBlock}{1}{문제}
\kfQStem{이 시는 첫 연과 마지막 연이 똑같은 문장으로 반복되지는 않지만, ‘빼앗긴 들’과 ‘봄’이라는 핵심 시어를 중심으로 질문과 대답의 관계를 이룹니다. 이처럼 시의 시작과 끝을 유사하게 구성하여 주제를 강조하는 기법을 무엇이라고 하는가?}
\kfAnswerLines{1}
\end{kfQBlock}

\begin{kfQBlock}{2}{문제}
\kfQStem{아름다운 봄에 대한 황홀경에서 차가운 현실로 돌아오게 만드는 시상 전환 부분의 시작 시구는 무엇인가?}
\kfAnswerLines{1}
\end{kfQBlock}

\begin{kfQBlock}{3}{문제}
\kfQStem{이 시의 구조는 일제강점기 어떤 계층의 인물이 겪는 고뇌와 내적 갈등을 효과적으로 보여주는가?}
\kfAnswerLines{1}
\end{kfQBlock}

\begin{kfQBlock}{4}{문제}
\kfQStem{첫 연에서 화자가 던진 질문은 무엇인가?}
\kfAnswerLines{1}
\end{kfQBlock}

\begin{kfQBlock}{5}{문제}
\kfQStem{화자가 마지막 연에서 내린 절망적인 대답은 무엇인가?}
\kfAnswerLines{1}
\end{kfQBlock}

\begin{kfQBlock}{6}{문제}
\kfQStem{첫 연과 마지막 연은 어떤 형식으로 서로 호응하고 있는가?}
\kfAnswerLines{1}
\end{kfQBlock}

\par\noindent\textbf{\color{kfPrimary} 해설 단원 3}\par\smallskip
\kfPassageNote{해설 3}{%
이 시에서 가장 중요한 상징은 '빼앗긴 들'과 '봄'이다. '빼앗긴 들'은 일제에 나라를 빼앗긴 우리나라를 의미하고, '봄'은 자연의 계절이면서 동시에 조국 광복에 대한 희망을 상징한다.

시 전체에는 우리나라의 전통적이고 정겨운 풍경들이 많이 등장한다. "가르마 같은 논길"은 논 사이의 좁은 길을 머리 가르마에 비유한 것이고, "삼단 머리"는 옛날 여성들의 머리 모양을 말한다. "피마자기름을 바른 사람"도 전통적인 우리나라 여성들의 모습을 그린 것이다.

"살찐 가슴처럼 부드러운 흙", "아기를 달래는 노래" 같은 표현들은 우리 땅을 어머니처럼 따뜻하고 풍요로운 존재로 그리고 있다. 이런 다양한 상징들이 어우러져 단순한 자연 풍경이 아닌, 조국에 대한 깊은 사랑과 그리움을 표현하고 있다.
}

\begin{kfQBlock}{1}{문제}
\kfQStem{이 시에서 '빼앗긴 들'이 상징하는 것은 무엇인가?}
\kfAnswerLines{1}
\end{kfQBlock}

\begin{kfQBlock}{2}{문제}
\kfQStem{이 시에서 계절로서의 의미와 함께 '조국 광복에 대한 희망'을 상징하는 시어는 무엇인가?}
\kfAnswerLines{1}
\end{kfQBlock}

\begin{kfQBlock}{3}{문제}
\kfQStem{논 사이의 좁은 길을 머리카락을 나눈 모양에 빗대어 표현한 시어는 무엇인가?}
\kfAnswerLines{1}
\end{kfQBlock}

\begin{kfQBlock}{4}{문제}
\kfQStem{"살찐 가슴처럼 부드러운 흙"이라는 표현은 우리 땅을 어떤 존재로 그리고 있는가?}
\kfAnswerLines{1}
\end{kfQBlock}

\begin{kfQBlock}{5}{문제}
\kfQStem{이 시에 나오는 "가르마 같은 논길"과 "삼단 머리"는 어떤 존재의 이미지를 떠올리게 하는가?}
\kfAnswerLines{1}
\end{kfQBlock}

\par\noindent\textbf{\color{kfPrimary} 해설 단원 4}\par\smallskip
\kfPassageNote{해설 4}{%
「빼앗긴 들에도 봄은 오는가」에서 주목할 만한 표현 기법 중 하나는 대유법이다. 대유법이란 전체를 부분으로, 또는 부분을 전체로 표현하는 방법이다.

이 시에서 '들'은 우리나라 전체를 의미하는 대유법적 표현이다. 실제로는 나라 전체가 빼앗겼지만, 그 중에서도 '들판'이라는 부분을 들어서 전체 국토를 나타내고 있다. 들판은 우리 민족이 대대로 농사지으며 살아온 삶의 터전이자, 가장 친근하고 정겨운 공간이기 때문에 조국을 상징하기에 적합한 소재이다.

이처럼 대유법을 사용하면 추상적이고 큰 개념을 구체적이고 친근한 이미지로 표현할 수 있어서, 독자들이 더 생생하게 느끼고 깊이 공감할 수 있게 된다.
}

\begin{kfQBlock}{1}{문제}
\kfQStem{전체를 부분으로, 또는 부분을 전체로 표현하는 방법을 무엇이라고 하는가?}
\kfAnswerLines{1}
\end{kfQBlock}

\begin{kfQBlock}{2}{문제}
\kfQStem{'들'이라는 부분은 실제로는 어떤 전체를 나타내고 있는가?}
\kfAnswerLines{1}
\end{kfQBlock}

\begin{kfQBlock}{3}{문제}
\kfQStem{들판이 조국을 상징하기에 적합한 이유 중 하나는, 우리 민족이 대대로 살아온 무엇이기 때문인가?}
\kfAnswerLines{1}
\end{kfQBlock}

\begin{kfQBlock}{4}{문제}
\kfQStem{이 시에서 '들판'이라는 부분으로 표현된 전체 국토가 어떤 상황에 처해있다고 설명되었는가?}
\kfAnswerLines{1}
\end{kfQBlock}

\kfSecHeader{kfAreaLiterature}{문제}{}

\begin{kfQBlock}{1}{}
\kfQStem{이 시의 전체적인 내용과 흐름을 가장 잘 설명한 것은 무엇인가?}
\begin{kfChoices}
\kfChoice 봄의 아름다움에 기뻐하다가, 봄이 짧다는 사실에 슬퍼하고 있다.
\kfChoice 나라를 빼앗긴 슬픔을 느끼다가, 봄의 풍경을 보며 희망을 되찾고 있다.
\kfChoice 자연의 아름다움에 흠뻑 빠져 기뻐하다가, 이 땅이 빼앗긴 현실을 깨닫고 절망하고 있다.
\kfChoice 고향의 봄을 그리워하다가, 다시 돌아갈 수 없다는 사실에 체념하고 있다.
\kfChoice 봄을 즐기는 사람들을 부러워하다가, 자신도 그들처럼 되고 싶다고 다짐하고 있다.
\end{kfChoices}
\end{kfQBlock}
\begin{kfQBlock}{2}{}
\kfQStem{말하는 이가 봄의 들판을 걸으며 느끼는 주된 감정으로 알맞지 \uline{않은} 것은 무엇인가?}
\begin{kfChoices}
\kfChoice 종달새 소리와 개울의 춤을 보며 느끼는 흥겨움
\kfChoice 땀 흘려 일하고 싶은 국토에 대한 애정
\kfChoice 푸른 하늘과 들판을 보며 느끼는 아름다움
\kfChoice 삶을 포기하고 싶은 절망감과 무기력함
\kfChoice 자신의 처지를 모르고 마냥 즐거워하는 것에 대한 부끄러움
\end{kfChoices}
\end{kfQBlock}
\begin{kfQBlock}{3}{}
\kfQStem{이 시에서 ‘빼앗긴 들’이 상징하는 의미로 적절하지 \uline{않은} 것은 무엇인가?}
\begin{kfChoices}
\kfChoice 다른 사람에게 팔려나간 한 개인의 사유지
\kfChoice 자유와 주권을 잃은 식민지 조선의 현실
\kfChoice 우리 민족이 되찾아야 할 삶의 터전
\kfChoice 봄이 와도 진정한 봄을 누릴 수 없는 억압된 공간
\kfChoice 일제에게 주권을 빼앗긴 우리의 국토
\end{kfChoices}
\end{kfQBlock}
\begin{kfQBlock}{4}{}
\kfQStem{“푸른 웃음과 푸른 슬픔이 섞인 사이로”라는 구절에 담긴 말하는 이의 심정으로 적절하지 \uline{않은} 것은 무엇인가?}
\begin{kfChoices}
\kfChoice 기쁜 척, 슬픈 척하며 자신의 진짜 감정을 일부러 속이는 거짓된 마음
\kfChoice 자연의 생명력이 주는 환희와 식민지 현실의 비애가 뒤섞인 양가감정
\kfChoice 봄을 누리고 싶지만 ‘빼앗긴 들’이라는 현실에 마음 깊이 자리한 비통함
\kfChoice 봄을 만끽하면서도 끝내 잊을 수 없는 민족적 슬픔
\kfChoice 봄의 아름다움에 빠진 기쁨과 나라를 빼앗긴 슬픔이 함께 어우러진 복잡한 마음
\end{kfChoices}
\end{kfQBlock}
\begin{kfQBlock}{5}{}
\kfQStem{말하는 이가 봄의 아름다움에 흠뻑 빠져 있다가, 다시 현실을 깨닫고 스스로를 “우스워라”라고 말하는 이유는 무엇인가?}
\begin{kfChoices}
\kfChoice 길을 잃어 엉뚱한 들길로 잘못 들어선 자신이 어이없어서
\kfChoice 종달새가 마치 아가씨처럼 구는 모습이 너무 우스워서
\kfChoice 나라를 빼앗긴 비참한 현실을 잊고 즐거워한 자신이 부끄러워서
\kfChoice 자신의 노래 솜씨가 형편없음을 새삼 깨달아서
\kfChoice 친구들과의 약속을 잊고 혼자 들판에 나와 있어서
\end{kfChoices}
\end{kfQBlock}
\begin{kfQBlock}{6}{}
\kfQStem{이 시의 첫 구절과 마지막 구절에 대한 설명으로 가장 알맞은 것은 무엇인가?}
\begin{kfChoices}
\kfChoice 질문으로 시작해서 희망적인 대답으로 끝나고 있다.
\kfChoice 슬픈 질문으로 시작해서 더욱 절망적인 대답으로 끝나고 있다.
\kfChoice 슬픈 질문으로 시작해서 긍정적인 깨달음으로 끝나고 있다.
\kfChoice 간절한 소망으로 시작해서 그 소망이 이루어지며 끝나고 있다.
\kfChoice 기쁜 다짐으로 시작해서 슬픈 후회로 끝나고 있다.
\end{kfChoices}
\end{kfQBlock}
\begin{kfQBlock}{7}{}
\kfQStem{이 시에서 대상에 대한 말하는 이의 정서 및 태도에 대한 설명으로 알맞지 \uline{않은} 것은 무엇인가?}
\begin{kfChoices}
\kfChoice 하늘과 들판을 향해 왜 침묵하냐며 원망을 드러낸다.
\kfChoice 잘 자란 보리밭을 보며 고마움을 느낀다.
\kfChoice 마른 논을 적셔주는 개울을 '착하다'고 말한다.
\kfChoice 흙을 ‘살찐 가슴’에 비유하며 부드럽고 풍요로운 대상으로 여긴다.
\kfChoice 종달새를 '아씨'에 비유하며 반가움을 표현한다.
\end{kfChoices}
\end{kfQBlock}
\begin{kfQBlock}{8}{}
\kfQStem{말하는 이가 “내 손에 호미를 쥐어 다오”라고 말한 가장 큰 이유는 무엇인가?}
\begin{kfChoices}
\kfChoice 농사를 직접 지어서 많은 돈을 벌어보고 싶어서
\kfChoice 가난한 농부들이 겪는 힘든 일상을 체험해 보고 싶어서
\kfChoice 빼앗긴 우리 땅의 주인이 되어 직접 땀 흘리며 가꾸고 싶어서
\kfChoice 오래 써서 낡고 부서진 헌 호미를 새것으로 바꾸고 싶어서
\kfChoice 호미를 들어 보이며 일본에 저항하는 용기를 드러내고 싶어서
\end{kfChoices}
\end{kfQBlock}
\begin{kfQBlock}{1}{}
\kfQStem{“푸른 웃음과 푸른 슬픔”이라는 표현에 담긴 말하는 이의 복잡한 마음을 <조건>에 맞게 서술하시오.}
\kfAnswerNote{답안}{4}
\end{kfQBlock}
\begin{kfQBlock}{2}{}
\kfQStem{이 시의 첫 연과 마지막 연은 서로 질문하고 답하는 관계에 있습니다. 이러한 구조가 주는 효과를 서술하시오.}
\kfAnswerNote{답안}{4}
\end{kfQBlock}


\kfStartNonlit{비문학 영역 — [인문] 내 마음의 다섯 가지 욕구, 현실요법}


\kfSecHeader{kfAreaNonlit}{지문}{[인문] 내 마음의 다섯 가지 욕구, 현실요법}
\kfPassageNote{[인문] 내 마음의 다섯 가지 욕구, 현실요법}{%
사람은 누구나 마음속에 바라는 것이 있다. 현실요법이라는 상담 방법에서는, 사람의 모든 행동이 이런 욕구를 이루기 위해 스스로 선택하는 것이라고 본다. 그런데 어떤 선택이 문제가 되면, 마음속 욕구들의 균형을 맞춰서 새로운 선택을 하는 것이 중요하다고 말한다.

현실요법에서는 사람에게 다섯 가지 기본적인 욕구가 있다고 설명한다. 첫째는 살고 싶은 마음이다. 건강을 지키고, 안전하게 살고 싶어 하는 마음이다. 이 마음이 강한 사람은 규칙을 잘 지키고, 필요한 것을 아끼며 모으는 것을 좋아한다. 둘째는 사랑하고 싶은 마음이다. 다른 사람과 함께 지내고, 사랑을 주고받고 싶어 하는 마음이다. 이런 사람은 다른 사람을 잘 도와주지만, 사랑을 받지 못하면 상처를 받기 쉽다. 셋째는 인정받고 싶은 마음이다. 다른 사람에게 잘한다는 말을 듣고 싶고, 스스로 뿌듯함을 느끼고 싶은 마음이다. 이런 사람은 공부나 일에 열심히 하고, 자기 생각을 잘 말하는 편이다. 넷째는 자유롭고 싶은 마음이다. 누구에게도 간섭받지 않고, 스스로 선택하고 싶은 마음이다. 이 마음이 강한 사람은 혼자 있는 걸 좋아하고, 다른 사람과 너무 가까워지는 것을 불편해할 수도 있다. 다섯째는 즐겁게 살고 싶은 마음이다. 새로운 것을 배우거나, 놀면서 재미를 느끼고 싶어 하는 마음이다. 이런 사람은 취미 활동을 즐기고, 항상 밝고 호기심이 많다.

사람마다 이 다섯 가지 마음의 크기는 다르다. 그래서 행동하는 방식도 다르고, 갈등이 생기기도 한다. 예를 들어, 사랑하고 싶은 마음은 크지만 인정받고 싶은 마음은 약한 사람은, 다른 사람의 부탁을 거절하지 못해 힘들 수 있다. 이럴 때는 인정받고 싶은 마음을 키워서, 자기 생각을 솔직하게 말할 수 있도록 도와준다. 또 어떤 사람은 자유롭고 싶은 마음과 인정받고 싶은 마음이 둘 다 강한 경우도 있다. 이럴 때는 자기 생각만 고집하면 친구들과 자주 다툴 수 있다. 이럴 때는 다른 사람의 생각도 기꺼이 듣고, 마음을 헤아리는 연습이 필요하다.

현실요법에서는 자기 욕구만 채우려고 하지 않고, 다른 사람의 욕구도 방해하지 않으면서 좋은 선택을 하도록 돕는다. 사람은 남이 시키는 대로만 사는 존재가 아니라, 스스로 선택하고 행동할 수 있는 존재라는 것이 현실요법의 가장 중요한 생각이다. 요즘에는 이 상담 방법이 많은 사람들에게 관심을 받고, 여러 심리 상담에서 넓게 사용되고 있다.
}


\kfSecHeader{kfAreaNonlit}{활동}{어휘 학습}
\noindent\textit{\small 다음 표의 '의미'를 읽고, 그에 해당하는 '어휘'를 지문에서 찾아 적으시오.}\par\medskip
\begin{kfVocab}
\kfVocabItem{1}{어떤 것을 이루거나 얻고 싶어 하는 마음}
\kfVocabItem{2}{여러 요소가 어느 한쪽으로 기울지 않고 고른 상태}
\kfVocabItem{3}{다른 사람에게 능력을 칭찬받고 싶어 하는 마음}
\kfVocabItem{4}{다른 사람의 일에 끼어들어 이래라저래라 하는 것}
\kfVocabItem{5}{생각이나 의견이 달라 서로 부딪치는 것}
\kfVocabItem{6}{즐거운 마음으로}
\kfVocabItem{7}{상대방의 마음이나 상황을 짐작하여 살피는 것}
\end{kfVocab}

\kfSecHeader{kfAreaNonlit}{활동}{비문학 문장 독해}
\noindent\textit{\small 지문을 문장 단위로 정독하면서, 문장별로 주어진 물음에 답해 보자.}\par\medskip
\begin{kfSentenceUnit}{1}{사람은 누구나 마음속에 바라는 것이 있다.}
\kfSentQ{1.}{누구나 마음속에 무엇이 있는가?}
\begin{kfChoices}
\kfChoice 미워하는 것
\kfChoice 두려워하는 것
\kfChoice 바라는 것
\kfChoice 후회하는 것
\end{kfChoices}
\end{kfSentenceUnit}

\begin{kfSentenceUnit}{2}{현실요법이라는 상담 방법에서는, 사람의 모든 행동이 이런 욕구를 이루기 위해 스스로 선택하는 것이라고 본다.}
\kfSentQ{2.}{'현실요법'은 무엇인가?}
\begin{kfChoices}
\kfChoice 운동 방법
\kfChoice 식이 요법
\kfChoice 상담 방법
\kfChoice 명상 방법
\end{kfChoices}
\kfSentQ{3.}{'이런 욕구'가 가리키는 것은 무엇인가?}
\begin{kfChoices}
\kfChoice 마음속에 바라는 것
\kfChoice 타인의 기대
\kfChoice 사회적 규칙
\kfChoice 신체적 고통
\end{kfChoices}
\kfSentQ{4.}{현실요법에서는 사람의 모든 행동을 무엇이라고 보는가?}
\begin{kfChoices}
\kfChoice 본능에 의한 무의식적 반응
\kfChoice 욕구를 이루기 위해 스스로 선택하는 것
\kfChoice 환경에 의해 결정되는 수동적 결과
\kfChoice 타인의 강요에 의한 행동
\end{kfChoices}
\end{kfSentenceUnit}

\begin{kfSentenceUnit}{3}{그런데 어떤 선택이 문제가 되면, 마음속 욕구들의 균형을 맞춰서 새로운 선택을 하는 것이 중요하다고 말한다.}
\kfSentQ{5.}{어떤 경우에 새로운 선택이 필요한가?}
\begin{kfChoices}
\kfChoice 모든 일이 순조로울 때
\kfChoice 욕구가 완전히 충족되었을 때
\kfChoice 어떤 선택이 문제가 될 때
\kfChoice 타인이 칭찬할 때
\end{kfChoices}
\kfSentQ{6.}{현실요법에서 문제가 되는 선택을 해결하는 방법은 무엇인가?}
\begin{kfChoices}
\kfChoice 욕구를 억누르고 참는 것
\kfChoice 마음속 욕구들의 균형을 맞추는 것
\kfChoice 상대방의 요구를 무조건 들어주는 것
\kfChoice 욕구를 포기하는 것
\end{kfChoices}
\end{kfSentenceUnit}

\begin{kfSentenceUnit}{4}{현실요법에서는 사람에게 다섯 가지 기본적인 욕구가 있다고 설명한다.}
\kfSentQ{7.}{현실요법에 따르면 사람에게는 몇 가지 기본적인 욕구 있는가?}
\begin{kfChoices}
\kfChoice 세 가지
\kfChoice 네 가지
\kfChoice 다섯 가지
\kfChoice 여섯 가지
\end{kfChoices}
\kfSentQ{8.}{이 문장을 통해 앞으로 어떤 내용이 나올 것을 예상할 수 있는가?}
\begin{kfChoices}
\kfChoice 현실요법의 역사와 유래
\kfChoice 상담가가 되는 방법
\kfChoice 다섯 가지 기본적인 욕구에 대한 구체적인 설명
\kfChoice 심리 상담의 종류와 비용
\end{kfChoices}
\end{kfSentenceUnit}

\begin{kfSentenceUnit}{5}{첫째는 살고 싶은 마음이다.}
\kfSentQ{9.}{첫 번째 기본적인 욕구는 무엇인가?}
\begin{kfChoices}
\kfChoice 사랑하고 싶은 마음
\kfChoice 살고 싶은 마음
\kfChoice 인정받고 싶은 마음
\kfChoice 즐겁게 살고 싶은 마음
\end{kfChoices}
\end{kfSentenceUnit}

\begin{kfSentenceUnit}{6}{건강을 지키고, 안전하게 살고 싶어 하는 마음이다.}
\kfSentQ{10.}{'살고 싶은 마음'이란 구체적으로 어떤 마음을 말하는가?}
\begin{kfChoices}
\kfChoice 건강을 지키고 안전하게 살고 싶은 마음
\kfChoice 다른 사람과 잘 지내고 싶은 마음
\kfChoice 남에게 칭찬받고 싶은 마음
\kfChoice 새롭고 재미있는 것을 찾는 마음
\end{kfChoices}
\end{kfSentenceUnit}

\begin{kfSentenceUnit}{7}{이 마음이 강한 사람은 규칙을 잘 지키고, 필요한 것을 아끼며 모으는 것을 좋아한다.}
\kfSentQ{11.}{'이 마음'이 가리키는 것은 무엇인가?}
\begin{kfChoices}
\kfChoice 사랑하고 싶은 마음
\kfChoice 인정받고 싶은 마음
\kfChoice 살고 싶은 마음
\kfChoice 자유롭고 싶은 마음
\end{kfChoices}
\kfSentQ{12.}{'살고 싶은 마음'이 강한 사람의 특징으로 알맞은 것은?}
\begin{kfChoices}
\kfChoice 다른 사람을 잘 도와주고, 함께 지내는 것을 좋아한다.
\kfChoice 규칙을 잘 지키고, 필요한 것을 아끼며 모으는 것을 좋아한다.
\kfChoice 공부나 일에 열심히 하고, 자기 생각을 잘 말하는 편이다.
\kfChoice 혼자 있는 것을 좋아하고, 간섭받는 것을 싫어한다.
\end{kfChoices}
\end{kfSentenceUnit}

\begin{kfSentenceUnit}{8}{둘째는 사랑하고 싶은 마음이다.}
\kfSentQ{13.}{두 번째 기본적인 욕구는 무엇인가?}
\begin{kfChoices}
\kfChoice 사랑하고 싶은 마음
\kfChoice 인정받고 싶은 마음
\kfChoice 자유롭고 싶은 마음
\kfChoice 살고 싶은 마음
\end{kfChoices}
\end{kfSentenceUnit}

\begin{kfSentenceUnit}{9}{다른 사람과 함께 지내고, 사랑을 주고받고 싶어 하는 마음이다.}
\kfSentQ{14.}{'사랑하고 싶은 마음'에 대한 설명으로 알맞은 것은?}
\begin{kfChoices}
\kfChoice 건강을 지키고 안전하게 살고 싶어 하는 마음
\kfChoice 다른 사람과 함께 지내며 사랑을 주고받고 싶어 하는 마음
\kfChoice 남에게 인정받고 스스로 뿌듯함을 느끼고 싶은 마음
\kfChoice 간섭받지 않고 스스로 선택하고 싶은 마음
\end{kfChoices}
\end{kfSentenceUnit}

\begin{kfSentenceUnit}{10}{이런 사람은 다른 사람을 잘 도와주지만, 사랑을 받지 못하면 상처를 받기 쉽다.}
\kfSentQ{15.}{'이런 사람'이 가리키는 것은 누구인가?}
\begin{kfChoices}
\kfChoice 살고 싶은 마음이 강한 사람
\kfChoice 사랑하고 싶은 마음이 강한 사람
\kfChoice 인정받고 싶은 마음이 강한 사람
\kfChoice 자유롭고 싶은 마음이 강한 사람
\end{kfChoices}
\kfSentQ{16.}{사랑하고 싶은 마음이 강한 사람의 장점으로 알맞은 것은?}
\begin{kfChoices}
\kfChoice 규칙을 철저하게 잘 지킨다.
\kfChoice 자기 생각을 분명하게 말한다.
\kfChoice 다른 사람을 잘 도와준다.
\kfChoice 혼자서도 외로움을 느끼지 않는다.
\end{kfChoices}
\kfSentQ{17.}{사랑하고 싶은 마음이 강한 사람의 특징으로 알맞은 것은?}
\begin{kfChoices}
\kfChoice 다른 사람과 너무 가까워지는 것을 불편해한다.
\kfChoice 친구들과 자주 다투고 갈등을 일으킨다.
\kfChoice 자기 물건을 남에게 빌려주지 않는다.
\kfChoice 사랑을 받지 못하면 상처를 받기 쉽다.
\end{kfChoices}
\end{kfSentenceUnit}

\begin{kfSentenceUnit}{11}{셋째는 인정받고 싶은 마음이다.}
\kfSentQ{18.}{세 번째 기본적인 욕구는 무엇인가?}
\begin{kfChoices}
\kfChoice 사랑하고 싶은 마음
\kfChoice 자유롭고 싶은 마음
\kfChoice 인정받고 싶은 마음
\kfChoice 즐겁게 살고 싶은 마음
\end{kfChoices}
\end{kfSentenceUnit}

\begin{kfSentenceUnit}{12}{다른 사람에게 잘한다는 말을 듣고 싶고, 스스로 뿌듯함을 느끼고 싶은 마음이다.}
\kfSentQ{19.}{'인정받고 싶은 마음'에 대한 설명으로 알맞은 것은?}
\begin{kfChoices}
\kfChoice 안전하고 건강하게 오래 살고 싶어 한다.
\kfChoice 다른 사람과 어울려 사이좋게 지내고 싶어 한다.
\kfChoice 다른 사람에게 잘한다는 말을 듣고 싶어 한다.
\kfChoice 간섭받지 않고 자유롭게 행동하고 싶어 한다.
\end{kfChoices}
\end{kfSentenceUnit}

\begin{kfSentenceUnit}{13}{이런 사람은 공부나 일에 열심히 하고, 자기 생각을 잘 말하는 편이다.}
\kfSentQ{20.}{'이런 사람'이 가리키는 것은 누구인가?}
\begin{kfChoices}
\kfChoice 사랑하고 싶은 마음이 강한 사람
\kfChoice 자유롭고 싶은 마음이 강한 사람
\kfChoice 인정받고 싶은 마음이 강한 사람
\kfChoice 즐겁게 살고 싶은 마음이 강한 사람
\end{kfChoices}
\kfSentQ{21.}{'인정받고 싶은 마음'이 강한 사람의 특징으로 알맞은 것은?}
\begin{kfChoices}
\kfChoice 규칙을 잘 지키고, 필요한 것을 아끼며 모으는 것을 좋아한다.
\kfChoice 공부나 일에 열심히 하고, 자기 생각을 잘 말하는 편이다.
\kfChoice 다른 사람을 잘 도와주고, 함께 지내는 것을 좋아한다.
\kfChoice 취미 활동을 즐기고, 항상 밝고 호기심이 많다.
\end{kfChoices}
\end{kfSentenceUnit}

\begin{kfSentenceUnit}{14}{넷째는 자유롭고 싶은 마음이다.}
\kfSentQ{22.}{네 번째 기본적인 욕구는 무엇인가?}
\begin{kfChoices}
\kfChoice 자유롭고 싶은 마음
\kfChoice 사랑하고 싶은 마음
\kfChoice 인정받고 싶은 마음
\kfChoice 즐겁게 살고 싶은 마음
\end{kfChoices}
\end{kfSentenceUnit}

\begin{kfSentenceUnit}{15}{누구에게도 간섭받지 않고, 스스로 선택하고 싶은 마음이다.}
\kfSentQ{23.}{'자유롭고 싶은 마음'에 대한 설명으로 알맞은 것은?}
\begin{kfChoices}
\kfChoice 건강을 지키고 안전하게 살고 싶어 하는 마음
\kfChoice 다른 사람과 함께 지내며 사랑을 주고받고 싶어 하는 마음
\kfChoice 남에게 인정받고 스스로 뿌듯함을 느끼고 싶은 마음
\kfChoice 간섭받지 않고 스스로 선택하고 싶은 마음
\end{kfChoices}
\end{kfSentenceUnit}

\begin{kfSentenceUnit}{16}{이 마음이 강한 사람은 혼자 있는 걸 좋아하고, 다른 사람과 너무 가까워지는 것을 불편해할 수도 있다.}
\kfSentQ{24.}{'이 마음'이 가리키는 것은 무엇인가?}
\begin{kfChoices}
\kfChoice 사랑하고 싶은 마음
\kfChoice 자유롭고 싶은 마음
\kfChoice 살고 싶은 마음
\kfChoice 인정받고 싶은 마음
\end{kfChoices}
\kfSentQ{25.}{'자유롭고 싶은 마음'이 강한 사람의 특징으로 알맞은 것은?}
\begin{kfChoices}
\kfChoice 규칙을 잘 지키고, 필요한 것을 아끼며 모으는 것을 좋아한다.
\kfChoice 다른 사람을 잘 도와주고, 함께 지내는 것을 좋아한다.
\kfChoice 혼자 있는 것을 좋아하고, 간섭받는 것을 싫어한다.
\kfChoice 공부나 일에 열심히 하고, 자기 생각을 잘 말하는 편이다.
\end{kfChoices}
\end{kfSentenceUnit}

\begin{kfSentenceUnit}{17}{다섯째는 즐겁게 살고 싶은 마음이다.}
\kfSentQ{26.}{다섯 번째 기본적인 욕구는 무엇인가?}
\begin{kfChoices}
\kfChoice 즐겁게 살고 싶은 마음
\kfChoice 사랑하고 싶은 마음
\kfChoice 인정받고 싶은 마음
\kfChoice 자유롭고 싶은 마음
\end{kfChoices}
\end{kfSentenceUnit}

\begin{kfSentenceUnit}{18}{새로운 것을 배우거나, 놀면서 재미를 느끼고 싶어 하는 마음이다.}
\kfSentQ{27.}{'즐겁게 살고 싶은 마음'에 대한 설명으로 알맞은 것은?}
\begin{kfChoices}
\kfChoice 건강을 지키고 안전하게 살고 싶어 하는 마음
\kfChoice 남에게 인정받고 스스로 뿌듯함을 느끼고 싶은 마음
\kfChoice 간섭받지 않고 스스로 선택하고 싶은 마음
\kfChoice 새로운 것을 배우거나, 놀면서 재미를 느끼고 싶어 하는 마음
\end{kfChoices}
\end{kfSentenceUnit}

\begin{kfSentenceUnit}{19}{이런 사람은 취미 활동을 즐기고, 항상 밝고 호기심이 많다.}
\kfSentQ{28.}{'이런 사람'이 가리키는 것은 누구인가?}
\begin{kfChoices}
\kfChoice 사랑하고 싶은 마음이 강한 사람
\kfChoice 즐겁게 살고 싶은 마음이 강한 사람
\kfChoice 자유롭고 싶은 마음이 강한 사람
\kfChoice 인정받고 싶은 마음이 강한 사람
\end{kfChoices}
\kfSentQ{29.}{'즐겁게 살고 싶은 마음'이 강한 사람의 특징으로 알맞은 것은?}
\begin{kfChoices}
\kfChoice 규칙을 잘 지키고, 필요한 것을 아끼며 모으는 것을 좋아한다.
\kfChoice 혼자 있는 것을 좋아하고, 간섭받는 것을 싫어한다.
\kfChoice 취미 활동을 즐기고, 항상 밝고 호기심이 많다.
\kfChoice 다른 사람을 잘 도와주고, 함께 지내는 것을 좋아한다.
\end{kfChoices}
\end{kfSentenceUnit}

\begin{kfSentenceUnit}{20}{사람마다 이 다섯 가지 마음의 크기는 다르다.}
\kfSentQ{30.}{'이 다섯 가지 마음'이 가리키는 것은 무엇인가?}
\begin{kfChoices}
\kfChoice 살고, 사랑하고, 인정받고, 자유롭고, 즐겁게 살고 싶은 마음
\kfChoice 보고, 듣고, 맛보고, 냄새 맡고, 느끼는 다섯 가지 감각
\kfChoice 기쁨, 슬픔, 분노, 즐거움, 두려움의 다섯 가지 감정
\kfChoice 성실함, 정직함, 용기, 지혜, 절제의 다섯 가지 덕목
\end{kfChoices}
\kfSentQ{31.}{다섯 가지 마음의 공통점으로 알맞은 것은?}
\begin{kfChoices}
\kfChoice 사람마다 크기가 같다.
\kfChoice 누구나 가지고 있다.
\kfChoice 나이가 들면 사라진다.
\kfChoice 교육을 통해 생긴다.
\end{kfChoices}
\kfSentQ{32.}{다섯 가지 마음의 다른 점으로 알맞은 것은?}
\begin{kfChoices}
\kfChoice 종류가 사람마다 다르다.
\kfChoice 한 가지만 가질 수 있다.
\kfChoice 사람마다 크기가 다르다.
\kfChoice 나이에 따라 종류가 변한다.
\end{kfChoices}
\end{kfSentenceUnit}

\begin{kfSentenceUnit}{21}{그래서 행동하는 방식도 다르고, 갈등이 생기기도 한다.}
\kfSentQ{33.}{'그래서'의 구체적인 내용으로 알맞은 것은?}
\begin{kfChoices}
\kfChoice 사람들이 서로 사랑하기 때문에
\kfChoice 다섯 가지 마음이 모두 중요해서
\kfChoice 사람마다 행동하는 방식이 같아서
\kfChoice 사람마다 이 다섯 가지 마음의 크기가 달라서
\end{kfChoices}
\kfSentQ{34.}{다섯 가지 마음의 크기가 다른 결과 나타나는 현상으로 알맞은 것은?}
\begin{kfChoices}
\kfChoice 서로 더 잘 이해하게 된다.
\kfChoice 행동하는 방식이 다르고, 갈등이 생긴다.
\kfChoice 마음의 평화를 얻게 된다.
\kfChoice 모든 욕구가 사라지게 된다.
\end{kfChoices}
\end{kfSentenceUnit}

\begin{kfSentenceUnit}{22}{예를 들어, 사랑하고 싶은 마음은 크지만 인정받고 싶은 마음은 약한 사람은, 다른 사람의 부탁을 거절하지 못해 힘들 수 있다.}
\kfSentQ{35.}{이 문장은 앞 문장의 어떤 내용에 대한 구체적인 예시인가?}
\begin{kfChoices}
\kfChoice 마음의 크기가 달라 갈등이 생기는 경우
\kfChoice 마음의 크기가 모두 같은 경우
\kfChoice 욕구를 모두 충족한 행복한 경우
\kfChoice 갈등 없이 원만하게 해결된 경우
\end{kfChoices}
\kfSentQ{36.}{예시에서 제시된 사람의 마음 상태로 알맞은 것은?}
\begin{kfChoices}
\kfChoice 인정받고 싶은 마음은 크지만 사랑하고 싶은 마음은 약하다.
\kfChoice 자유롭고 싶은 마음과 즐겁게 살고 싶은 마음이 강하다.
\kfChoice 사랑하고 싶은 마음은 크지만 인정받고 싶은 마음은 약하다.
\kfChoice 살고 싶은 마음이 가장 강하다.
\end{kfChoices}
\kfSentQ{37.}{사랑하고 싶은 마음은 크지만 인정받고 싶은 마음이 약하면 어떤 문제가 생기는가?}
\begin{kfChoices}
\kfChoice 친구들과 자주 다투고 화해하지 않는다.
\kfChoice 자기 생각만 고집하여 외톨이가 된다.
\kfChoice 규칙을 너무 엄격하게 지켜서 피곤해한다.
\kfChoice 다른 사람의 부탁을 거절하지 못해 힘들어한다.
\end{kfChoices}
\end{kfSentenceUnit}

\begin{kfSentenceUnit}{23}{이럴 때는 인정받고 싶은 마음을 키워서, 자기 생각을 솔직하게 말할 수 있도록 도와준다.}
\kfSentQ{38.}{'이럴 때'는 언제를 가리키는가?}
\begin{kfChoices}
\kfChoice 부탁을 거절하지 못해 힘들어할 때
\kfChoice 친구들과 싸우고 나서 화해하고 싶을 때
\kfChoice 자신의 욕구를 모두 충족하여 행복할 때
\kfChoice 새로운 취미를 배우고 싶을 때
\end{kfChoices}
\kfSentQ{39.}{부탁을 거절 못 하는 사람에게 현실요법은 어떤 마음을 키우도록 도와주는가?}
\begin{kfChoices}
\kfChoice 사랑하고 싶은 마음
\kfChoice 인정받고 싶은 마음
\kfChoice 자유롭고 싶은 마음
\kfChoice 즐겁게 살고 싶은 마음
\end{kfChoices}
\end{kfSentenceUnit}

\begin{kfSentenceUnit}{24}{또 어떤 사람은 자유롭고 싶은 마음과 인정받고 싶은 마음이 둘 다 강한 경우도 있다.}
\kfSentQ{40.}{이 문장이 설명하는, 마음속에서 갈등이 생길 수 있는 경우는 무엇인가?}
\begin{kfChoices}
\kfChoice 사랑하고 싶은 마음과 살고 싶은 마음이 둘 다 약한 경우
\kfChoice 즐겁게 살고 싶은 마음만 아주 강한 경우
\kfChoice 자유롭고 싶은 마음과 인정받고 싶은 마음이 둘 다 강한 경우
\kfChoice 모든 욕구가 균형을 이루고 있는 경우
\end{kfChoices}
\end{kfSentenceUnit}

\begin{kfSentenceUnit}{25}{이럴 때는 자기 생각만 고집하면 친구들과 자주 다툴 수 있다.}
\kfSentQ{41.}{자유롭고 싶은 마음과 인정받고 싶은 마음이 둘 다 강한 사람이 어떤 행동을 하면 문제가 생기는가?}
\begin{kfChoices}
\kfChoice 자기 생각만 고집하면
\kfChoice 다른 사람의 의견을 너무 많이 들으면
\kfChoice 규칙을 너무 철저히 지키면
\kfChoice 욕구를 억누르고 참으면
\end{kfChoices}
\kfSentQ{42.}{자유롭고 싶은 마음과 인정받고 싶은 마음이 둘 다 강한 사람이 자기 생각만 고집하면 겪을 수 있는 문제는 무엇인가?}
\begin{kfChoices}
\kfChoice 혼자서 외로움을 느낄 수 있다.
\kfChoice 자신의 일을 미루게 된다.
\kfChoice 친구들과 자주 다툴 수 있다.
\kfChoice 건강이 나빠질 수 있다.
\end{kfChoices}
\end{kfSentenceUnit}

\begin{kfSentenceUnit}{26}{이럴 때는 다른 사람의 생각도 기꺼이 듣고, 마음을 헤아리는 연습이 필요하다.}
\kfSentQ{43.}{'이럴 때'는 언제를 가리키는가?}
\begin{kfChoices}
\kfChoice 사랑하고 싶은 마음이 가장 클 때
\kfChoice 자유와 인정 욕구가 모두 강할 때
\kfChoice 살고 싶은 마음만 강할 때
\kfChoice 모든 욕구가 사라졌을 때
\end{kfChoices}
\kfSentQ{44.}{자기 생각만 고집하는 사람에게 필요한 연습은 무엇인가?}
\begin{kfChoices}
\kfChoice 자기 주장을 더 논리적으로 펼치는 연습
\kfChoice 혼자 있는 시간을 늘리는 연습
\kfChoice 욕구를 포기하는 연습
\kfChoice 다른 사람의 생각을 듣고 마음을 헤아리는 연습
\end{kfChoices}
\end{kfSentenceUnit}

\begin{kfSentenceUnit}{27}{현실요법에서는 자기 욕구만 채우려고 하지 않고, 다른 사람의 욕구도 방해하지 않으면서 좋은 선택을 하도록 돕는다.}
\kfSentQ{45.}{현실요법이 추구하는 목표로 알맞은 것은?}
\begin{kfChoices}
\kfChoice 자신의 욕구를 무조건 최우선으로 채우는 것
\kfChoice 다른 사람을 위해 자신의 모든 욕구를 희생하는 것
\kfChoice 모든 욕구를 없애고 마음을 비우는 것
\kfChoice 타인의 욕구를 존중하며 좋은 선택을 돕는 것
\end{kfChoices}
\kfSentQ{46.}{현실요법에서 피해야 할 것으로 알맞은 것은?}
\begin{kfChoices}
\kfChoice 자기 욕구만 채우려고 하는 것
\kfChoice 스스로 선택하고 행동하는 것
\kfChoice 다른 사람의 이야기를 듣는 것
\kfChoice 자신의 꿈을 위해 노력하는 것
\end{kfChoices}
\end{kfSentenceUnit}

\begin{kfSentenceUnit}{28}{사람은 남이 시키는 대로만 사는 존재가 아니라, 스스로 선택하고 행동할 수 있는 존재라는 것이 현실요법의 가장 중요한 생각이다.}
\kfSentQ{47.}{현실요법이 생각하는 인간의 가장 중요한 특징은 무엇인가?}
\begin{kfChoices}
\kfChoice 환경에 의해 지배받는 존재
\kfChoice 본능에 따라 움직이는 존재
\kfChoice 스스로 선택하고 행동할 수 있는 존재
\kfChoice 남이 시키는 대로 따르는 존재
\end{kfChoices}
\kfSentQ{48.}{이 문장에서 대조되는 두 가지 인간관으로 알맞은 것은?}
\begin{kfChoices}
\kfChoice 착한 사람과 나쁜 사람
\kfChoice 성공한 사람과 실패한 사람
\kfChoice 친구가 많은 사람과 적은 사람
\kfChoice 시키는 대로 사는 존재와 스스로 선택하는 존재
\end{kfChoices}
\end{kfSentenceUnit}

\begin{kfSentenceUnit}{29}{요즘에는 이 상담 방법이 많은 사람들에게 관심을 받고, 여러 심리 상담에서 넓게 사용되고 있다.}
\kfSentQ{49.}{'이 상담 방법'이 가리키는 것은 무엇인가?}
\begin{kfChoices}
\kfChoice 현실요법
\kfChoice 정신분석
\kfChoice 행동수정
\kfChoice 최면요법
\end{kfChoices}
\kfSentQ{50.}{현실요법은 요즘 어떻게 평가받고 있는가?}
\begin{kfChoices}
\kfChoice 비과학적이라며 비판받고 있다.
\kfChoice 너무 어려워서 전문가들만 사용한다.
\kfChoice 효과가 없어서 점차 사라지고 있다.
\kfChoice 많은 사람에게 관심을 받고 넓게 사용되고 있다.
\end{kfChoices}
\end{kfSentenceUnit}

\kfSecHeader{kfAreaNonlit}{활동}{구조도}
\noindent\textit{\small 글의 흐름을 살펴보고, 구조도의 빈칸을 알맞게 채워 보자.}\par\medskip
\par\medskip\needspace{6\baselineskip}
\begin{tcolorbox}[enhanced, breakable=false,
  colback=kfPaper, colframe=kfRule, boxrule=0.4pt, arc=2pt,
  left=10pt, right=10pt, top=8pt, bottom=8pt,
  title={\footnotesize\bfseries\color{kfMute} 글의 구조도},
  coltitle=kfMute, colbacktitle=kfPrimaryLight!40,
  attach boxed title to top left={yshift=-1.5mm, xshift=6pt},
  boxed title style={colframe=kfRule, boxrule=0.3pt, arc=1pt}]
\setstretch{1.4}\footnotesize
지침: 글의 전체적인 논리 흐름을 파악하며 빈칸에 알맞은 내용을 채우시오. \\

\noindent \textbullet\ 1문단: 현실요법의 기본 생각 \\
\noindent \quad\textbullet\ 인간의 모든 행동 = ( ① )을 이루기 위한 ( ② ) \\
\noindent \quad\textbullet\ 문제 해결 방법 = 마음속 욕구들의 ( ③ ) 맞추기 \\
\noindent \textbullet\ 2문단: 다섯 가지 기본 욕구 \\
\noindent \quad\textbullet\ 욕구 1: ( ④ ) 싶은 마음 \\
\noindent \quad\textbullet\ 욕구 2: ( ⑤ ) 싶은 마음 \\
\noindent \quad\textbullet\ 욕구 3: ( ⑥ )받고 싶은 마음 \\
\noindent \quad\textbullet\ 욕구 4: ( ⑦ ) 싶은 마음 \\
\noindent \quad\textbullet\ 욕구 5: ( ⑧ ) 살고 싶은 마음 \\
\noindent \textbullet\ 3문단: 갈등의 원인과 해결 \\
\noindent \quad\textbullet\ 원인: 사람마다 다섯 가지 마음의 ( ⑨ )가 다르기 때문 \\
\noindent \quad\textbullet\ 사례 및 해결책 제시 \\
\noindent \textbullet\ 4문단: 현실요법의 목표와 현재 \\
\noindent \quad\textbullet\ 핵심 목표: 다른 사람을 존중하며 ( ⑩ ) 좋은 선택을 하도록 돕는 것 \\
\noindent \quad\textbullet\ 현재: 많은 ( ⑪ )을 받으며 넓게 사용되고 있음 \\
\end{tcolorbox}\par\medskip
\kfSecHeader{kfAreaNonlit}{문제}{}

\begin{kfQBlock}{1}{}
\kfQStem{윗글의 중심 내용으로 적절하지 \uline{않은} 것은?}
\begin{kfChoices}
\kfChoice 현실요법이 제시하는 다섯 가지 기본 욕구
\kfChoice 인간의 행동을 결정하는 뇌 과학과 신경 회로의 작동 원리
\kfChoice 욕구의 균형을 통해 더 나은 선택과 관계를 만드는 길
\kfChoice 현실요법이 인간의 행동과 정서를 설명하는 관점
\kfChoice 기본 욕구를 이해함으로써 자신의 행동을 점검하는 방법
\end{kfChoices}
\end{kfQBlock}
\begin{kfQBlock}{2}{}
\kfQStem{윗글의 내용과 일치하지 \uline{않는} 것은?}
\begin{kfChoices}
\kfChoice 현실요법은 사람의 모든 행동을 스스로의 선택으로 본다.
\kfChoice 모든 사람의 다섯 가지 마음의 크기는 똑같다.
\kfChoice '살고 싶은 마음'이 강한 사람은 규칙을 잘 지키는 경향이 있다.
\kfChoice 현실요법은 다른 사람의 욕구도 존중하는 선택을 하도록 돕는다.
\kfChoice 현실요법은 요즘 여러 심리 상담에서 사용되고 있다.
\end{kfChoices}
\end{kfQBlock}
\begin{kfQBlock}{3}{}
\kfQStem{<보기>의 민수가 가장 키워야 할 마음으로 적절한 것은?}
\begin{kfEvidence}민수는 친구들과 어울려 노는 것을 가장 좋아한다. 하지만 친구들이 무리한 부탁을 할 때도 거절하지 못하고 다 들어주다가 결국 자기 할 일을 못 해서 속상할 때가 많다.\end{kfEvidence}
\begin{kfChoices}
\kfChoice 살고 싶은 마음
\kfChoice 사랑하고 싶은 마음
\kfChoice 인정받고 싶은 마음
\kfChoice 자유롭고 싶은 마음
\kfChoice 즐겁게 살고 싶은 마음
\end{kfChoices}
\end{kfQBlock}
\begin{kfQBlock}{4}{}
\kfQStem{다음 중 '즐겁게 살고 싶은 마음'이 강한 사람의 특징으로 가장 적절한 것은?}
\begin{kfChoices}
\kfChoice 취미 활동을 즐기고 항상 밝고 호기심이 많다.
\kfChoice 다른 사람을 잘 도와주지만 상처받기 쉽다.
\kfChoice 공부나 일에 열심히 하고 자기 생각을 잘 말한다.
\kfChoice 혼자 있는 것을 좋아하고 간섭받기를 싫어한다.
\kfChoice 규칙을 잘 지키고 필요한 것을 아껴 모은다.
\end{kfChoices}
\end{kfQBlock}
\begin{kfQBlock}{5}{}
\kfQStem{현실요법이 생각하는 '인간'에 대한 설명으로 가장 중요한 것은?}
\begin{kfChoices}
\kfChoice 다른 사람의 도움 없이는 살아갈 수 없는 존재
\kfChoice 다섯 가지 마음의 크기가 정해져 있어 변하지 않는 존재
\kfChoice 남이 시키는 대로 행동하는 것이 가장 편한 존재
\kfChoice 스스로 선택하고 행동할 수 있는 힘을 가진 존재
\kfChoice 다른 사람의 욕구를 먼저 이루어주어야 하는 존재
\end{kfChoices}
\end{kfQBlock}
\begin{kfQBlock}{6}{}
\kfQStem{'자유롭고 싶은 마음'과 '인정받고 싶은 마음'이 둘 다 강한 사람이 겪을 수 있는 문제점으로 가장 적절한 것은?}
\begin{kfChoices}
\kfChoice 다른 사람의 부탁을 거절하지 못해 힘들어한다.
\kfChoice 혼자 있는 것을 너무 좋아해서 친구가 없다.
\kfChoice 자기 생각만 고집하여 친구들과 자주 다툰다.
\kfChoice 규칙을 너무 중요하게 생각해서 답답해 보인다.
\kfChoice 새로운 것을 배우는 데에만 시간을 너무 많이 쓴다.
\end{kfChoices}
\end{kfQBlock}
\begin{kfQBlock}{7}{}
\kfQStem{윗글에 나타난 현실요법의 상담 목표를 가장 정확히 진술한 것은?}
\begin{kfChoices}
\kfChoice 다섯 가지 욕구 중 한 가지를 선택해 그 욕구만 집중적으로 키우도록 돕는다.
\kfChoice 자기 욕구만 채우려 하지 않고 타인의 욕구도 방해하지 않으면서 좋은 선택을 하도록 돕는다.
\kfChoice 문제가 되는 욕구를 마음에서 완전히 없애 갈등의 원인을 제거하도록 돕는다.
\kfChoice 상담자의 지시에 따라 정해진 행동을 반복함으로써 욕구를 통제하도록 돕는다.
\kfChoice 다섯 가지 욕구 중 가장 약한 욕구를 진단해 그것을 키우라고 알려 주는 데 그친다.
\end{kfChoices}
\end{kfQBlock}
\begin{kfQBlock}{1}{}
\kfQStem{사람의 모든 행동이 마음속 욕구를 이루기 위한 선택이라고 보는 상담 방법은 무엇인가?}
\kfAnswerLines{1}
\end{kfQBlock}
\begin{kfQBlock}{2}{}
\kfQStem{건강을 지키고 안전하게 살고 싶어 하는 기본적인 마음은 무엇인가?}
\kfAnswerLines{1}
\end{kfQBlock}
\begin{kfQBlock}{3}{}
\kfQStem{다른 사람과 함께 지내고 사랑을 주고받고 싶어 하는 마음은 무엇인가?}
\kfAnswerLines{1}
\end{kfQBlock}
\begin{kfQBlock}{4}{}
\kfQStem{다른 사람에게 잘한다는 말을 듣고 스스로 뿌듯함을 느끼고 싶은 마음은 무엇인가?}
\kfAnswerLines{1}
\end{kfQBlock}
\begin{kfQBlock}{5}{}
\kfQStem{누구에게도 간섭받지 않고 스스로 선택하고 싶은 마음은 무엇인가?}
\kfAnswerLines{1}
\end{kfQBlock}
\begin{kfQBlock}{1}{}
\kfQStem{현실요법에서 갈등이 생기는 이유와 해결 방법을 다음 조건에 맞게 써보세요.}
\kfAnswerNote{답안}{4}
\end{kfQBlock}



\clearpage

\kfChapterCover{2}{러셀1 (챕터2)}{Chapter 2}{이번 챕터의 학습 내용을 시작합니다.}

\kfStartGrammar{문법 영역 — 문장의 뼈대를 세우는 말 (명사, 대명사, 수사)}


\kfSecHeader{kfAreaGrammar}{개념}{문장의 뼈대를 세우는 말 (명사, 대명사, 수사)}

하나의 문장이 만들어지려면 ‘누가’ 또는 ‘무엇이’에 해당하는 주인공이 반드시 있어야 한다. 이러한 주인공 역할을 하는 대표적인 말에는 세 가지가 있는데, 바로 명사, 대명사, 수사다.

명사는 세상 모든 것의 ‘이름’을 나타내는 말이다. 나의 이름도 명사고, 우리 엄마 이름도 명사다. 우리가 매일 보는 ‘학교, 컴퓨터, 의자’ 같은 물건의 이름도 명사고, 눈에는 안 보이지만 머릿속으로 생각할 수 있는 ‘사랑, 행복, 우정’ 같은 이름도 전부 명사다. 한마디로 이름표를 붙일 수 있는 것은 모두 명사라고 할 수 있다.

대명사는 명사 대신 쓰이는 말로, 앞에 나온 이름을 반복하는 것을 피하고 싶을 때 나타나는 말이다. 예를 들어 ‘다인이는 그림을 잘 그린다. 그리고 다인이는 마음씨도 착하다.’처럼 같은 이름을 반복하면 어색하다. 이때 두 번째 ‘다인이’를 ‘그녀’로 바꿔 ‘다인이는 그림을 잘 그린다. 그리고 그녀는 마음씨도 착하다.’처럼 자연스럽게 만들어 준다. 이처럼 ‘나, 너, 우리, 그, 그녀’나 ‘이것, 저것, 거기’ 같은 말들이 바로 대명사다.

수사는 이름 그대로 ‘숫자’와 관련이 깊다. 수사는 수량이나 순서를 나타내는 말이다. '하나', '둘', '셋'처럼 개수를 세는 말과 '첫째', '둘째', '셋째'처럼 순서를 나타내는 말이 모두 수사에 속한다.
\par\medskip\begin{itemize}[leftmargin=18pt, itemsep=2pt, topsep=2pt, label=\textbullet]
\item 명사 — 구체적(학교·의자) / 추상적(사랑·우정)
\item 대명사 — 인칭(나·너·그녀) / 지시(이것·저것·거기)
\item 수사 — 양수사(하나·둘·셋) / 서수사(첫째·둘째·셋째)
\end{itemize}

\kfSecHeader{kfAreaGrammar}{활동}{내용 확인}
\noindent\textit{\small 아래 표의 빈칸에 알맞은 말을 채워 문법 내용을 확인해 보자.}\par\medskip
* 문장의 뼈대를 세우는 말 *

문장의 주인공에 어떤 것들이 있는지 아래 표를 채우며 알아보자. 

\par\medskip\needspace{4\baselineskip}
\noindent\begin{adjustbox}{max width=\linewidth, center}
\renewcommand{\arraystretch}{1.45}
\arrayrulecolor{kfRule}
\begin{tabular}{|>{\raggedright\arraybackslash}p{\dimexpr 0.2338\linewidth-12pt\relax}|>{\raggedright\arraybackslash}p{\dimexpr 0.5195\linewidth-12pt\relax}|>{\raggedright\arraybackslash}p{\dimexpr 0.2468\linewidth-12pt\relax}|}
\hline
\rowcolor{kfPrimaryLight!50}\textbf{\color{kfPrimary} 주인공의 종류} & \textbf{\color{kfPrimary} 설 명} & \textbf{\color{kfPrimary} 예 시} \\\hline
[ ① \_\_\_\_\_\_ ] & 세상 모든 것의 이름을 나타내는 말 (이름표를 붙일 수 있는 모든 것!) & 학교, 컴퓨터, 사랑, 행복 \\\hline
[ ② \_\_\_\_\_\_ ] & 앞에 나온 이름을 대신 부르는 말 (같은 말 반복을 피하게 해줌!) & 나, 너, 우리, 그, 그녀, 이것 \\\hline
[ ③ \_\_\_\_\_\_ ] & 수량(개수)이나 순서를 나타내는 말 & 하나, 둘, 셋, 첫째, 둘째 \\\hline
\end{tabular}
\arrayrulecolor{black}
\end{adjustbox}\par\medskip

\kfSecHeader{kfAreaGrammar}{활동}{분석 훈련}
\noindent\textit{\small  명사, 대명사, 수사 찾기 훈련! 문장에서 명사(이름), 대명사(이름 대신), 수사(숫자)에 해당하는 단어를 모두 찾아 쓰시오.}\par\medskip
\begin{enumerate}[leftmargin=22pt, itemsep=6pt, topsep=4pt, label=\textbf{\arabic*.}]
\item 하늘에 무지개가 예쁘게 떠 있다.
\item 우리는 공원에서 그를 만났다.
\item 가게에 가서 사과 하나를 샀다.
\item 우리는 저기에서 배를 보았다.
\item 나의 꿈은 첫째도 희망, 둘째도 희망이다.
\item 여기에 있는 책은 다 재미있다.
\item 철수는 운동장에서 친구와 축구를 했다.
\item 이것은 도하의 연필이다.
\item 하나보다 둘이 더 낫다는 말이 있다.
\item 이분은 우리 할머니시고, 저분은 우리 할아버지시다.
\end{enumerate}

\kfSecHeader{kfAreaGrammar}{문제}{}

\begin{kfQBlock}{1}{}
\kfQStem{다음 중 품사가 \uline{다른} 하나는?}
\begin{kfChoices}
\kfChoice 행복
\kfChoice 우정
\kfChoice 우리
\kfChoice 학교
\kfChoice 연필
\end{kfChoices}
\end{kfQBlock}
\begin{kfQBlock}{2}{}
\kfQStem{‘수사’에 대한 설명으로 옳은 것은?}
\begin{kfChoices}
\kfChoice 사물의 이름을 나타낸다.
\kfChoice 명사를 대신하여 사용된다.
\kfChoice 수량이나 순서를 나타낸다.
\kfChoice 다른 말을 꾸며주는 역할을 한다.
\kfChoice 놀람이나 느낌, 대답을 나타낸다.
\end{kfChoices}
\end{kfQBlock}
\begin{kfQBlock}{3}{}
\kfQStem{다음 밑줄 친 단어의 품사가 명사가 \uline{아닌} 것은?}
\begin{kfChoices}
\kfChoice 나는 \uline{사과}를 좋아한다.
\kfChoice \uline{이것}은 내 책이다.
\kfChoice \uline{사랑}은 위대하다.
\kfChoice \uline{철수}는 내 친구다.
\kfChoice \uline{컴퓨터}가 고장 났다.
\end{kfChoices}
\end{kfQBlock}
\begin{kfQBlock}{4}{}
\kfQStem{다음 중 수량을 나타내는 수사가 \uline{아닌} 것은?}
\begin{kfChoices}
\kfChoice 하나
\kfChoice 일곱
\kfChoice 아홉
\kfChoice 넷째
\kfChoice 열
\end{kfChoices}
\end{kfQBlock}
\begin{kfQBlock}{5}{}
\kfQStem{다음 중 본문에서 ‘눈으로 볼 수 있는 이름(구체명사)’으로 분류한 것에 해당하지 \uline{않는} 것은?}
\begin{kfChoices}
\kfChoice 책상
\kfChoice 우정
\kfChoice 연필
\kfChoice 바다
\kfChoice 나무
\end{kfChoices}
\end{kfQBlock}
\begin{kfQBlock}{6}{}
\kfQStem{다음 문장에 사용된 대명사는 모두 몇 개인가?}
\begin{kfChoices}
\kfChoice 1개
\kfChoice 2개
\kfChoice 3개
\kfChoice 4개
\kfChoice 5개
\end{kfChoices}
\end{kfQBlock}
\begin{kfQBlock}{7}{}
\kfQStem{다음 중 품사의 연결이 바르지 \uline{않은} 것은?}
\begin{kfChoices}
\kfChoice 의자 - 명사
\kfChoice 그녀 - 대명사
\kfChoice 다섯 - 수사
\kfChoice 행복 - 대명사
\kfChoice 저것 - 대명사
\end{kfChoices}
\end{kfQBlock}
\begin{kfQBlock}{8}{}
\kfQStem{다음 중 장소를 가리키는 대명사는?}
\begin{kfChoices}
\kfChoice 무엇
\kfChoice 누구
\kfChoice 거기
\kfChoice 언제
\kfChoice 저희
\end{kfChoices}
\end{kfQBlock}
\begin{kfQBlock}{9}{}
\kfQStem{다음 문장에서 수사는 총 몇 번 사용되었는가?}
\begin{kfEvidence}우리는 사과 하나와 배 둘을 샀고, 형은 달리기에서 첫째로 들어왔다.\end{kfEvidence}
\begin{kfChoices}
\kfChoice 1개
\kfChoice 2개
\kfChoice 3개
\kfChoice 4개
\kfChoice 5개
\end{kfChoices}
\end{kfQBlock}
\begin{kfQBlock}{10}{}
\kfQStem{다음 <보기>에 대한 설명으로 옳은 것은?}
\begin{kfEvidence}가. 그는 책 하나를 샀다. 나. 그것은 그녀의 가방이다.\end{kfEvidence}
\begin{kfChoices}
\kfChoice '가'의 ‘그’와 ‘나'의 ‘그것’은 모두 명사다.
\kfChoice '가'의 ‘책’과 '나'의 ‘가방’은 모두 대명사다.
\kfChoice '가'의 ‘하나’는 명사이고, '나'의 ‘그녀’는 수사다.
\kfChoice '가'의 ‘그’와 '나'의 ‘그것’은 모두 대명사다.
\kfChoice '가'의 ‘하나’는 대명사이고, '나'의 ‘가방’은 수사다.
\end{kfChoices}
\end{kfQBlock}
\begin{kfQBlock}{1}{}
\kfQStem{‘의자’, ‘하늘’처럼 구체적인 사물의 이름을 나타내는 품사는 무엇인가?}
\kfAnswerLines{1}
\end{kfQBlock}
\begin{kfQBlock}{2}{}
\kfQStem{명사를 대신하여 쓰이는 품사는 무엇인가?}
\kfAnswerLines{1}
\end{kfQBlock}
\begin{kfQBlock}{3}{}
\kfQStem{수량이나 순서를 나타내는 품사는 무엇인가?}
\kfAnswerLines{1}
\end{kfQBlock}
\begin{kfQBlock}{4}{}
\kfQStem{눈에 보이지 않는 ‘희망’, ‘슬픔’ 등의 단어는 어떤 품사에 속하는가?}
\kfAnswerLines{1}
\end{kfQBlock}
\begin{kfQBlock}{5}{}
\kfQStem{‘다인이는 착하다. 그녀는 인기가 많다.’에서 ‘그녀’의 품사는 무엇인가?}
\kfAnswerLines{1}
\end{kfQBlock}
\begin{kfQBlock}{6}{}
\kfQStem{‘사과 하나와 배 둘을 샀다.’에서 ‘하나’와 ‘둘’의 품사는 무엇인가?}
\kfAnswerLines{1}
\end{kfQBlock}
\begin{kfQBlock}{7}{}
\kfQStem{‘첫째도 건강, 둘째도 건강’에서 ‘첫째’, ‘둘째’의 품사는 무엇인가?}
\kfAnswerLines{1}
\end{kfQBlock}
\begin{kfQBlock}{8}{}
\kfQStem{‘너와 나’에서 ‘너’와 ‘나’의 품사는 무엇인가?}
\kfAnswerLines{1}
\end{kfQBlock}
\begin{kfQBlock}{9}{}
\kfQStem{‘컴퓨터’의 품사는 무엇인가?}
\kfAnswerLines{1}
\end{kfQBlock}
\begin{kfQBlock}{10}{}
\kfQStem{‘저기가 우리 집이야.’에서 ‘저기’의 품사는 무엇인가?}
\kfAnswerLines{1}
\end{kfQBlock}
\begin{kfQBlock}{1}{}
\kfQStem{다음 문장에서 명사, 대명사, 수사를 각각 모두 찾아 쓰시오.}
\begin{kfEvidence}우리 셋이 힘을 합쳐 저 문제를 풀었다.\end{kfEvidence}
\kfAnswerNote{답안}{4}
\end{kfQBlock}


\kfStartConcept{개념 영역 — 이야기가 펼쳐지는 시간과 공간}


\kfSecHeader{kfAreaConcept}{개념}{이야기가 펼쳐지는 시간과 공간}

소설의 배경은 등장인물들이 활동하고 사건이 일어나는 시간과 공간을 말한다. 배경은 단순히 이야기가 벌어지는 무대 역할만 하는 것이 아니라, 인물의 행동과 감정에 영향을 주고 작품의 분위기를 만드는 중요한 역할을 한다.

배경은 크게 시간적 배경과 공간적 배경으로 나누어진다. 시간적 배경은 이야기가 일어나는 때를 나타내며, 계절이나 하루 중 특정 시간, 역사적 시대 등을 포함한다. 공간적 배경은 이야기가 펼쳐지는 장소를 말하며, 나라나 지역 같은 넓은 공간부터 집이나 방 같은 좁은 공간까지 모두 해당한다.

이효석의 「메밀꽃 필 무렵」을 예로 들어 살펴본다. 이 작품의 시간적 배경은 달빛이 쏟아지는 여름밤이고, 공간적 배경은 메밀꽃이 피어있는 산길이다. 이런 배경은 허생원이 과거의 사랑 이야기를 꺼내기에 알맞은 낭만적이고 서정적인 분위기를 만들어준다. 만약 추운 겨울밤이었다면 허생원이 그런 이야기를 할 마음이 들었을까? 또한 좁은 길이라는 공간적 배경은 세 사람이 일렬로 서게 만들어, 허생원의 이야기가 동이에게 제대로 들리지 않게 하는 자연스러운 장치가 된다.

이처럼 배경은 인물의 행동에 필연성을 부여하고, 작품의 분위기를 조성하며, 때로는 앞으로 일어날 일을 암시하는 소설의 중요한 구성 요소이다.
\par\medskip\begin{itemize}[leftmargin=18pt, itemsep=2pt, topsep=2pt, label=\textbullet]
\item 시간 — 「메밀꽃 필 무렵」의 '달빛 쏟아지는 여름밤'(낭만적 분위기)
\item 공간 — 같은 작품의 '메밀꽃 핀 좁은 산길'(세 사람의 동선 결정)
\item 분위기 조성 — '비 오는 일제 강점기 평양역' → 식민지 침체감
\end{itemize}

\kfSecHeader{kfAreaConcept}{문제}{}

\begin{kfQBlock}{1}{}
\kfQStem{소설의 ‘배경’에 대한 설명으로 적절하지 \uline{않은} 것은 무엇인가?}
\begin{kfChoices}
\kfChoice 작품의 주제·인물·사건 가운데 ‘인물’만을 가리키는 요소이다.
\kfChoice 시간적·공간적 배경뿐만 아니라 사회적인 배경까지 포함한다.
\kfChoice 작품의 분위기를 형성하고 인물의 심리에 영향을 준다.
\kfChoice 사건의 의미를 구체화하고 주제를 효과적으로 드러낸다.
\kfChoice 인물들이 활동하는 특정한 시간과 공간을 가리킨다.
\end{kfChoices}
\end{kfQBlock}
\begin{kfQBlock}{2}{}
\kfQStem{박지원의 「허생전」이 ‘묵적골’·‘남산’ 같은 실제 지명을 사용한 이유로 적절하지 \uline{않은} 것은 무엇인가?}
\begin{kfChoices}
\kfChoice 이야기에 현실감을 부여해 진짜처럼 느끼게 하기 위해서이다.
\kfChoice 독자들에게 서울의 지리를 자세하게 가르쳐 주기 위해서이다.
\kfChoice 인물의 가난한 처지를 구체적인 공간으로 강조하기 위해서이다.
\kfChoice 사회 비판이라는 주제를 보다 설득력 있게 전달하기 위해서이다.
\kfChoice 당시 조선 사회의 모습을 사실적으로 보여 주기 위해서이다.
\end{kfChoices}
\end{kfQBlock}
\begin{kfQBlock}{3}{}
\kfQStem{현진건의 「운수 좋은 날」에서 ‘비’가 내리는 배경이 주는 효과로 적절하지 \uline{않은} 것은 무엇인가?}
\begin{kfChoices}
\kfChoice 주인공의 마음을 상쾌하고 즐겁게 만들어 희망을 준다.
\kfChoice 사건 전체를 음울한 정조로 감싸는 분위기를 형성한다.
\kfChoice 김 첨지의 불안한 심리를 시각적으로 환기한다.
\kfChoice ‘운수 좋은 날’이라는 제목과 대비되어 반어적 효과를 강화한다.
\kfChoice 어둡고 불길한 분위기를 만들어 비극적 결말을 암시한다.
\end{kfChoices}
\end{kfQBlock}
\begin{kfQBlock}{4}{}
\kfQStem{<보기>의 ‘소나기’ 배경이 이야기 전개에 한 역할로 적절하지 \uline{않은} 것은 무엇인가?}
\begin{kfEvidence}징검다리를 건너는데 소나기가 쏟아졌다. ... 비를 피할 곳을 찾다가 마침 보이는 수수단 속으로 두 아이는 뛰어들었다.\end{kfEvidence}
\begin{kfChoices}
\kfChoice 두 사람이 다투고 서로를 원망하게 만드는 갈등의 원인이 된다.
\kfChoice 둘만의 시간을 만들어 짧지만 깊은 정을 나누게 한다.
\kfChoice 사건의 분위기와 정서를 응축적으로 보여 주는 장치이다.
\kfChoice 작품의 제목이자 핵심 사건으로 작용한다.
\kfChoice 서먹한 두 사람을 한 공간에 머물게 해 가까워지는 계기를 만든다.
\end{kfChoices}
\end{kfQBlock}
\begin{kfQBlock}{5}{}
\kfQStem{<보기>의 ‘장마’가 상징하는 의미로 적절하지 \uline{않은} 것은 무엇인가?}
\begin{kfEvidence}이 밤중에 억수로 내리는 비를 맞아 가며 마을을 활보하는 사람들은... 전쟁이 북으로 물러갔다고는 하지만 아직도 빨치산들이 읍내 경찰서를 습격하고 불을 지를 만큼 어수선한 때였다.\end{kfEvidence}
\begin{kfChoices}
\kfChoice 6·25 전쟁이라는 우리 민족의 불행과 시련
\kfChoice 농사에 꼭 필요한 풍요로움과 여름의 활기
\kfChoice 길고 음울하게 이어지는 분단·전쟁의 고통
\kfChoice 작품의 음울한 분위기를 자아내는 시대적 배경
\kfChoice 이념 갈등으로 어두워진 한 가족·민족의 처지
\end{kfChoices}
\end{kfQBlock}
\begin{kfQBlock}{6}{}
\kfQStem{김유정의 「봄·봄」처럼 계절을 통해 작품의 분위기를 드러내는 배경은 어떤 종류에 해당하는가?}
\begin{kfChoices}
\kfChoice 시간적 배경
\kfChoice 시대적 배경
\kfChoice 공간적 배경
\kfChoice 심리적 배경
\kfChoice 사회적 배경
\end{kfChoices}
\end{kfQBlock}
\begin{kfQBlock}{7}{}
\kfQStem{허균의 「홍길동전」에서 서자를 차별하는 '조선 시대의 신분 제도'는 어떤 배경에 해당하는가?}
\begin{kfChoices}
\kfChoice 시간적 배경
\kfChoice 시대적 배경
\kfChoice 공간적 배경
\kfChoice 자연적 배경
\kfChoice 심리적 배경
\end{kfChoices}
\end{kfQBlock}
\begin{kfQBlock}{8}{}
\kfQStem{다음 중 소설 속 배경의 기능으로 알맞지 \uline{않은} 것은 무엇인가?}
\begin{kfChoices}
\kfChoice 작품 전체의 독특한 분위기를 자연스럽게 조성한다.
\kfChoice 이야기에 마치 실제처럼 느껴지는 사실성을 부여한다.
\kfChoice 일어나는 사건 전개에 그럴듯한 필연성을 부여한다.
\kfChoice 등장인물의 성격을 직접적으로 설명한다.
\kfChoice 작품 전체가 전하려는 핵심 주제를 슬쩍 암시하기도 한다.
\end{kfChoices}
\end{kfQBlock}
\begin{kfQBlock}{9}{}
\kfQStem{이효석의 「메밀꽃 필 무렵」에서 ‘달빛 비치는 메밀꽃밭’ 배경이 주는 분위기로 적절하지 \uline{않은} 것은 무엇인가?}
\begin{kfChoices}
\kfChoice 긴급하고 위험한 사건이 벌어지는 긴장된 분위기
\kfChoice 신비롭고 환상적인 자연의 분위기
\kfChoice 인물의 회상을 자연스럽게 이끌어 내는 정감 어린 분위기
\kfChoice 달빛과 메밀꽃의 흰빛이 어우러진 시각적 아름다움의 분위기
\kfChoice 서정적이고 낭만적인 분위기
\end{kfChoices}
\end{kfQBlock}
\begin{kfQBlock}{10}{}
\kfQStem{소설의 배경에 대한 설명으로 적절하지 \uline{않은} 것은 무엇인가?}
\begin{kfChoices}
\kfChoice 때로는 작품의 주제를 암시하는 상징적 의미를 지닌다.
\kfChoice 인물·사건과 무관한 단순한 배경 무대일 뿐이다.
\kfChoice 인물의 심리·성격·사건 전개에 큰 영향을 미친다.
\kfChoice 작품 전체의 분위기를 형성하는 핵심적인 요소이다.
\kfChoice 시간·공간·사회적 배경 등 여러 차원으로 나타난다.
\end{kfChoices}
\end{kfQBlock}
\begin{kfQBlock}{1}{}
\kfQStem{소설에서 등장인물들이 활동하고 사건이 일어나는 시간과 공간을 무엇이라고 하는가?}
\kfAnswerLines{1}
\end{kfQBlock}
\begin{kfQBlock}{2}{}
\kfQStem{배경은 크게 '공간적 배경'과 어떤 배경으로 나누어지는가?}
\kfAnswerLines{1}
\end{kfQBlock}
\begin{kfQBlock}{3}{}
\kfQStem{이야기가 일어나는 '때'를 나타내는 배경은 무엇인가?}
\kfAnswerLines{1}
\end{kfQBlock}
\begin{kfQBlock}{4}{}
\kfQStem{이야기가 펼쳐지는 '장소'를 나타내는 배경은 무엇인가?}
\kfAnswerLines{1}
\end{kfQBlock}
\begin{kfQBlock}{5}{}
\kfQStem{「메밀꽃 필 무렵」의 시간적 배경은 달빛이 쏟아지는 어떤 계절의 밤인가?}
\kfAnswerLines{1}
\end{kfQBlock}
\begin{kfQBlock}{6}{}
\kfQStem{「메밀꽃 필 무렵」의 공간적 배경은 무엇이 피어있는 산길인가?}
\kfAnswerLines{1}
\end{kfQBlock}
\begin{kfQBlock}{7}{}
\kfQStem{「메밀꽃 필 무렵」의 배경은 허생원이 과거 이야기를 꺼내기에 알맞은 어떤 분위기를 만들어주는가?}
\kfAnswerLines{1}
\end{kfQBlock}
\begin{kfQBlock}{1}{}
\kfQStem{소설에서 배경이 하는 중요한 역할을 서술하시오.}
\kfAnswerNote{답안}{4}
\end{kfQBlock}


\kfStartLit{문학 영역 — 작품}


\kfSecHeader{kfAreaLiterature}{지문}{작품}
\kfPassageNote{작품}{%
\textbackslash{}[앞부분 줄거리\textbackslash{}] 왼손잡이인 허 생원은 나귀와 함께 장에서 장으로 다니며 물건을 파는 장돌뱅이입니다. 봉평 장에서 장사를 마친 후 조 선달과 함께 주막집에서 술을 마시다가, 젊은 장돌뱅이 동이가 주막 아주머니와 농담치는 모습을 보고 화가 나서 동이를 꾸짖었습니다. 동이는 말대꾸하지 않고 그냥 나갔지만, 잠시 후 허 생원의 나귀가 도망갔다고 알려줍니다. 허 생원은 동이가 도와준 것이 고마웠습니다. 그들은 함께 대화 장을 향해 메밀꽃이 핀 달밤에 길을 떠났습니다.

허 생원은 여자와는 인연이 멀었습니다. 얼굴에 얽은 자국이 있는 모습을 들고 여자 앞에 나설 용기도 없었고, 여자 쪽에서 정을 보낸 적도 없었으며, 쓸쓸하고 비뚤어진 인생을 살았습니다. 주막집을 생각만 해도 어린애처럼 얼굴이 빨개지고 발밑이 떨리며 그 자리에서 깜짝 놀라버립니다.

주막집 문을 들어서서 술자리에서 정말로 동이를 만났을 때에는 어찌 된 일인지 갑자기 화가 나버렸습니다. 상 위에 붉은 얼굴을 들고 제법 여자와 농담치는 것을 보고서는 견딜 수 없었던 것입니다.

"저 녀석이 제법 놀기 좋아하는구나, 보기 싫다. 머리에 피도 안 마른 녀석이 낮부터 술 마시고 여자와 농담이야. 장돌뱅이 망신만 시키고 돌아다니는구나. 그 꼴에 우리들과 한몫 하자는 셈이지."

동이 앞을 막아서면서부터 꾸짖었습니다. '걱정도 팔자다' 하는 듯이 뻔히 쳐다보는 흥분된 눈동자와 마주칠 때, 그 순간에 따귀를 하나 때리지 않고는 참을 수 없었습니다. 동이도 화를 내며 벌떡 일어서기는 했지만, 허 생원은 조금도 얼굴색이 변하지 않고 마음먹은 대로 다 말했습니다.

"어디서 주워 온 건방진 애인지는 모르겠지만, 너에게도 아버지 어머니가 있겠지. 그 사나운 꼴 보면 마음 좋겠다. 장사란 제대로 해야 되는 거지, 여자가 다 뭐야, 나가거라, 빨리 꼴 치워."

그러나 한마디도 대답하지 않고 조용히 나가는 모습을 보니 오히려 불쌍한 생각이 들었습니다. 아직도 서먹서먹한 사이인데 너무 심하게 하지 않았을까 하고 마음이 무거워졌습니다.

(중략)

그렇다고 해도 딱 한 번의 첫 일을 잊을 수는 없었습니다. 뒤에도 앞에도 없는 단 한 번의 신기한 인연이었습니다. 봉평에 다니기 시작한 젊은 시절의 일이었지만 그것을 생각할 때만은 그도 살아온 보람을 느꼈습니다.

"달밤이었는데 어떻게 해서 그렇게 됐는지 지금 생각해도 도무지 알 수 없어."

허 생원은 오늘 밤도 또 그 이야기를 꺼내려고 합니다. 조 선달은 친구가 된 이래 귀에 못이 박히도록 들어 왔습니다. 그렇다고 싫어할 수도 없었지만 허 생원은 모르는 척하고 같은 이야기를 계속 반복하곤 했습니다.

"달밤에는 그런 이야기가 어울리거든."

조 선달 쪽을 바라보았지만, 물론 미안해서가 아니라 달빛에 감동해서였습니다. 조금 이지러졌지만 보름을 막 지난 달이 부드러운 빛을 가득 흘리고 있습니다. 대화까지는 칠십 리의 밤길로, 고개를 둘이나 넘고 개울을 하나 건너고 벌판과 산길을 걸어야 됩니다. 길은 지금 긴 산허리에 걸려 있습니다. 밤중을 지난 무렵인지 죽은 듯이 고요한 속에서 짐승 같은 달의 숨소리가 손에 잡힐 듯이 들리며, 콩 포기와 옥수수 잎새가 한층 달빛에 푸르게 젖었습니다.

산허리는 온통 메밀밭이어서 피기 시작한 꽃이 소금을 뿌린 듯이 하얀 달빛에 숨이 막힐 지경입니다. 붉은 대궁이 향기처럼 애잔하고 나귀들의 걸음도 시원합니다. 길이 좁은 까닭에 세 사람은 나귀를 타고 한 줄로 늘어섰습니다.

방울 소리가 시원스럽게 딸랑딸랑 메밀밭으로 흘러갑니다. 앞장선 허 생원의 이야기 소리는 뒤에 있는 동이에게는 똑똑히 들리지 않았지만, 그는 그대로 상쾌한 기분으로 외롭지 않았습니다.

"꼭 이런 날 밤이었네. 여관방이 무더워서 잠이 안 와서 말이야. 밤중이 되어서 혼자 일어나 개울가에 목욕하러 나갔지. 봉평은 지금이나 그때나 마찬가지로 보이는 곳마다 메밀밭이어서 개울가가 온통 하얀 꽃이야. 돌밭에 옷을 벗어도 좋을 것을, 달이 너무도 밝은 까닭에 옷을 벗으러 물방앗간으로 들어가지 않았나. 이상한 일도 많지. 거기서 갑자기 성 서방네 처녀와 마주쳤단 말이네. 봉평에서 제일가는 미인이었지."

"운명에 있었나 보지."

그냥 대답하면서 말을 아끼는 듯이 한참이나 담배를 빨 뿐이었습니다. 진한 자줏빛 연기가 밤 공기 속에 흘러서는 사라졌습니다.

"날 기다린 것은 아니었지만 그렇다고 달리 기다리는 사람이 있는 것도 아니었네. 처녀는 울고 있단 말야. 짐작은 하고 있었지만 성 서방네는 한창 어려워서 망할 판인 때였지. 한집안 일이니 딸에게 걱정이 없을 리 있겠나. 좋은 곳에만 있으면 시집도 보내겠지만 시집은 죽어도 싫다고 하고... 그런데 처녀란 울 때처럼 마음을 끄는 때가 있을까. 처음에는 놀라는 눈치였지만 걱정 있을 때는 마음이 부드러워지기도 쉬운 듯해서 이러저러 이야기가 되었네... 생각하면 무섭고도 기막힌 밤이었어."

"제천으로 도망간 건 그 다음 날이었나?"

"다음 장날에는 벌써 온 집안이 사라진 뒤였네. 장터는 소문으로 시끄러워서 기껏해야 술집에 팔려가는 것이 보통이라고 처녀의 뒷이야기가 많이 나돌았단 말이야. 제천 장터를 몇 번이나 뒤졌겠나. 하지만 처녀의 모습은 꿩 구워 먹은 자리야. 첫날밤이 마지막 밤이었지. 그때부터 봉평이 마음에 들어서 반평생을 다니게 되었네. 평생 잊을 수 있겠나."

"좋았지. 그렇게 신기한 일이란 쉽지 않어. 보통은 못난 것 얻어서 새끼 낳고, 걱정 늘고 생각만 해도 진절머리가 나지... 그런데 늙을 때까지 장돌뱅이로 지내기도 힘든 일 아닌가? 난 가을까지만 하고 이 생활과도 작별하려네. 대화쯤에 작은 가게나 하나 차리고 식구들을 부르겠어. 일 년 내내 뚜벅뚜벅 걷기란 쉬운 일이 아니야."

"옛 처녀나 만나면 같이 살까... 난 죽을 때까지 이 길 걷고 저 달 볼 테야."

산길을 벗어나니 큰길로 나왔습니다. 뒤에 있던 동이도 앞으로 나서서 나귀들은 가로 늘어섰습니다.

"총각도 젊겠다, 지금이 한창 때겠지. 주막집에서는 그만 실수를 해서 그 꼴이 되었지만 너무 생각 말게."

"아, 아닙니다. 오히려 부끄러워요. 여자란 지금 무슨 소용인가요. 자나 깨나 어머니 생각뿐인데요."

허 생원의 이야기로 풀이 죽은 탓인지 동이의 말투는 한풀 꺾인 것이었습니다.

"아버지 어머니란 말에 가슴이 터지는 것 같았지만 제겐 아버지가 없어요. 피붙이라고는 어머니 하나뿐인걸요."

"돌아가셨나?" 

"처음부터 없어요." 

"그런 일이 세상에..."

생원과 선달이 크게 웃으니, 동이는 진지하게 말할 수밖에 없었습니다.

(중략)

"어머니의 친정이 원래부터 제천이었나?"

"아니요, 시원하게 말은 안 해 주지만 봉평이라는 것만은 들었죠."

"봉평? 그래 그 아버지 성은 무엇이고?"

"알 수 있나요. 도무지 듣지를 못했으니까."

"그, 그렇겠지."

하고 중얼거리며 흐려지는 눈을 깜빡이다가 허 생원은 실수로 발을 헛디뎠습니다. 앞으로 넘어지면서 몸째 풍덩 빠져 버렸습니다.

허우적거릴수록 몸을 잡을 수 없어서 동이가 소리를 치며 가까이 왔을 때에는 벌써 꽤 흘러갔었습니다. 옷째 흠뻑 젖으니 물에 젖은 개보다도 참혹한 꼴이었습니다.

동이는 물속에서 어른을 쉽게 업을 수 있었습니다. 젖었다고는 해도 여윈 몸이라 젊은이 등에는 오히려 가벼웠습니다.

허 생원은 젖은 옷을 어느 정도 짜서 입었습니다. 이가 덜덜 떨리고 가슴이 떨리며 몹시도 추웠지만 마음은 알 수 없이 둥실둥실 가벼웠습니다.

"여관까지 빨리들 가세나. 뜰에 불을 피우고 푹 쉬어. 나귀에겐 따뜻한 물을 끓여 주고. 내일 대화 장 보고는 제천이다."

"생원도 제천으로?"

"오래간만에 가 보고 싶어. 같이 가려나, 동이?"

나귀가 걷기 시작했을 때 동이의 채찍은 왼손에 있었습니다. 오랫동안 할머니처럼 눈이 어둡던 허 생원도 이번만은 동이의 왼손잡이가 눈에 띄지 않을 수 없었습니다.

걸음도 가볍고 방울 소리가 밤 벌판에 한층 맑게 울렸습니다.

달이 꽤 기울어졌습니다.
}


\kfSecHeader{kfAreaLiterature}{활동}{문학 문장 독해}
\noindent\textit{\small 작품 속 문장을 자세히 읽고, 문장별로 주어진 물음에 답해 보자.}\par\medskip
\begin{kfSentenceUnit}{1}{"나귀가 걷기 시작했을 때 동이의 채찍은 왼손에 있었습니다. 오랫동안 할머니처럼 눈이 어둡던 허 생원도 이번만은 동이의 왼손잡이가 눈에 띄지 않을 수 없었습니다."}
\kfSentQ{1.}{다음 문장에서 서술어를 기준으로 각각의 주어는 무엇인가? (1) 서술어 '걷기 시작했을 때'의 주어는 무엇인가?}
\begin{kfChoices}
\kfChoice 나귀가
\kfChoice 동이가
\kfChoice 허 생원이
\kfChoice 조 선달이
\end{kfChoices}
\kfSentQ{1.}{(2) 서술어 '있었습니다'의 주어는 무엇인가?}
\begin{kfChoices}
\kfChoice 나귀의 채찍이
\kfChoice 동이의 채찍이
\kfChoice 허 생원의 채찍이
\kfChoice 조 선달의 채찍이
\end{kfChoices}
\kfSentQ{1.}{(3) 서술어 '어둡던'의 주어는 무엇인가?}
\begin{kfChoices}
\kfChoice 밤이
\kfChoice 길눈이
\kfChoice 눈이
\kfChoice 마음이
\end{kfChoices}
\kfSentQ{1.}{(4) 서술어 '띄지 않을 수 없었습니다'의 주어는 무엇인가?}
\begin{kfChoices}
\kfChoice 나귀의 걸음걸이가
\kfChoice 동이의 왼손잡이가
\kfChoice 허 생원의 얼굴이
\kfChoice 달빛이
\end{kfChoices}
\end{kfSentenceUnit}

\begin{kfSentenceUnit}{2}{"뒤에도 앞에도 없는 단 한 번의 신기한 인연이었습니다. 그것을 생각할 때만은 그도 살아온 보람을 느꼈습니다."}
\kfSentQ{2.}{'그것'이 가리키는 내용은 무엇인가?}
\begin{kfChoices}
\kfChoice 동이와의 만남
\kfChoice 단 한 번의 신기한 인연
\kfChoice 나귀를 잃어버린 일
\kfChoice 장사를 망친 일
\end{kfChoices}
\end{kfSentenceUnit}

\kfSecHeader{kfAreaLiterature}{활동}{해설 학습}
\noindent\textit{\small 작품 해설을 단원별로 읽고, 각 단원에 딸린 물음에 답해 보자.}\par\medskip
\par\noindent\textbf{\color{kfPrimary} 해설 단원 1}\par\smallskip
\kfPassageNote{해설 1}{%
「메밀꽃 필 무렵」에서 가장 인상적인 것은 달빛 아래 하얗게 핀 메밀꽃밭의 모습이다. '소금을 뿌린 듯이' 하얀 메밀꽃과 부드러운 달빛이 만들어내는 풍경은 정말 아름답다.

작가는 이런 자연 풍경을 통해 등장인물들의 마음을 표현하고 있다. 이 아름다운 배경은 단순히 예쁜 풍경이 아니라 허 생원이 젊은 시절 성 서방네 처녀와 만났던 그날 밤을 떠올리게 해주는 중요한 역할을 한다.

또한 '짐승 같은 달의 숨소리', '방울 소리가 시원스럽게 흘러간다'는 표현처럼 보이는 것을 들리는 것으로, 들리는 것을 느끼는 것으로 나타내는 특별한 표현 방법을 사용했다. 이런 감각적인 표현들은 바로 '공감각적 심상'이라고 부르는데, 하나의 감각을 다른 감각으로 전이시켜 표현하는 문학적 기법이다. 예를 들어 '방울 소리가 시원스럽게'에서 '시원스럽게'는 본래 촉각으로 느끼는 것인데, 이것을 청각인 '소리'와 연결시켜 표현한 것이다.

이처럼 감각을 전이시킨 표현들이 모여서 마치 한 편의 아름다운 시를 읽는 것 같은 느낌을 주며, 독자들을 꿈같은 분위기 속으로 이끌어준다.
}

\begin{kfQBlock}{1}{문제}
\kfQStem{소설에서 가장 인상적인 풍경은 달빛 아래 어떤 꽃밭의 모습인가?}
\kfAnswerLines{1}
\end{kfQBlock}

\begin{kfQBlock}{2}{문제}
\kfQStem{하얀 메밀꽃이 핀 모습을 무엇을 뿌린 것 같다고 비유했는가?}
\kfAnswerLines{1}
\end{kfQBlock}

\begin{kfQBlock}{3}{문제}
\kfQStem{아름다운 자연 풍경은 허 생원이 과거의 어떤 인물을 떠올리게 하는 중요한 역할을 하는가?}
\kfAnswerLines{1}
\end{kfQBlock}

\begin{kfQBlock}{4}{문제}
\kfQStem{‘짐승 같은 달의 숨소리’는 눈에 보이는 것을 어떤 감각으로 바꾼 표현인가?}
\kfAnswerLines{1}
\end{kfQBlock}

\begin{kfQBlock}{5}{문제}
\kfQStem{작품의 서정적인 분위기를 만드는 가장 핵심적인 배경 두 가지는 무엇인가?}
\kfAnswerLines{1}
\end{kfQBlock}

\par\noindent\textbf{\color{kfPrimary} 해설 단원 2}\par\smallskip
\kfPassageNote{해설 2}{%
이 소설은 현재의 이야기와 과거의 이야기, 즉 두 개의 시간이 교차되는 특별한 구조를 가지고 있다. 하나는 허 생원과 동이, 조 선달이 봉평에서 대화 장으로 가는 현재의 이야기이다.

다른 하나는 허 생원이 젊은 시절 성 서방네 처녀와 만났던 과거의 이야기이다. 현재 이야기에서는 허 생원이 동이를 꾸짖다가 화해하고, 점점 동이가 자신의 아들일지도 모른다는 생각을 하게 된다.

과거 이야기에서는 허 생원이 달밤에 물방앗간에서 성 서방네 처녀와 우연히 만나 하룻밤을 보낸 추억을 들려준다. 이 두 이야기는 따로 놀지 않고 서로 연결되어 있다.

마지막에 동이가 허 생원처럼 '왼손잡이'라는 것이 밝혀지면서 과거의 그 처녀가 동이의 어머니이고, 허 생원이 동이의 아버지일 가능성을 강하게 보여준다.
}

\begin{kfQBlock}{1}{문제}
\kfQStem{현재 이야기에서 허 생원 일행이 가고 있는 목적지는 어디인가?}
\kfAnswerLines{1}
\end{kfQBlock}

\begin{kfQBlock}{2}{문제}
\kfQStem{과거 이야기에서 허 생원이 성 서방네 처녀와 인연을 맺은 장소는 어디인가?}
\kfAnswerLines{1}
\end{kfQBlock}

\begin{kfQBlock}{3}{문제}
\kfQStem{마지막에 밝혀진 단서를 통해, 과거의 성 서방네 처녀는 누구의 어머니일 가능성이 높아졌는가?}
\kfAnswerLines{1}
\end{kfQBlock}

\begin{kfQBlock}{4}{문제}
\kfQStem{허 생원과 동이가 부자 관계일 수 있음을 보여주는 가장 결정적인 신체적 특징은 무엇인가?}
\kfAnswerLines{1}
\end{kfQBlock}

\par\noindent\textbf{\color{kfPrimary} 해설 단원 3}\par\smallskip
\kfPassageNote{해설 3}{%
허 생원은 겉으로는 무뚝뚝하고 거친 장돌뱅이처럼 보이지만, 속으로는 매우 순수하고 낭만적인 마음을 가진 사람이다. 평생 여자와 인연이 없었다고 하면서도 젊은 시절 단 하룻밤의 추억을 소중히 간직하고 있다.

처음에 허 생원은 동이가 주막에서 여자와 농담치는 모습을 보고 화를 내며 따귀까지 때린다. 하지만 동이가 말대꾸 없이 조용히 나가는 모습을 보고 미안한 마음이 든다.

그 후 동이가 도망간 나귀를 알려주자 고마운 마음을 갖게 된다. 여행을 하면서 허 생원은 동이에게 사과하고, 동이의 이야기를 들으면서 점점 가까워진다.

동이는 아버지 없이 어머니 혼자 키운 아이로, 어머니에 대한 효심이 깊고 순박한 청년이다. 둘 사이의 갈등이 화해로 바뀌면서 부자의 정을 느끼게 되는 과정이 자연스럽게 그려진다.
}

\begin{kfQBlock}{1}{문제}
\kfQStem{허 생원은 겉으로는 무뚝뚝하지만, 속으로는 어떤 마음을 가진 사람인가?}
\kfAnswerLines{1}
\end{kfQBlock}

\begin{kfQBlock}{2}{문제}
\kfQStem{허 생원은 주막에서 동이에게 화가 나 어떤 행동을 했는가?}
\kfAnswerLines{1}
\end{kfQBlock}

\begin{kfQBlock}{3}{문제}
\kfQStem{동이가 말대꾸 없이 조용히 나가자 허 생원은 어떤 마음이 들었는가?}
\kfAnswerLines{1}
\end{kfQBlock}

\begin{kfQBlock}{4}{문제}
\kfQStem{도망간 나귀를 알려준 동이에게 허 생원은 어떤 마음을 갖게 되었는가?}
\kfAnswerLines{1}
\end{kfQBlock}

\begin{kfQBlock}{5}{문제}
\kfQStem{동이는 아버지 없이 자라면서 누구에 대한 효심이 깊은 청년인가?}
\kfAnswerLines{1}
\end{kfQBlock}

\begin{kfQBlock}{6}{문제}
\kfQStem{두 사람이 점점 가까워지면서 무엇을 느끼게 되는가?}
\kfAnswerLines{1}
\end{kfQBlock}

\kfSecHeader{kfAreaLiterature}{문제}{}

\begin{kfQBlock}{1}{}
\kfQStem{허 생원이 주막에서 동이에게 화를 내고 때린 가장 큰 이유는 무엇인가?}
\begin{kfChoices}
\kfChoice 동이가 빌려 간 돈을 갚지 않아서
\kfChoice 동이가 허 생원의 나귀를 괴롭혔기 때문에
\kfChoice 동이가 자신의 오랜 친구인 조 선달에게 무례하게 굴어서
\kfChoice 어린 동이가 여자와 농담하는 모습이 꼴사납고 못마땅해서
\kfChoice 동이가 장사를 제대로 하지 않고 게으름을 피워서
\end{kfChoices}
\end{kfQBlock}
\begin{kfQBlock}{2}{}
\kfQStem{이 소설에서 ‘달밤과 메밀꽃밭’이라는 배경이 하는 역할로 적절하지 \uline{않은} 것은 무엇인가?}
\begin{kfChoices}
\kfChoice 앞으로 일어날 위험하고 잔혹한 사건을 강하게 암시한다.
\kfChoice 흰빛의 메밀꽃과 달빛이 어우러진 서정적·환상적 정서를 환기한다.
\kfChoice 인물의 회상과 현재가 자연스럽게 이어지도록 돕는 매개 공간이다.
\kfChoice 동이와 허 생원 사이의 정서적 거리를 좁혀 주는 배경이 된다.
\kfChoice 허 생원이 과거의 사랑 이야기를 떠올리게 하는 낭만적인 분위기를 만든다.
\end{kfChoices}
\end{kfQBlock}
\begin{kfQBlock}{3}{}
\kfQStem{허 생원이 평생 잊지 못하고 소중하게 간직하고 있는 추억은 무엇인가?}
\begin{kfChoices}
\kfChoice 젊은 시절 장사를 해서 아주 큰돈을 한꺼번에 벌었던 일
\kfChoice 친구 조 선달을 처음 만나 평생의 우정을 쌓게 된 일
\kfChoice 돌아가신 어머니가 정성껏 해주시던 음식을 먹었던 일
\kfChoice 봉평의 물방앗간에서 성 서방네 처녀와 인연을 맺었던 일
\kfChoice 자신의 나귀를 타고 처음으로 장터에 함께 나갔던 일
\end{kfChoices}
\end{kfQBlock}
\begin{kfQBlock}{4}{}
\kfQStem{허 생원의 이야기를 듣는 조 선달의 역할로 적절하지 \uline{않은} 것은 무엇인가?}
\begin{kfChoices}
\kfChoice 허 생원의 이야기가 거짓임을 지적하며 강하게 비판하는 역할을 한다.
\kfChoice 청자로서 인물의 회상을 부드럽게 끌어내는 보조적 역할을 한다.
\kfChoice 허 생원의 이야기를 들어 줌으로써 인물의 내면을 드러내도록 돕는다.
\kfChoice 인물 사이의 대화 구조를 통해 사건 전개에 자연스러움을 더한다.
\kfChoice 허 생원의 말에 맞장구를 쳐 주며 이야기가 자연스럽게 이어지게 한다.
\end{kfChoices}
\end{kfQBlock}
\begin{kfQBlock}{5}{}
\kfQStem{허 생원이 동이에게 친밀감과 정을 느끼게 되는 결정적인 계기는 무엇인가?}
\begin{kfChoices}
\kfChoice 동이가 허 생원의 도망간 나귀를 찾아주었을 때
\kfChoice 동이가 허 생원의 매우 무거운 짐을 대신 들어주었을 때
\kfChoice 동이가 자신에게 아버지가 없다고 고백했을 때
\kfChoice 허 생원이 개울에 빠졌을 때 동이가 그를 업어주었을 때
\kfChoice 동이가 허 생원에게 맛있는 음식을 대접했을 때
\end{kfChoices}
\end{kfQBlock}
\begin{kfQBlock}{6}{}
\kfQStem{동이가 허 생원의 아들일지도 모른다고 암시하는 단서로 적절하지 \uline{않은} 것은 무엇인가?}
\begin{kfChoices}
\kfChoice 동이가 허 생원처럼 왼손잡이라는 점
\kfChoice 동이가 허 생원과 거울에 비춘 듯 똑같이 닮은 얼굴을 하고 있다는 점
\kfChoice 어머니가 ‘제천’으로 옮겨 가게 된 사연이 허 생원의 옛사랑과 겹치는 점
\kfChoice 허 생원이 동이에게 묘한 친밀감과 정을 느끼게 되는 점
\kfChoice 동이의 어머니가 봉평 출신이라는 점
\end{kfChoices}
\end{kfQBlock}
\begin{kfQBlock}{7}{}
\kfQStem{허 생원이 동이와 함께 가기로 결심한 새로운 목적지는 어디인가?}
\begin{kfChoices}
\kfChoice 동이의 어머니가 계신 제천
\kfChoice 동이 어머니의 고향인 봉평
\kfChoice 큰 장이 열리는 충주
\kfChoice 자신의 고향인 강릉
\kfChoice 돈을 많이 벌 수 있는 서울
\end{kfChoices}
\end{kfQBlock}
\begin{kfQBlock}{8}{}
\kfQStem{이 글의 표현상 특징으로 알맞지 \uline{않은} 것은 무엇인가?}
\begin{kfChoices}
\kfChoice "소금을 뿌린 듯이"와 같은 비유적 표현을 사용한다.
\kfChoice "짐승 같은 달의 숨소리"와 같은 감각적인 표현을 사용한다.
\kfChoice 인물 간의 대화를 통해 사건을 자연스럽게 전개한다.
\kfChoice 어려운 한자어를 많이 사용하여 지적인 분위기를 만든다.
\kfChoice 현재의 사건 속에 과거의 회상을 삽입하여 인물의 사연을 깊이 있게 보여준다.
\end{kfChoices}
\end{kfQBlock}
\begin{kfQBlock}{1}{}
\kfQStem{이 소설에서 ‘달밤의 메밀꽃밭’이라는 배경이 하는 역할을 <조건>에 맞게 서술하시오.}
\kfAnswerNote{답안}{4}
\end{kfQBlock}
\begin{kfQBlock}{2}{}
\kfQStem{허 생원은 어떤 단서들을 통해 동이가 자신의 아들일지도 모른다고 짐작하게 되는지 두 가지를 서술하시오.}
\kfAnswerNote{답안}{4}
\end{kfQBlock}


\kfStartNonlit{비문학 영역 — [사회] 어떤 방법이 가장 공정할까? 여러 가지 투표 방식}


\kfSecHeader{kfAreaNonlit}{지문}{[사회] 어떤 방법이 가장 공정할까? 여러 가지 투표 방식}
\kfPassageNote{[사회] 어떤 방법이 가장 공정할까? 여러 가지 투표 방식}{%
우리 마을에 병원, 학교, 경찰서 중 하나를 짓는다고 상상해 보자. 사람들마다 원하는 것이 다르기 때문에, 모두의 생각을 하나로 모으는 과정이 필요하다. 이처럼 여러 사람의 의견을 모아 단체의 최종 결정을 내리는 방법에는 여러 가지가 있다. 각 투표 방식의 특징은 무엇이고 어떤 장점과 단점이 있는지 알아보자.   가장 널리 쓰이는 투표 방식은 단순 과반수제이다. 투표에 참여한 사람 중 절반이 넘는 사람이 찬성하는 의견을 선택하는 방법이다. 가장 간단하고 빠르게 결정을 내릴 수 있다는 장점이 있다. 하지만 이 방식은 사람들이 각 의견을 얼마나 강하게 원하거나 싫어하는지는 보여주지 못한다. 그래서 다수의 사람들이 찬성했다는 이유만으로 소수의 사람들이 크게 피해를 보는 결정이 내려질 수 있다. 또한, 어떤 후보끼리 먼저 대결을 붙이느냐에 따라 최종 결과가 달라지는 현상이 나타나기도 한다.   점수 투표제는 사람들에게 일정한 점수를 나누어 주고, 원하는 의견에 자유롭게 점수를 나누어 투표하게 하는 방식이다. 예를 들어 10점을 가지고 병원에 7점, 학교에 3점을 주는 식이다. 이 방법은 각자가 의견을 얼마나 좋아하는지 그 마음의 크기를 자세히 표현할 수 있다. 하지만 다른 사람의 투표 결과를 예상하고, 이기기 위해 자신의 진짜 생각과 다르게 점수를 몰아주는 전략적인 행동이 나타날 수 있다는 단점이 있다.   보르다 투표제는 모든 후보의 순위를 매겨서 점수를 주는 방식이다. 후보가 세 명이면 1등에게 3점, 2등에게 2점, 3등에게 1점을 주는 식이다. 그 후 모든 점수를 합쳐 가장 높은 점수를 얻은 후보가 선택된다. 이 방식은 한두 명에게만 인기가 많은 의견보다는, 여러 사람에게 두루두루 괜찮다고 평가받는 무난한 의견이 뽑힐 가능성이 높다.   이처럼 투표에는 다양한 방법이 있으며 각각 장점과 단점이 뚜렷하다. 따라서 어떤 상황에서 투표를 하는지에 따라 가장 알맞은 방식을 선택하는 지혜가 필요하다.
}


\kfSecHeader{kfAreaNonlit}{활동}{어휘 학습}
\noindent\textit{\small 다음 표의 '의미'를 읽고, 그에 해당하는 '어휘'를 지문에서 찾아 적으시오.}\par\medskip
\begin{kfVocab}
\kfVocabItem{1}{전체 수의 절반을 넘는 수}
\kfVocabItem{2}{수가 많은 쪽}
\kfVocabItem{3}{수가 적은 쪽}
\kfVocabItem{4}{원하는 의견에 자유롭게 점수를 나누어 투표하게 하는 방식}
\kfVocabItem{5}{이기기 위해 자신의 진짜 생각과 다르게 행동하는 것}
\kfVocabItem{6}{모든 후보의 순위를 매겨서 점수를 주는 방식}
\kfVocabItem{7}{특별히 좋지도 나쁘지도 않은 평범한 것}
\end{kfVocab}

\kfSecHeader{kfAreaNonlit}{활동}{비문학 문장 독해}
\noindent\textit{\small 지문을 문장 단위로 정독하면서, 문장별로 주어진 물음에 답해 보자.}\par\medskip
\begin{kfSentenceUnit}{1}{우리 마을에 병원, 학교, 경찰서 중 하나를 짓는다고 상상해 보자.}
\kfSentQ{1.}{이 문장은 독자에게 어떤 상황을 상상해 보라고 제안하는가?}
\begin{kfChoices}
\kfChoice 자신이 투표에 당선되는 상황
\kfChoice 마을에 필요한 건물 중 하나를 짓는 상황
\kfChoice 새로운 법을 만드는 상황
\kfChoice 마을 대표를 뽑는 상황
\end{kfChoices}
\kfSentQ{2.}{이 문장을 통해 앞으로 어떤 내용이 나올 것을 예상할 수 있는가?}
\begin{kfChoices}
\kfChoice 마을의 문제점을 해결하는 방법
\kfChoice 여러 기부금을 모으는 방법
\kfChoice 여러 선택지 중 하나를 고르는 방법에 대한 설명
\kfChoice 건물을 짓는 순서에 대한 설명
\end{kfChoices}
\end{kfSentenceUnit}

\begin{kfSentenceUnit}{2}{사람들마다 원하는 것이 다르기 때문에, 모두의 생각을 하나로 모으는 과정이 필요하다.}
\kfSentQ{3.}{이 문장에서 '원인'에 해당하는 내용은?}
\begin{kfChoices}
\kfChoice 사람들마다 원하는 것이 다르기 때문에
\kfChoice 모두의 생각을 하나로 모으는 과정이 필요하다
\kfChoice 병원이나 학교를 지어야 한다
\kfChoice 투표를 해야 한다
\end{kfChoices}
\kfSentQ{4.}{모두의 생각을 하나로 모으는 과정이 필요한 이유는 무엇인가?}
\begin{kfChoices}
\kfChoice 건물을 짓는 비용이 부족해서
\kfChoice 사람들마다 원하는 것이 다르기 때문이다.
\kfChoice 시간이 부족해서
\kfChoice 법으로 정해져 있어서
\end{kfChoices}
\end{kfSentenceUnit}

\begin{kfSentenceUnit}{3}{이처럼 여러 사람의 의견을 모아 단체의 최종 결정을 내리는 방법에는 여러 가지가 있다.}
\kfSentQ{5.}{'이처럼'이 가리키는 것은 무엇인가?}
\begin{kfChoices}
\kfChoice 모두의 생각을 하나로 모으는 과정
\kfChoice 사람들끼리 서로 싸우는 과정
\kfChoice 건물을 짓는 과정
\kfChoice 투표를 거부하는 과정
\end{kfChoices}
\kfSentQ{6.}{여러 가지가 있다고 한 것은 무엇인가?}
\begin{kfChoices}
\kfChoice 마을에 필요한 건물의 종류
\kfChoice 마을 사람들의 직업
\kfChoice 여러 사람의 의견을 모아 단체의 최종 결정을 내리는 방법
\kfChoice 투표에 참여하는 사람들의 나이
\end{kfChoices}
\end{kfSentenceUnit}

\begin{kfSentenceUnit}{4}{각 투표 방식의 특징은 무엇이고 어떤 장점과 단점이 있는지 알아보자.}
\kfSentQ{7.}{이 글이 앞으로 설명할 내용으로 알맞은 것은?}
\begin{kfChoices}
\kfChoice 투표의 역사와 유래
\kfChoice 각 투표 방식의 특징, 장점, 단점
\kfChoice 민주주의의 발전 과정
\kfChoice 선거 운동의 방법
\end{kfChoices}
\end{kfSentenceUnit}

\begin{kfSentenceUnit}{5}{가장 널리 쓰이는 투표 방식은 단순 과반수제이다.}
\kfSentQ{8.}{가장 널리 쓰이는 투표 방식의 이름은 무엇인가?}
\begin{kfChoices}
\kfChoice 점수 투표제
\kfChoice 보르다 투표제
\kfChoice 단순 과반수제
\kfChoice 만장일치제
\end{kfChoices}
\end{kfSentenceUnit}

\begin{kfSentenceUnit}{6}{투표에 참여한 사람 중 절반이 넘는 사람이 찬성하는 의견을 선택하는 방법이다.}
\kfSentQ{9.}{단순 과반수제에서 의견이 선택되기 위한 조건은 무엇인가?}
\begin{kfChoices}
\kfChoice 모든 사람이 찬성하는 것
\kfChoice 투표 참여자 중 절반이 넘는 사람이 찬성하는 것
\kfChoice 가장 많은 점수를 받는 것
\kfChoice 순위가 가장 높은 것
\end{kfChoices}
\kfSentQ{10.}{'절반이 넘는다'는 것을 다른 말로 표현하면 무엇인가?}
\begin{kfChoices}
\kfChoice 소수
\kfChoice 다수
\kfChoice 과반수
\kfChoice 전원
\end{kfChoices}
\end{kfSentenceUnit}

\begin{kfSentenceUnit}{7}{가장 간단하고 빠르게 결정을 내릴 수 있다는 장점이 있다.}
\kfSentQ{11.}{단순 과반수제의 장점은 무엇인가?}
\begin{kfChoices}
\kfChoice 소수 의견을 잘 반영한다.
\kfChoice 가장 공정한 방법이다.
\kfChoice 간단하고 빠르게 결정을 내릴 수 있다.
\kfChoice 점수를 자유롭게 줄 수 있다.
\end{kfChoices}
\end{kfSentenceUnit}

\begin{kfSentenceUnit}{8}{하지만 이 방식은 사람들이 각 의견을 얼마나 강하게 원하거나 싫어하는지는 보여주지 못한다.}
\kfSentQ{12.}{'하지만'이 나타내는 관계는 무엇인가?}
\begin{kfChoices}
\kfChoice 앞 내용에 대한 예시
\kfChoice 앞 내용과 이어지는 결과
\kfChoice 앞 내용과 반대되는 내용
\kfChoice 앞 내용에 대한 원인
\end{kfChoices}
\kfSentQ{13.}{'이 방식'이 가리키는 것은 무엇인가?}
\begin{kfChoices}
\kfChoice 점수 투표제
\kfChoice 보르다 투표제
\kfChoice 단순 과반수제
\kfChoice 다수결의 원칙
\end{kfChoices}
\kfSentQ{14.}{단순 과반수제가 보여주지 못하는 것은 무엇인가?}
\begin{kfChoices}
\kfChoice 찬성하는 사람의 수
\kfChoice 사람들이 의견을 얼마나 강하게 원하거나 싫어하는지
\kfChoice 반대하는 사람의 수
\kfChoice 투표에 참여한 총 인원 수
\end{kfChoices}
\end{kfSentenceUnit}

\begin{kfSentenceUnit}{9}{그래서 다수의 사람들이 찬성했다는 이유만으로 소수의 사람들이 크게 피해를 보는 결정이 내려질 수 있다.}
\kfSentQ{15.}{'그래서'의 구체적인 내용은 무엇인가?}
\begin{kfChoices}
\kfChoice 투표 시간이 너무 오래 걸려서
\kfChoice 사람들의 선호 강도를 보여주지 못해서
\kfChoice 후보가 너무 많아서
\kfChoice 계산 방법이 복잡해서
\end{kfChoices}
\kfSentQ{16.}{단순 과반수제 때문에 생길 수 있는 문제는 무엇인가?}
\begin{kfChoices}
\kfChoice 결정이 너무 늦어질 수 있다.
\kfChoice 소수의 사람들이 크게 피해를 볼 수 있다.
\kfChoice 비용이 많이 든다.
\kfChoice 투표율이 낮아질 수 있다.
\end{kfChoices}
\kfSentQ{17.}{단순 과반수제에서 소수의 사람들이 피해를 볼 수 있는 이유는 무엇인가?}
\begin{kfChoices}
\kfChoice 다수결로만 결정되기 때문에
\kfChoice 소수의 의견이 너무 강력해서
\kfChoice 투표 방식이 너무 복잡해서
\kfChoice 후보자가 한 명뿐이라서
\end{kfChoices}
\end{kfSentenceUnit}

\begin{kfSentenceUnit}{10}{또한, 어떤 후보끼리 먼저 대결을 붙이느냐에 따라 최종 결과가 달라지는 현상이 나타나기도 한다.}
\kfSentQ{18.}{무엇이 달라지면 최종 결과가 바뀔 수 있는가?}
\begin{kfChoices}
\kfChoice 투표 용지의 색깔
\kfChoice 투표 장소
\kfChoice 후보끼리 대결하는 순서
\kfChoice 투표함의 모양
\end{kfChoices}
\kfSentQ{19.}{단순 과반수제의 또 다른 문제점은 무엇인가?}
\begin{kfChoices}
\kfChoice 대결 순서에 따라 결과가 달라질 수 있다.
\kfChoice 모든 사람이 만족하는 결과를 얻기 힘들다.
\kfChoice 투표 집계에 시간이 오래 걸린다.
\kfChoice 비밀 투표가 보장되지 않는다.
\end{kfChoices}
\end{kfSentenceUnit}

\begin{kfSentenceUnit}{11}{점수 투표제는 사람들에게 일정한 점수를 나누어 주고, 원하는 의견에 자유롭게 점수를 나누어 투표하게 하는 방식이다.}
\kfSentQ{20.}{'점수 투표제'는 무엇을 어떻게 나누어 주는 방식인가?}
\begin{kfChoices}
\kfChoice 순위를 매겨 점수를 주는 방식
\kfChoice 한 사람에게 한 표씩 주는 방식
\kfChoice 일정한 점수를 원하는 의견에 자유롭게 나누어 주는 방식
\kfChoice 무작위로 점수를 배정하는 방식
\end{kfChoices}
\end{kfSentenceUnit}

\begin{kfSentenceUnit}{12}{예를 들어 10점을 가지고 병원에 7점, 학교에 3점을 주는 식이다.}
\kfSentQ{21.}{이 문장은 앞 문장의 어떤 내용에 대한 구체적인 예시인가?}
\begin{kfChoices}
\kfChoice 단순 과반수제
\kfChoice 점수 투표제
\kfChoice 보르다 투표제
\kfChoice 다수결의 원칙
\end{kfChoices}
\end{kfSentenceUnit}

\begin{kfSentenceUnit}{13}{이 방법은 각자가 의견을 얼마나 좋아하는지 그 마음의 크기를 자세히 표현할 수 있다.}
\kfSentQ{22.}{'이 방법'이 가리키는 것은 무엇인가?}
\begin{kfChoices}
\kfChoice 단순 과반수제
\kfChoice 점수 투표제
\kfChoice 보르다 투표제
\kfChoice 비밀 투표제
\end{kfChoices}
\kfSentQ{23.}{점수 투표제의 장점은 무엇인가?}
\begin{kfChoices}
\kfChoice 가장 빠르게 결과를 알 수 있다.
\kfChoice 의견을 얼마나 좋아하는지 마음의 크기를 자세히 표현할 수 있다.
\kfChoice 계산이 매우 간단하다.
\kfChoice 소수 의견을 무시할 수 있다.
\end{kfChoices}
\end{kfSentenceUnit}

\begin{kfSentenceUnit}{14}{하지만 다른 사람의 투표 결과를 예상하고, 이기기 위해 자신의 진짜 생각과 다르게 점수를 몰아주는 전략적인 행동이 나타날 수 있다는 단점이 있다.}
\kfSentQ{24.}{'하지만'이 나타내는 관계는 무엇인가?}
\begin{kfChoices}
\kfChoice 앞 문장에 대한 구체적인 예시
\kfChoice 앞 문장과 대조되는 내용
\kfChoice 앞 문장의 원인
\kfChoice 앞 문장과 동일한 내용 반복
\end{kfChoices}
\kfSentQ{25.}{점수 투표제의 단점은 무엇인가?}
\begin{kfChoices}
\kfChoice 동점자가 나올 경우 해결하기 어렵다.
\kfChoice 투표 용지를 만드는 비용이 많이 든다.
\kfChoice 자신의 진짜 생각과 다르게 점수를 몰아주는 전략적인 행동이 나타날 수 있다.
\kfChoice 모든 사람의 의견을 듣지 못한다.
\end{kfChoices}
\kfSentQ{26.}{점수 투표제에서 나타날 수 있는 '전략적인 행동'이란 무엇인가?}
\begin{kfChoices}
\kfChoice 다른 사람의 투표를 방해하는 행동
\kfChoice 투표 결과를 조작하는 행동
\kfChoice 여러 후보에게 골고루 점수를 주는 행동
\kfChoice 자신의 진짜 생각과 다르게 이기기 위해 점수를 몰아주는 행동
\end{kfChoices}
\kfSentQ{27.}{전략적인 행동이 나타나는 이유는 무엇인가?}
\begin{kfChoices}
\kfChoice 다른 사람의 투표 결과를 예상하고 이기기 위해서
\kfChoice 투표 규칙을 잘 몰라서
\kfChoice 실수로 잘못 투표해서
\kfChoice 투표 시간을 줄이기 위해서
\end{kfChoices}
\end{kfSentenceUnit}

\begin{kfSentenceUnit}{15}{보르다 투표제는 모든 후보의 순위를 매겨서 점수를 주는 방식이다.}
\kfSentQ{28.}{'보르다 투표제'는 어떻게 점수를 주는 방식인가?}
\begin{kfChoices}
\kfChoice 가장 좋아하는 후보 한 명에게만 점수를 주는 방식
\kfChoice 모든 후보의 순위를 매겨서 점수를 주는 방식
\kfChoice 싫어하는 후보에게 마이너스 점수를 주는 방식
\kfChoice 무작위로 점수를 배정하는 방식
\end{kfChoices}
\end{kfSentenceUnit}

\begin{kfSentenceUnit}{16}{후보가 세 명이면 1등에게 3점, 2등에게 2점, 3등에게 1점을 주는 식이다.}
\kfSentQ{29.}{이 문장은 앞 문장의 어떤 내용에 대한 설명인가?}
\begin{kfChoices}
\kfChoice 보르다 투표제에 대한 구체적인 예시
\kfChoice 보르다 투표제의 단점
\kfChoice 보르다 투표제의 역사
\kfChoice 보르다 투표제의 문제점
\end{kfChoices}
\end{kfSentenceUnit}

\begin{kfSentenceUnit}{17}{그 후 모든 점수를 합쳐 가장 높은 점수를 얻은 후보가 선택된다.}
\kfSentQ{30.}{보르다 투표제에서 최종적으로 선택되는 후보는 어떤 후보인가?}
\begin{kfChoices}
\kfChoice 1등 표를 가장 많이 받은 후보
\kfChoice 가장 적은 점수를 받은 후보
\kfChoice 절반 이상의 찬성을 얻은 후보
\kfChoice 모든 점수를 합쳐 가장 높은 점수를 얻은 후보
\end{kfChoices}
\end{kfSentenceUnit}

\begin{kfSentenceUnit}{18}{이 방식은 한두 명에게만 인기가 많은 의견보다는, 여러 사람에게 두루두루 괜찮다고 평가받는 무난한 의견이 뽑힐 가능성이 높다.}
\kfSentQ{31.}{'이 방식'이 가리키는 것은 무엇인가?}
\begin{kfChoices}
\kfChoice 단순 과반수제
\kfChoice 점수 투표제
\kfChoice 보르다 투표제
\kfChoice 만장일치제
\end{kfChoices}
\kfSentQ{32.}{보르다 투표제의 특징은 무엇인가?}
\begin{kfChoices}
\kfChoice 가장 인기가 많은 의견이 항상 뽑힌다.
\kfChoice 여러 사람에게 두루두루 괜찮다고 평가받는 무난한 의견이 뽑힐 가능성이 높다.
\kfChoice 소수의 의견이 완전히 무시된다.
\kfChoice 계산 방법이 가장 간단하다.
\end{kfChoices}
\kfSentQ{33.}{보르다 투표제에서 뽑힐 가능성이 높은 의견의 특징은 무엇인가?}
\begin{kfChoices}
\kfChoice 호불호가 강하게 갈리는 의견
\kfChoice 소수의 사람만 좋아하는 의견
\kfChoice 여러 사람에게 무난하다고 평가받는 의견
\kfChoice 반대하는 사람이 많은 의견
\end{kfChoices}
\end{kfSentenceUnit}

\begin{kfSentenceUnit}{19}{이처럼 투표에는 다양한 방법이 있으며 각각 장점과 단점이 뚜렷하다.}
\kfSentQ{34.}{'이처럼'이 가리키는 것은 무엇인가?}
\begin{kfChoices}
\kfChoice 투표의 단점들
\kfChoice 지금까지 설명한 다양한 투표 방법들
\kfChoice 투표의 역사적 유래
\kfChoice 선거 운동 방식
\end{kfChoices}
\kfSentQ{35.}{단순 과반수제, 점수 투표제, 보르다 투표제의 공통적인 특징은 무엇인가?}
\begin{kfChoices}
\kfChoice 모두 완벽한 투표 방식이다.
\kfChoice 각각 장점과 단점이 뚜렷하다.
\kfChoice 모두 계산하기 복잡하다.
\kfChoice 모두 소수 의견을 무시한다.
\end{kfChoices}
\end{kfSentenceUnit}

\begin{kfSentenceUnit}{20}{따라서 어떤 상황에서 투표를 하는지에 따라 가장 알맞은 방식을 선택하는 지혜가 필요하다.}
\kfSentQ{36.}{투표 방식을 선택할 때 고려해야 할 기준은 무엇인가?}
\begin{kfChoices}
\kfChoice 투표하는 시간
\kfChoice 투표 용지의 크기
\kfChoice 어떤 상황에서 투표를 하는지
\kfChoice 투표 관리자의 선호도
\end{kfChoices}
\end{kfSentenceUnit}

\kfSecHeader{kfAreaNonlit}{활동}{구조도}
\noindent\textit{\small 글의 흐름을 살펴보고, 구조도의 빈칸을 알맞게 채워 보자.}\par\medskip
\par\medskip\needspace{6\baselineskip}
\begin{tcolorbox}[enhanced, breakable=false,
  colback=kfPaper, colframe=kfRule, boxrule=0.4pt, arc=2pt,
  left=10pt, right=10pt, top=8pt, bottom=8pt,
  title={\footnotesize\bfseries\color{kfMute} 글의 구조도},
  coltitle=kfMute, colbacktitle=kfPrimaryLight!40,
  attach boxed title to top left={yshift=-1.5mm, xshift=6pt},
  boxed title style={colframe=kfRule, boxrule=0.3pt, arc=1pt}]
\setstretch{1.4}\footnotesize
지침: 글의 전체적인 논리 흐름을 파악하며 빈칸에 알맞은 내용을 채우시오. \\

\noindent \textbullet\ 1문단: 투표 방식의 필요성 \\
\noindent \quad\textbullet\ 배경: 사람들마다 ( ① )이 다름 \\
\noindent \quad\textbullet\ 필요성: ( ② )을 하나로 모으는 과정이 필요 \\
\noindent \quad\textbullet\ 주제 소개: 다양한 ( ③ ) 방식의 특징과 장단점 \\
\noindent \textbullet\ 2문단: 방법 1 - ( ④ ) 과반수제 \\
\noindent \quad\textbullet\ 방법: ( ⑤ )이 넘는 찬성으로 결정 \\
\noindent \quad\textbullet\ 장점: 간단하고 ( ⑥ ) \\
\noindent \quad\textbullet\ 단점 1: 좋아하는 크기를 반영하지 못해 ( ⑦ ) 의 사람들이 피해를 봄 \\
\noindent \quad\textbullet\ 단점 2: 대결 순서에 따라 ( ⑧ )가 바뀔 수 있음 \\
\noindent \textbullet\ 3문단: 방법 2 - ( ⑨ ) 투표제 \\
\noindent \quad\textbullet\ 방법: 주어진 ( ⑩ )를 자유롭게 나누어 투표 \\
\noindent \quad\textbullet\ 장점: 각자가 의견을 얼마나 좋아하는지 그 마음의 ( ⑪ )를 자세히 표현할 수 있음  단점: ( ⑫ )인 행동이 나타날 수 있음 \\
\noindent \textbullet\ 4문단: 방법 3 - ( ⑬ ) 투표제 \\
\noindent \quad\textbullet\ 방법: 모든 후보의 ( ⑭ )를 매겨 점수 부여 \\
\noindent \quad\textbullet\ 특징: ( ⑮ )한 의견이 뽑힐 가능성이 높음 \\
\noindent \textbullet\ 5문단: 결론 \\
\noindent \quad\textbullet\ 요약: 투표 방식마다 ( ⑯ )과 단점이 뚜렷함 \\
\noindent \quad\textbullet\ 주장: ( ⑰ )에 맞는 방식을 선택하는 지혜가 필요 \\
\end{tcolorbox}\par\medskip
\kfSecHeader{kfAreaNonlit}{문제}{}

\begin{kfQBlock}{1}{}
\kfQStem{윗글의 중심 내용으로 적절하지 \uline{않은} 것은?}
\begin{kfChoices}
\kfChoice 투표의 역사와 민주주의의 전반적인 발전 과정
\kfChoice 다수결 투표·결선 투표·보르다 투표제 같은 다양한 의사결정 방법
\kfChoice 상황에 따라 적합한 투표 방식이 달라진다는 사실
\kfChoice 각 투표 방식이 가질 수 있는 장점과 한계
\kfChoice 여러 가지 투표 방식의 특징과 장단점 비교
\end{kfChoices}
\end{kfQBlock}
\begin{kfQBlock}{2}{}
\kfQStem{윗글의 내용과 일치하지 \uline{않는} 것은?}
\begin{kfChoices}
\kfChoice 단순 과반수제는 가장 간단하고 빠른 결정 방식이다.
\kfChoice 점수 투표제는 사람들이 의견을 얼마나 좋아하는지 알 수 있다.
\kfChoice 보르다 투표제는 가장 인기가 많은 의견이 항상 뽑히도록 보장한다.
\kfChoice 단순 과반수제는 소수의 의견이 무시될 수 있는 위험이 있다.
\kfChoice 투표 방식을 선택할 때는 상황을 고려하는 지혜가 필요하다.
\end{kfChoices}
\end{kfQBlock}
\begin{kfQBlock}{3}{}
\kfQStem{<보기>의 상황에서 나타날 수 있는 문제점을 윗글에서 찾은 것으로 가장 적절한 것은?}
\begin{kfEvidence}우리 반 30명 중 16명이 점심 메뉴로 떡볶이를 골랐고, 14명은 피자를 골랐다. 그런데 피자를 고른 친구들 중에는 떡볶이를 아주 싫어하는 친구들이 많았다.\end{kfEvidence}
\begin{kfChoices}
\kfChoice 점수 투표제의 전략적인 행동
\kfChoice 보르다 투표제의 무난한 의견 선택
\kfChoice 단순 과반수제에서 소수의 의견이 무시되는 문제
\kfChoice 점수 투표제에서 점수를 나누기 어려운 문제
\kfChoice 보르다 투표제에서 순위를 매기기 힘든 문제
\end{kfChoices}
\end{kfQBlock}
\begin{kfQBlock}{4}{}
\kfQStem{다음 중 '보르다 투표제'에 대한 설명으로 옳은 것은?}
\begin{kfChoices}
\kfChoice 투표자의 절반이 넘는 표를 얻어야 한다.
\kfChoice 모든 후보의 순위를 매겨 점수를 합산한다.
\kfChoice 각자 가진 점수를 원하는 후보에게 몰아줄 수 있다.
\kfChoice 가장 간단하고 빠르게 결정을 내릴 수 있다.
\kfChoice 가장 인기 있는 후보가 항상 유리하다.
\end{kfChoices}
\end{kfQBlock}
\begin{kfQBlock}{5}{}
\kfQStem{각 투표 방식과 그 특징을 바르게 짝지은 것은?}
\begin{kfChoices}
\kfChoice 단순 과반수제 - 선호도의 크기를 알 수 있다.
\kfChoice 점수 투표제 - 가장 빠르게 결정할 수 있다.
\kfChoice 보르다 투표제 - 무난한 의견이 선택될 가능성이 높다.
\kfChoice 단순 과반수제 - 전략적인 투표 행동이 나타나기 쉽다.
\kfChoice 점수 투표제 - 소수의 의견이 무시되기 쉽다.
\end{kfChoices}
\end{kfQBlock}
\begin{kfQBlock}{6}{}
\kfQStem{<보기>의 상황에 가장 적합한 투표 방식과 그 이유를 고르시오.}
\begin{kfEvidence}우리 반 학급 회의에서 소풍 장소를 정하려고 한다. 후보지는 놀이공원, 박물관, 수목원 세 곳이다. 한두 명이라도 아주 싫어하는 곳은 피하고, 되도록 많은 친구들이 "그곳도 괜찮아"라고 생각하는 곳으로 정하고 싶다.\end{kfEvidence}
\begin{kfChoices}
\kfChoice 단순 과반수제 - 가장 빠르고 간단하게 결론을 내릴 수 있어서
\kfChoice 점수 투표제 - 자신이 원하는 후보에 점수를 몰아줄 수 있어서
\kfChoice 보르다 투표제 - 다수에게 무난한 평가를 받는 안을 뽑기 좋아서
\kfChoice 단순 과반수제 - 소수의 반대 의견도 확실하게 반영할 수 있어서
\kfChoice 점수 투표제 - 결과를 왜곡하는 전략적인 행동을 막을 수 있어서
\end{kfChoices}
\end{kfQBlock}
\begin{kfQBlock}{1}{}
\kfQStem{여러 사람의 의견을 모아 단체의 최종 결정을 내리는 대표적인 방법을 무엇이라고 하는가?}
\kfAnswerLines{1}
\end{kfQBlock}
\begin{kfQBlock}{2}{}
\kfQStem{가장 널리 쓰이며, 투표 참여자 중 절반이 넘는 사람이 찬성하면 결정되는 방식은 무엇인가?}
\kfAnswerLines{1}
\end{kfQBlock}
\begin{kfQBlock}{3}{}
\kfQStem{단순 과반수제에서 다수의 의견 때문에 피해를 볼 수 있는 사람들은 누구인가?}
\kfAnswerLines{1}
\end{kfQBlock}
\begin{kfQBlock}{4}{}
\kfQStem{사람들에게 일정한 점수를 주고 원하는 의견에 자유롭게 나누어 투표하게 하는 방식은 무엇인가?}
\kfAnswerLines{1}
\end{kfQBlock}
\begin{kfQBlock}{5}{}
\kfQStem{모든 후보의 순위를 매겨서 점수를 주는 투표 방식은 무엇인가?}
\kfAnswerLines{1}
\end{kfQBlock}
\begin{kfQBlock}{1}{}
\kfQStem{'단순 과반수제'와 '점수 투표제'가 각각 가진 가장 큰 장점을 서술하시오.}
\kfAnswerNote{답안}{4}
\end{kfQBlock}



\clearpage

\kfChapterCover{3}{러셀1 (챕터3)}{Chapter 3}{이번 챕터의 학습 내용을 시작합니다.}

\kfStartGrammar{문법 영역 — 단어를 이어주는 풀 (조사)}


\kfSecHeader{kfAreaGrammar}{개념}{단어를 이어주는 풀 (조사)}

단어들을 그냥 순서대로 늘어놓는다고 해서 문장이 되지는 않는다. 예를 들어 ‘동생’, ‘빵’, ‘먹다’처럼 단어만 흩어져 있으면, 누가 무엇을 하는지 관계가 분명하지 않아 어색하다.

하지만 이 단어들 사이에 '이'와 '을'이라는 특별한 풀을 붙여주면 ‘동생이 빵을 먹는다.’라는 자연스러운 문장이 완성된다. 이처럼 다른 말 뒤에 주로 붙어서 그 말과 다른 말과의 문법적 관계를 나타내거나 특별한 의미를 더해주는 단어를 '조사'라고 한다. 여기서 조사 ‘이’는 ‘동생’이 문장의 주인공임을, ‘을’은 ‘빵’이 먹는 행동의 대상임을 명확히 알려준다. 이처럼 문장에서의 문법적 자격을 나타내는 조사를 격조사라고 한다. 이 외에도 조사는 '은', '는', '만', '도'처럼 특별한 의미를 더해주는 보조사나, '와', '과'처럼 두 단어를 이어주는 접속 조사 등이 있다.

이렇게 흩어진 단어를 묶어 문장에 명확한 뜻을 불어넣는 풀, 이것이 바로 ‘조사’다. 조사의 가장 큰 특징은 명사, 대명사, 수사 뒤에 붙어야 의미가 생긴다는 점이다. 풀이 혼자서는 아무것도 붙일 수 없듯이, 조사도 홀로 쓰일 수는 없다.
\par\medskip\begin{itemize}[leftmargin=18pt, itemsep=2pt, topsep=2pt, label=\textbullet]
\item 격조사 — '동생이 빵을 먹는다'(이=주격, 을=목적격)
\item 보조사 — '나는 갔다 / 나도 안다 / 너만 와라'(은·도·만)
\item 접속 조사 — '사과와 배를 샀다'(와)
\end{itemize}

\kfSecHeader{kfAreaGrammar}{활동}{내용 확인}
\noindent\textit{\small 아래 표의 빈칸에 알맞은 말을 채워 문법 내용을 확인해 보자.}\par\medskip
* 단어를 이어주는 풀, 조사 *

아래 표의 빈칸을 채우며 조사의 특징과 종류를 알아보자.

1. 콕콕! 핵심 내용 확인하기

* 단어를 이어주는 풀, 조사 *

아래 표의 빈칸을 채우며 조사의 특징과 종류를 완성해 보세요. 

\par\medskip\needspace{4\baselineskip}
\noindent\begin{adjustbox}{max width=\linewidth, center}
\renewcommand{\arraystretch}{1.45}
\arrayrulecolor{kfRule}
\begin{tabular}{|>{\raggedright\arraybackslash}p{\dimexpr 0.0328\linewidth-12pt\relax}|>{\raggedright\arraybackslash}p{\dimexpr 0.5902\linewidth-12pt\relax}|>{\raggedright\arraybackslash}p{\dimexpr 0.3770\linewidth-12pt\relax}|}
\hline
\rowcolor{kfPrimaryLight!50}\textbf{\color{kfPrimary} 구 분} & \textbf{\color{kfPrimary} 설 명} & \textbf{\color{kfPrimary} 예 시} \\\hline
정의 & [ ① \_\_\_\_\_\_ ]  다른 말 뒤에 붙어 문법적 관계를 나타내거나 특별한 의미를 더해주는 말 & 학교\uline{에}, 밥\uline{을}, 너\uline{만}, 하늘\uline{과} 땅 \\\hline
특징 & 홀로 쓰일 수 없고, 주로 [ ② \_\_\_\_\_\_, \_\_\_\_\_\_, \_\_\_\_\_\_ ] 뒤에 붙어 쓰인다. & 나\uline{도} 간다. \\\hline
종류 & [ ③ \_\_\_\_\_\_ ]  문장에서의 문법적 자격을 나타냄 & 동생\uline{이} 빵\uline{을} 먹는다. \\\hline
\rule[-1.0em]{0pt}{2.4em} & [ ④ \_\_\_\_\_\_ ]  특별한 의미를 더해줌 & 나는 학생, 나\uline{도} 학생, 나\uline{만} 학생 \\\hline
\rule[-1.0em]{0pt}{2.4em} & [ ⑤ \_\_\_\_\_\_ ]  두 단어를 같은 자격으로 이어줌 & 너\uline{와} 나, 사과\uline{랑} 배 \\\hline
\end{tabular}
\arrayrulecolor{black}
\end{adjustbox}\par\medskip

\kfSecHeader{kfAreaGrammar}{활동}{분석 훈련}
\noindent\textit{\small  조사 찾기 훈련! 문장에서 조사에 해당하는 단어를 모두 찾아 쓰시오.}\par\medskip
\begin{enumerate}[leftmargin=22pt, itemsep=6pt, topsep=4pt, label=\textbf{\arabic*.}]
\item 하늘이 푸르고 바다가 넓다.
\item 나는 밥을 먹었다.
\item 철수와 영희는 같은 학교에 다닌다.
\item 이것은 연필이다.
\item 너만 나를 믿어준다.
\item 학교에서 집까지 멀다.
\item 나도 너처럼 할 수 있다.
\item 사과랑 바나나를 샀다.
\item 오늘은 날씨가 좋다.
\item 아버지께서 공원에 가신다.
\end{enumerate}

\kfSecHeader{kfAreaGrammar}{문제}{}

\begin{kfQBlock}{1}{}
\kfQStem{조사에 대한 설명으로 옳지 \uline{않은} 것은?}
\begin{kfChoices}
\kfChoice 명사 뒤에 붙어 쓰인다.
\kfChoice 홀로 쓰일 수 없다.
\kfChoice 문장에서 '풀' 역할을 한다.
\kfChoice 문장의 맨 처음에 올 수 있다.
\kfChoice 문법적 관계를 표시한다.
\end{kfChoices}
\end{kfQBlock}
\begin{kfQBlock}{2}{}
\kfQStem{밑줄 친 단어 중 품사가 \uline{다른} 하나는?}
\begin{kfChoices}
\kfChoice 하늘\uline{이} 매우 푸르다.
\kfChoice 밥\uline{을} 먹는다.
\kfChoice 너\uline{만} 기다린다.
\kfChoice \uline{매우} 빨리 달린다.
\kfChoice 학교\uline{에} 간다.
\end{kfChoices}
\end{kfQBlock}
\begin{kfQBlock}{3}{}
\kfQStem{<보기>의 빈칸에 공통으로 들어갈 수 있는 조사는?}
\begin{kfEvidence}* 나\_\_ 학생이다. * 사과\_\_ 맛있다.\end{kfEvidence}
\begin{kfChoices}
\kfChoice 는
\kfChoice 를
\kfChoice 만
\kfChoice 와
\kfChoice 까지
\end{kfChoices}
\end{kfQBlock}
\begin{kfQBlock}{4}{}
\kfQStem{밑줄 친 조사의 종류가 나머지와 \uline{다른} 하나는?}
\begin{kfChoices}
\kfChoice 너\uline{와} 나
\kfChoice 사과\uline{랑} 배
\kfChoice 밥\uline{과} 빵
\kfChoice 동생\uline{이} 웃는다.
\kfChoice 철수\uline{하고} 영희
\end{kfChoices}
\end{kfQBlock}
\begin{kfQBlock}{5}{}
\kfQStem{'특별한 의미'를 더해주는 보조사가 사용된 문장은?}
\begin{kfChoices}
\kfChoice 나는 너만 좋아한다.
\kfChoice 밥을 먹는다.
\kfChoice 학교에 간다.
\kfChoice 너와 내가 친구다.
\kfChoice 하늘이 맑다.
\end{kfChoices}
\end{kfQBlock}
\begin{kfQBlock}{6}{}
\kfQStem{‘동생이 빵을 먹는다’에 대한 설명으로 옳은 것은?}
\begin{kfChoices}
\kfChoice 조사 '이'는 동작의 대상임을 나타낸다.
\kfChoice 조사 '을'은 문장의 주인공임을 나타낸다.
\kfChoice 조사 '이'와 '을'은 모두 보조사이다.
\kfChoice '이'와 '을'은 문장에서의 자격을 나타낸다.
\kfChoice 조사 '이'와 '을'은 홀로 쓰일 수 있다.
\end{kfChoices}
\end{kfQBlock}
\begin{kfQBlock}{7}{}
\kfQStem{'두 단어를 이어주는' 접속 조사가 쓰인 문장은?}
\begin{kfChoices}
\kfChoice 그는 천재이다.
\kfChoice 밥을 먹어라.
\kfChoice 사과와 배를 샀다.
\kfChoice 비가 온다.
\kfChoice 너는 학생이다.
\end{kfChoices}
\end{kfQBlock}
\begin{kfQBlock}{8}{}
\kfQStem{<보기>의 빈칸에 들어갈 격조사가 바르게 짝지어진 것은?}
\begin{kfEvidence}* 철수\_\_ 학교에 간다. * 나는 사과\_\_ 먹는다.\end{kfEvidence}
\begin{kfChoices}
\kfChoice 가, 를
\kfChoice 는, 도
\kfChoice 가, 만
\kfChoice 는, 를
\kfChoice 와, 를
\end{kfChoices}
\end{kfQBlock}
\begin{kfQBlock}{1}{}
\kfQStem{명사, 대명사, 수사 뒤에 붙어 문법적 관계를 나타내는 품사는 무엇인가?}
\kfAnswerLines{1}
\end{kfQBlock}
\begin{kfQBlock}{2}{}
\kfQStem{문장에서의 문법적 자격을 나타내는 조사는 무엇인가?}
\kfAnswerLines{1}
\end{kfQBlock}
\begin{kfQBlock}{3}{}
\kfQStem{'은', '는', '만', '도'처럼 특별한 의미를 더하는 조사는 무엇인가?}
\kfAnswerLines{1}
\end{kfQBlock}
\begin{kfQBlock}{4}{}
\kfQStem{'와', '과'처럼 두 단어를 이어주는 조사는 무엇인가?}
\kfAnswerLines{1}
\end{kfQBlock}
\begin{kfQBlock}{5}{}
\kfQStem{‘동생이 밥을 먹는다’에서 주인공임을 나타내는 격조사는 무엇인가?}
\kfAnswerLines{1}
\end{kfQBlock}
\begin{kfQBlock}{6}{}
\kfQStem{‘동생이 밥을 먹는다’에서 행동의 대상임을 나타내는 격조사는 무엇인가?}
\kfAnswerLines{1}
\end{kfQBlock}
\begin{kfQBlock}{7}{}
\kfQStem{'이것만 너에게 줄게'에서 특별한 의미를 더하는 보조사는 무엇인가?}
\kfAnswerLines{1}
\end{kfQBlock}
\begin{kfQBlock}{8}{}
\kfQStem{'사과는 맛있지만 배는 맛없다'에서 대조의 의미를 더하는 보조사는 무엇인가?}
\kfAnswerLines{1}
\end{kfQBlock}
\begin{kfQBlock}{9}{}
\kfQStem{'철수야, 이리 와 봐'에서 누군가를 부를 때 사용하는 격조사는 무엇인가?}
\kfAnswerLines{1}
\end{kfQBlock}
\begin{kfQBlock}{10}{}
\kfQStem{'우리는 학교에 간다'에서 장소를 나타내는 격조사는 무엇인가?}
\kfAnswerLines{1}
\end{kfQBlock}
\begin{kfQBlock}{1}{}
\kfQStem{다음 <보기>의 흩어진 단어들을 조사를 사용하여 자연스러운 한 문장으로 만들어 보시오.}
\begin{kfEvidence}나, 사과, 바나나, 좋아하다\end{kfEvidence}
\kfAnswerNote{답안}{4}
\end{kfQBlock}


\kfStartConcept{개념 영역 — 시의 머리와 꼬리가 만나는 수미상관}


\kfSecHeader{kfAreaConcept}{개념}{시의 머리와 꼬리가 만나는 수미상관}

수미상관은 시의 첫 부분과 마지막 부분을 같거나 비슷한 시구로 구성하는 표현 기법이다. '수(首)'는 머리, '미(尾)'는 꼬리, '상관(相關)'은 서로 관련이 있다는 뜻으로, 시의 머리와 꼬리가 서로 호응한다는 의미이다. 마치 액자가 그림을 감싸듯이 시의 처음과 끝이 하나로 연결되어 완결된 구조를 만들어낸다.

수미상관에는 두 가지 형태가 있다. 첫째는 완전한 반복으로, 처음과 끝이 글자 하나 바꾸지 않고 똑같이 나오는 것이다. 한용운의 「나룻배와 행인」에서 "나는 나룻배 / 당신은 행인"이라는 구절이 처음과 끝에 그대로 반복되는 것이 그 예이다. 둘째는 변주된 반복으로, 기본 구조는 비슷하지만 일부 단어가 바뀌는 것이다. 김소월의 「산유화」에서 "산에는 꽃 피네"가 "산에는 꽃 지네"로 변주되어 꽃이 피고 지는 자연의 순환을 보여주는 것이 그 예이다.

수미상관은 여러 가지 효과를 준다. 먼저 시 전체에 안정감과 통일성을 부여하고, 시구를 반복하여 운율을 형성한다. 또한 시인이 가장 강조하고 싶은 주제나 감정을 처음과 끝에서 다시 한번 보여줌으로써 독자의 기억에 깊이 남기게 한다. 마지막으로 시를 다 읽고 나서 다시 처음으로 돌아가게 만들어 감동을 심화시키고 여운을 남기는 역할을 한다.
\par\medskip\begin{itemize}[leftmargin=18pt, itemsep=2pt, topsep=2pt, label=\textbullet]
\item 완전 반복 — '나는 나룻배 / 당신은 행인'(「나룻배와 행인」 처음·끝)
\item 변주 반복 — '산에는 꽃 피네' → '산에는 꽃 지네'(「산유화」)
\item 효과 — 통일감·운율 형성 + 주제 강조 + 여운 남김
\end{itemize}

\kfSecHeader{kfAreaConcept}{문제}{}

\begin{kfQBlock}{1}{}
\kfQStem{‘수미상관’에 대한 설명으로 적절하지 \uline{않은} 것은 무엇인가?}
\begin{kfChoices}
\kfChoice 비슷한 문장 구조를 짝지어 운율을 만드는 ‘대구법’과 같은 개념이다.
\kfChoice 시의 형식에 통일성과 안정감을 부여한다.
\kfChoice 핵심 정서나 주제를 한 번 더 강조하는 효과가 있다.
\kfChoice 시의 시작과 끝이 호응하여 깊은 여운을 남긴다.
\kfChoice 시의 첫 부분과 마지막 부분을 같거나 비슷한 시구절로 구성한다.
\end{kfChoices}
\end{kfQBlock}
\begin{kfQBlock}{2}{}
\kfQStem{‘수미상관(首尾相關)’의 한자 뜻풀이에 대한 설명으로 적절하지 \uline{않은} 것은 무엇인가?}
\begin{kfChoices}
\kfChoice ‘首’는 시의 머리(처음)를 가리킨다.
\kfChoice 머리와 꼬리가 서로 반대되는 의미로 대립한다는 뜻이다.
\kfChoice ‘相關’은 서로 관련을 맺고 호응한다는 뜻이다.
\kfChoice 머리와 꼬리가 서로 관련을 맺도록 구성한 시 형식을 일컫는다.
\kfChoice ‘尾’는 시의 꼬리(끝)를 가리킨다.
\end{kfChoices}
\end{kfQBlock}
\begin{kfQBlock}{3}{}
\kfQStem{수미상관이 시의 형식에 주는 효과로 적절하지 \uline{않은} 것은 무엇인가?}
\begin{kfChoices}
\kfChoice 시의 길이를 늘려 분량을 키우는 효과를 낸다.
\kfChoice 시작과 끝의 호응으로 형식미를 완성한다.
\kfChoice 시구의 반복으로 운율을 만든다.
\kfChoice 닫힌 구조를 형성하여 시의 완결성을 높인다.
\kfChoice 시 전체에 안정감과 통일성을 준다.
\end{kfChoices}
\end{kfQBlock}
\begin{kfQBlock}{4}{}
\kfQStem{수미상관이 운율을 형성하는 이유는 무엇인가?}
\begin{kfChoices}
\kfChoice 같거나 비슷한 시구를 반복하기 때문에
\kfChoice 마음을 울리는 슬픈 내용을 담고 있기 때문에
\kfChoice 처음부터 끝까지 시의 길이가 길어지기 때문에
\kfChoice 그림을 그리듯 장면을 자세히 묘사하기 때문에
\kfChoice 평소 잘 안 쓰는 어려운 단어를 사용하기 때문에
\end{kfChoices}
\end{kfQBlock}
\begin{kfQBlock}{5}{}
\kfQStem{수미상관이 시의 내용에 주는 효과로 적절하지 \uline{않은} 것은 무엇인가?}
\begin{kfChoices}
\kfChoice 시인이 강조하고 싶은 주제를 독자의 기억에 깊이 남긴다.
\kfChoice 시의 주제와 정반대되는 내용을 제시하여 반전의 재미를 준다.
\kfChoice 시의 주제를 처음과 끝에서 호응시켜 분명히 드러낸다.
\kfChoice 처음에 던진 의문이나 정서를 끝에서 한 번 더 마무리한다.
\kfChoice 핵심 정서를 한 번 더 환기시켜 강한 인상을 남긴다.
\end{kfChoices}
\end{kfQBlock}
\begin{kfQBlock}{6}{}
\kfQStem{수미상관이 시를 다 읽고 난 후 독자에게 주는 감동과 관련된 효과는 무엇인가?}
\begin{kfChoices}
\kfChoice 감동을 줄여주고 이성적으로 생각하게 한다.
\kfChoice 시의 내용이 모두 사실이라고 믿게 한다.
\kfChoice 다음 내용을 궁금하게 만들어 긴장감을 준다.
\kfChoice 감동을 더 깊게 하고 긴 여운을 남긴다.
\kfChoice 숨겨진 반전의 내용을 암시한다.
\end{kfChoices}
\end{kfQBlock}
\begin{kfQBlock}{7}{}
\kfQStem{다음 중 수미상관의 효과로 알맞지 \uline{않은} 것은 무엇인가?}
\begin{kfChoices}
\kfChoice 시 전체의 구조적 안정감을 높여 준다.
\kfChoice 시인이 전하려는 주제를 한층 강조해 준다.
\kfChoice 시 전체에 자연스러운 운율을 형성한다.
\kfChoice 시의 의미를 반대로 바꾸어 표현한다.
\kfChoice 다 읽은 독자에게 깊은 여운을 남겨 준다.
\end{kfChoices}
\end{kfQBlock}
\begin{kfQBlock}{8}{}
\kfQStem{다음 중 수미상관에 대한 설명으로 알맞지 \uline{않은} 것은 무엇인가?}
\begin{kfChoices}
\kfChoice 시의 머리와 꼬리가 서로 관련된다는 뜻이다.
\kfChoice 시의 처음과 중간 부분을 똑같이 만드는 기법이다.
\kfChoice 완전한 반복과 변주된 반복의 두 가지 형태가 있다.
\kfChoice 시 전체에 통일성을 부여하는 효과가 있다.
\kfChoice 감동을 심화시키고 여운을 남기는 역할을 한다.
\end{kfChoices}
\end{kfQBlock}
\begin{kfQBlock}{1}{}
\kfQStem{시의 첫 부분과 마지막 부분을 같거나 비슷하게 구성하는 표현 기법을 무엇이라고 하는가?}
\kfAnswerLines{1}
\end{kfQBlock}
\begin{kfQBlock}{2}{}
\kfQStem{'수미상관'에서 '수(首)'는 머리를, '미(尾)'는 무엇을 뜻하는가?}
\kfAnswerLines{1}
\end{kfQBlock}
\begin{kfQBlock}{3}{}
\kfQStem{수미상관은 시 전체에 통일성과 무엇을 부여하는가?}
\kfAnswerLines{1}
\end{kfQBlock}
\begin{kfQBlock}{4}{}
\kfQStem{시의 처음과 끝에 비슷한 구절을 반복하면 무엇이 형성되는가?}
\kfAnswerLines{1}
\end{kfQBlock}
\begin{kfQBlock}{5}{}
\kfQStem{시의 처음과 끝에서 같은 말을 반복하면 시인이 말하려는 무엇이 강조되는가?}
\kfAnswerLines{1}
\end{kfQBlock}
\begin{kfQBlock}{6}{}
\kfQStem{시를 다 읽고 나서 다시 처음으로 돌아가는 듯한 느낌을 주어, 감동을 더 깊게 하고 무엇을 남기는 역할을 하는가?}
\kfAnswerLines{1}
\end{kfQBlock}
\begin{kfQBlock}{7}{}
\kfQStem{김소월의 「산유화」에서 '꽃 피네'가 '꽃 지네'로 살짝 바뀌어 나타난 것은 자연의 어떤 원리를 보여주기 위함인가?}
\kfAnswerLines{1}
\end{kfQBlock}
\begin{kfQBlock}{1}{}
\kfQStem{시인들이 시의 처음과 끝을 비슷하게 만드는 수미상관 기법을 사용하는 이유는 무엇인지, 그 효과와 관련하여 <조건>에 맞게 서술하시오.}
\kfAnswerNote{답안}{4}
\end{kfQBlock}


\kfStartLit{문학 영역 — 작품}


\kfSecHeader{kfAreaLiterature}{지문}{작품}
\kfPassageNote{작품}{%
그리운 그의 얼굴 다시 찾을 수 없어도 \\[4pt]
화사한 그의 꽃 \\[4pt]
산에 언덕에 피어날지어이. \\[14pt]
그리운 그의 노래 다시 들을 수 없어도 \\[4pt]
맑은 그 숨결 \\[4pt]
들에 숲속에 살아갈지어이. \\[14pt]
쓸쓸한 마음으로 들길 더듬는 행인아. \\[14pt]
눈길 비었거든 바람 담을지네. \\[4pt]
바람 비었거든 인정 담을지네. \\[14pt]
그리운 그의 모습 다시 찾을 수 없어도 \\[4pt]
울고 간 그의 영혼 \\[4pt]
들에 언덕에 피어날지어이.
}


\kfSecHeader{kfAreaLiterature}{활동}{문학 문장 독해}
\noindent\textit{\small 작품 속 문장을 자세히 읽고, 문장별로 주어진 물음에 답해 보자.}\par\medskip
\begin{kfSentenceUnit}{1}{"그리운 그의 얼굴 다시 찾을 수 없어도 화사한 그의 꽃 산에 언덕에 피어날지어이."}
\kfSentQ{1.}{다음 문장에서 서술어를 기준으로 각각의 주어는 무엇인가? (1) 서술어 '찾을 수 없어도'의 주어는 무엇인가?}
\begin{kfChoices}
\kfChoice 내가
\kfChoice 꽃이
\kfChoice 바람이
\kfChoice 노래가
\end{kfChoices}
\kfSentQ{1.}{(2) 서술어 '피어날지어이'의 주어는 무엇인가?}
\begin{kfChoices}
\kfChoice 그의 얼굴이
\kfChoice 그의 꽃이
\kfChoice 산이
\kfChoice 언덕이
\end{kfChoices}
\end{kfSentenceUnit}

\begin{kfSentenceUnit}{2}{}
\kfSentQ{2.}{다음 예스러운 표현들의 현대어 풀이로 알맞은 것은? (1) "피어날지어이"}
\begin{kfChoices}
\kfChoice 피어나라
\kfChoice 피어날 것이다
\kfChoice 피어나고 싶다
\kfChoice 피어날까
\end{kfChoices}
\kfSentQ{2.}{(2) "살아갈지어이"}
\begin{kfChoices}
\kfChoice 살아갈 것이다
\kfChoice 살아가라
\kfChoice 살아가고 싶다
\kfChoice 살아갈까
\end{kfChoices}
\kfSentQ{2.}{(3) "담을지네"}
\begin{kfChoices}
\kfChoice 담아라
\kfChoice 담을 것이다
\kfChoice 담고 싶다
\kfChoice 담을까
\end{kfChoices}
\end{kfSentenceUnit}

\kfSecHeader{kfAreaLiterature}{활동}{해설 학습}
\noindent\textit{\small 작품 해설을 단원별로 읽고, 각 단원에 딸린 물음에 답해 보자.}\par\medskip
\par\noindent\textbf{\color{kfPrimary} 해설 단원 1}\par\smallskip
\kfPassageNote{해설 1}{%
「산에 언덕에」는 신동엽 시인이 1960년 4·19 혁명 때 희생된 젊은 학생들을 추모하며 쓴 시이다. 4·19 혁명은 이승만 정권의 독재에 맞서 학생들이 일으킨 민주화 운동으로, 많은 젊은이들이 목숨을 잃었다.

이 시는 신동엽이 4·19 혁명 때 희생된 학생들을 그리워하며 쓴 추모시이다. '그'는 나라를 위해 목숨을 바친 젊은 영웅들을 가리킨다.

화자는 '그'가 죽어서 더 이상 얼굴을 볼 수 없고 노래를 들을 수 없다는 슬픈 현실을 인정한다. 하지만 절망하지 않고, 그의 정신이 아름다운 '꽃'이 되어 산과 언덕에 피어나고, '숨결'이 되어 들과 숲에서 영원히 살아갈 것이라고 믿는다.

3연과 4연에서는 똑같이 슬퍼하는 '행인'에게 말을 걸어 위로한다. 눈으로 '그'를 찾을 수 없으면 자유로운 정신인 '바람'을, 그것도 없으면 따뜻한 마음인 '인정'을 담으라고 조언한다.

마지막 연에서는 '그'의 영혼이 반드시 이 땅에 다시 피어날 것이라는 확신을 보여준다. 죽음으로 끝나는 것이 아니라 새로운 생명으로 부활할 것이라는 희망적인 메시지를 전한다.
}

\begin{kfQBlock}{1}{문제}
\kfQStem{이 시는 신동엽 시인이 어떤 역사적 사건 때 희생된 학생들을 추모하며 쓴 시인가?}
\kfAnswerLines{1}
\end{kfQBlock}

\begin{kfQBlock}{2}{문제}
\kfQStem{시에서 화자가 그리워하는 '그'는 누구를 가리키는가?}
\kfAnswerLines{1}
\end{kfQBlock}

\begin{kfQBlock}{3}{문제}
\kfQStem{화자는 '그'가 죽었지만, 그의 정신이 무엇이 되어 산과 언덕에 피어날 것이라고 믿는가?}
\kfAnswerLines{1}
\end{kfQBlock}

\begin{kfQBlock}{4}{문제}
\kfQStem{'그'의 정신은 '꽃'과 함께 무엇이 되어 들과 숲에서 영원히 살아갈 것이라고 했는가?}
\kfAnswerLines{1}
\end{kfQBlock}

\begin{kfQBlock}{5}{문제}
\kfQStem{3연과 4연에서 화자는 자신과 똑같이 슬퍼하는 누구에게 말을 걸어 위로하는가?}
\kfAnswerLines{1}
\end{kfQBlock}

\begin{kfQBlock}{6}{문제}
\kfQStem{화자는 행인에게, 눈으로 '그'를 찾을 수 없다면 무엇을 담으라고 조언하는가?}
\kfAnswerLines{1}
\end{kfQBlock}

\begin{kfQBlock}{7}{문제}
\kfQStem{화자는 행인에게 '바람'도 느낄 수 없다면, 무엇을 담으라고 조언하는가?}
\kfAnswerLines{1}
\end{kfQBlock}

\par\noindent\textbf{\color{kfPrimary} 해설 단원 2}\par\smallskip
\kfPassageNote{해설 2}{%
이 시는 총 5연으로 구성되어 있으며, 각 연마다 조금씩 다른 내용을 담고 있다. 1연에서는 '그'의 얼굴을 다시 볼 수 없다는 현실을 인정하면서도, 그가 아름다운 꽃으로 피어나기를 바라는 마음을 표현한다.

2연에서는 '그'의 노래를 더 이상 들을 수 없지만, 그의 맑은 정신이 자연 속에서 계속 살아갈 것이라고 노래한다. 여기서 '노래'는 그가 가졌던 신념이나 꿈을 의미하고, '숨결'은 보이지 않지만 여전히 살아있는 그의 정신을 나타낸다.

3연과 4연에서는 '행인'이라는 새로운 인물이 등장한다. 행인은 '그'의 부재를 슬퍼하며 그의 흔적을 찾아 헤매는 사람이다. 화자는 이런 행인에게 말을 건네며, 눈으로 볼 수 없다면 바람을, 바람도 느낄 수 없다면 따뜻한 인정을 마음에 담으라고 위로한다.

마지막 5연에서는 다시 1연과 비슷한 구조로 돌아가면서, '그'의 영혼이 반드시 이 땅에 피어날 것이라는 확신을 표현한다. 전체적으로 이 시는 죽음이라는 슬픈 현실을 부정하지 않으면서도, 그 속에서 희망과 생명의 의미를 찾아내는 내용을 담고 있다.
}

\begin{kfQBlock}{1}{문제}
\kfQStem{이 시는 총 몇 개의 연으로 구성되어 있는가?}
\kfAnswerLines{1}
\end{kfQBlock}

\begin{kfQBlock}{2}{문제}
\kfQStem{1연에서 화자는 '그'가 무엇으로 피어나기를 바라고 있는가?}
\kfAnswerLines{1}
\end{kfQBlock}

\begin{kfQBlock}{3}{문제}
\kfQStem{2연에서 '그'가 가졌던 신념이나 꿈을 의미하는 시어는 무엇인가?}
\kfAnswerLines{1}
\end{kfQBlock}

\begin{kfQBlock}{4}{문제}
\kfQStem{'그'의 보이지 않지만 여전히 살아있는 정신을 나타내는 시어는 무엇인가?}
\kfAnswerLines{1}
\end{kfQBlock}

\begin{kfQBlock}{5}{문제}
\kfQStem{'그'가 없다는 사실, 즉 '그'의 부재를 슬퍼하며 헤매는 사람은 누구인가?}
\kfAnswerLines{1}
\end{kfQBlock}

\begin{kfQBlock}{6}{문제}
\kfQStem{화자는 행인에게, 눈으로 볼 수 없다면 바람을, 바람도 느낄 수 없다면 무엇을 마음에 담으라고 위로하는가?}
\kfAnswerLines{1}
\end{kfQBlock}

\begin{kfQBlock}{7}{문제}
\kfQStem{마지막 5연에서 화자는 '그'의 무엇이 반드시 피어날 것이라고 확신하는가?}
\kfAnswerLines{1}
\end{kfQBlock}

\par\noindent\textbf{\color{kfPrimary} 해설 단원 3}\par\smallskip
\kfPassageNote{해설 3}{%
이 시의 가장 큰 특징은 '반복법'이다. "-ㄹ지어이"라는 특별한 어미가 1연, 2연, 5연에서 반복되어 나타난다. 이는 예스러운 표현으로, 단순한 바람이 아니라 '반드시 그렇게 될 것이다'라는 강한 확신을 담고 있다. 이런 반복을 통해 시 전체에 운율감이 생기고, 화자의 간절한 마음이 더욱 강조된다.

'수미상관'이라는 표현 기법도 중요하다. 수미상관이란 첫 부분과 끝 부분이 비슷한 구조를 가지는 것을 말한다. 1연과 5연이 "그리운 그의 \textasciitilde{} 다시 찾을 수 없어도"로 시작하여 비슷한 패턴을 보인다. 이를 통해 시에 안정감을 주면서도 주제 의식을 더욱 강조하게 된다.

또한 '대구법'이 사용되었다. 대구법이란 비슷한 구조의 문장이나 구를 나란히 배치하여 리듬감을 만들고 의미를 강조하는 표현 방법이다. "눈길 비었거든 바람 담을지네. 바람 비었거든 인정 담을지네."에서 보듯이 "\textasciitilde{} 비었거든 \textasciitilde{} 담을지네"라는 똑같은 문장 구조를 반복하여 리듬감을 만들고 의미를 강조했다. 이런 대구법을 통해 화자가 행인에게 전하고 싶은 메시지가 더욱 또렷하게 전달된다.
}

\begin{kfQBlock}{1}{문제}
\kfQStem{"-ㄹ지어이"라는 어미를 계속 사용하는 표현 기법은 무엇인가?}
\kfAnswerLines{1}
\end{kfQBlock}

\begin{kfQBlock}{2}{문제}
\kfQStem{시의 첫 부분과 끝 부분이 비슷한 구조를 가지는 표현 기법을 무엇이라고 하는가?}
\kfAnswerLines{1}
\end{kfQBlock}

\begin{kfQBlock}{3}{문제}
\kfQStem{이 시에서는 몇 연과 몇 연이 수미상관 구조를 이루고 있는가?}
\kfAnswerLines{1}
\end{kfQBlock}

\begin{kfQBlock}{4}{문제}
\kfQStem{수미상관 기법은 시에 안정감을 주고 무엇을 강조하는 효과가 있는가?}
\kfAnswerLines{1}
\end{kfQBlock}

\begin{kfQBlock}{5}{문제}
\kfQStem{비슷한 구조의 문장이나 구를 나란히 배치하여 리듬감을 만드는 표현 방법을 무엇이라고 하는가?}
\kfAnswerLines{1}
\end{kfQBlock}

\begin{kfQBlock}{6}{문제}
\kfQStem{이 시에서 "\textasciitilde{}\textasciitilde{} 비었거든 \textasciitilde{}\textasciitilde{} 담을지네"라는 똑같은 문장 구조를 반복한 표현법은 무엇인가?}
\kfAnswerLines{1}
\end{kfQBlock}

\begin{kfQBlock}{7}{문제}
\kfQStem{대구법은 리듬감을 만들고 무엇을 강조하는 효과가 있는가?}
\kfAnswerLines{1}
\end{kfQBlock}

\par\noindent\textbf{\color{kfPrimary} 해설 단원 4}\par\smallskip
\kfPassageNote{해설 4}{%
이 시에 나오는 여러 시어들은 각각 깊은 의미를 담고 있다. '그'는 화자가 그리워하는 대상으로, 4·19 혁명에서 희생된 젊은 학생들을 상징한다. 단순히 한 개인이 아니라, 정의를 위해 목숨을 바친 모든 젊은 영혼들을 대표한다.

'꽃'은 죽음을 이긴 새로운 생명을 상징한다. 꽃은 아름답고 화사하며, 생명력이 넘치는 존재이다. 화자는 '그'의 정신이 꽃처럼 아름답게 다시 피어날 것이라고 믿고 있다. '숨결'은 보이지는 않지만 살아있는 것을 의미한다. 사람이 숨을 쉬어야 살 수 있듯이, '그'의 정신도 보이지 않는 숨결처럼 이 땅에서 계속 살아갈 것이라는 뜻이다.

'행인'은 길을 가다가 우연히 만나는 사람을 뜻하지만, 여기서는 '그'를 그리워하는 모든 사람들을 나타낸다. 화자와 마찬가지로 '그'의 부재를 슬퍼하며 그의 흔적을 찾아 헤매는 존재이다. '바람'은 자유롭고 시원한 것으로, 자유로운 정신을 상징한다. '인정'은 사람과 사람 사이의 따뜻한 마음을 의미한다.

'산', '언덕', '들', '숲'과 같은 자연 공간들은 '그'의 정신이 부활할 무대를 나타낸다. 이런 넓고 푸른 자연 속에서 '그'의 영혼이 영원히 살아갈 것이라는 희망을 담고 있다.
}

\begin{kfQBlock}{1}{문제}
\kfQStem{이 시에서 화자가 그리워하는 '그'는 누구를 상징하는가?}
\kfAnswerLines{1}
\end{kfQBlock}

\begin{kfQBlock}{2}{문제}
\kfQStem{죽음을 이긴 새로운 생명을 상징하며, '그'의 정신이 아름답게 부활한 모습을 나타내는 시어는 무엇인가?}
\kfAnswerLines{1}
\end{kfQBlock}

\begin{kfQBlock}{3}{문제}
\kfQStem{보이지는 않지만 살아있는 '그'의 정신을 의미하는 시어는 무엇인가?}
\kfAnswerLines{1}
\end{kfQBlock}

\begin{kfQBlock}{4}{문제}
\kfQStem{'그'를 그리워하며 그의 흔적을 찾아 헤매는 모든 사람들을 상징하는 시어는 무엇인가?}
\kfAnswerLines{1}
\end{kfQBlock}

\begin{kfQBlock}{5}{문제}
\kfQStem{'그'의 자유로운 정신을 상징하는 시어는 무엇인가?}
\kfAnswerLines{1}
\end{kfQBlock}

\begin{kfQBlock}{6}{문제}
\kfQStem{사람과 사람 사이의 따뜻한 마음을 의미하는 시어는 무엇인가?}
\kfAnswerLines{1}
\end{kfQBlock}

\begin{kfQBlock}{7}{문제}
\kfQStem{'산', '언덕', '들', '숲'과 같은 자연 공간은 무엇을 의미하는가?}
\kfAnswerLines{1}
\end{kfQBlock}

\kfSecHeader{kfAreaLiterature}{문제}{}

\begin{kfQBlock}{1}{}
\kfQStem{‘그’의 죽음에 대한 말하는 이의 태도로 적절하지 \uline{않은} 것은 무엇인가?}
\begin{kfChoices}
\kfChoice 그의 정신과 영혼이 꽃·숨결로 다시 피어날 것이라고 믿는다.
\kfChoice 죽음을 헛된 것으로 여겨 슬픔과 절망에 잠겨 있다.
\kfChoice 그의 희생을 추모하면서도 그 의미를 영원히 살아 있는 것으로 받아들인다.
\kfChoice 죽음을 통해 자연과 후손에게 이어지는 생명력을 확신한다.
\kfChoice 죽음을 절망으로만 보지 않고 부활의 가능성을 노래한다.
\end{kfChoices}
\end{kfQBlock}
\begin{kfQBlock}{2}{}
\kfQStem{이 시에서 ‘꽃’과 ‘숨결’이 공통적으로 상징하는 의미로 적절하지 \uline{않은} 것은 무엇인가?}
\begin{kfChoices}
\kfChoice ‘그’를 죽음으로 이끈 비극적 원인
\kfChoice 사라진 듯 보이지만 다시 살아나는 생명의 부활
\kfChoice ‘그’의 숭고한 정신이 자연 속에 영원히 깃들어 있음
\kfChoice 슬픔을 넘어 다시 피어나는 새로운 생명력
\kfChoice 죽은 ‘그’의 정신이 다시 태어난 모습
\end{kfChoices}
\end{kfQBlock}
\begin{kfQBlock}{3}{}
\kfQStem{3연에 나오는 ‘행인’은 어떤 사람이라고 할 수 있는가?}
\begin{kfChoices}
\kfChoice '그'의 죽음에 아무 관심이 없는 사람
\kfChoice '그'를 죽음으로 몰아간 나쁜 사람
\kfChoice 말하는 이와 마찬가지로 '그'를 그리워하는 사람
\kfChoice '그'의 원수를 갚기 위해 길을 떠나는 사람
\kfChoice '그'의 무덤을 찾아와 슬프게 우는 가족
\end{kfChoices}
\end{kfQBlock}
\begin{kfQBlock}{4}{}
\kfQStem{‘\textasciitilde{}ㄹ지어이’라는 말의 반복이 만들어 내는 효과로 적절하지 \uline{않은} 것은 무엇인가?}
\begin{kfChoices}
\kfChoice 화자가 자신의 주장이 틀렸을 수 있다는 겸손함을 강조한다.
\kfChoice 같은 종결 어미의 반복으로 운율을 만들어 낸다.
\kfChoice 단정적이고 권유적인 말투로 화자의 확신을 드러낸다.
\kfChoice 핵심 정서를 청자에게 강하게 각인시키는 효과를 낸다.
\kfChoice ‘그’가 꼭 다시 부활할 것이라는 간절한 소망과 강한 믿음을 강조한다.
\end{kfChoices}
\end{kfQBlock}
\begin{kfQBlock}{5}{}
\kfQStem{이 시의 수미상관 구조가 만들어 내는 효과로 적절하지 \uline{않은} 것은 무엇인가?}
\begin{kfChoices}
\kfChoice 시 전체에 안정감을 주고 주제를 한 번 더 강조한다.
\kfChoice 말하는 이의 마음이 약해졌음을 암시하는 신호 역할을 한다.
\kfChoice 처음과 끝의 호응으로 형식적 완결성을 갖추게 한다.
\kfChoice 핵심 정서를 반복해 깊은 여운을 남긴다.
\kfChoice 부활에 대한 화자의 굳은 믿음을 강하게 부각한다.
\end{kfChoices}
\end{kfQBlock}
\begin{kfQBlock}{6}{}
\kfQStem{이 시의 ‘그’가 4·19 혁명 때 희생된 젊은 학생들을 상징한다고 할 때, 이 시의 주제로 적절하지 \uline{않은} 것은 무엇인가?}
\begin{kfChoices}
\kfChoice 자연의 아름다움과 생명력만을 예찬하는 단순한 자연시
\kfChoice 희생된 이들의 정신이 ‘꽃·숨결’로 부활한다는 굳은 믿음
\kfChoice 비극적 죽음을 헛된 것으로 보지 않고 의미 있는 가치로 승화하려는 의지
\kfChoice 4·19 정신이 후대에도 이어져 살아 있음을 노래하는 추모와 다짐
\kfChoice 민주주의를 위해 희생된 숭고한 정신에 대한 추모
\end{kfChoices}
\end{kfQBlock}
\begin{kfQBlock}{7}{}
\kfQStem{"눈길 비었거든 바람 담을지네. / 바람 비었거든 인정 담을지네."와 같이 비슷한 문장 구조를 짝지어 운율을 만드는 표현법은 무엇인가?}
\begin{kfChoices}
\kfChoice 과장법
\kfChoice 반어법
\kfChoice 역설법
\kfChoice 대구법
\kfChoice 의인법
\end{kfChoices}
\end{kfQBlock}
\begin{kfQBlock}{8}{}
\kfQStem{이 시에 대한 설명으로 가장 알맞은 것은 무엇인가?}
\begin{kfChoices}
\kfChoice '그'의 죽음이라는 슬픈 현실에 절망하며 무거운 분위기로 끝을 맺는다.
\kfChoice '그'가 죽기 전에 마지막으로 남긴 유언의 내용을 자세히 전달하고 있다.
\kfChoice '그'의 죽음의 원인을 밝혀내고 그에 대한 복수를 굳게 다짐하고 있다.
\kfChoice '그'의 죽음을 헛되지 않은, 영원한 생명으로 여기는 희망적인 태도를 보인다.
\kfChoice '그'와 함께했던 어린 시절의 즐거웠던 일들을 한 가지씩 자세히 회상하고 있다.
\end{kfChoices}
\end{kfQBlock}
\begin{kfQBlock}{1}{}
\kfQStem{이 시는 “그리운 그의 \textasciitilde{} 없어도 / \textasciitilde{}ㄹ지어이”라는 비슷한 문장 구조를 반복합니다. 이러한 표현이 주는 효과를 서술하시오.}
\kfAnswerNote{답안}{4}
\end{kfQBlock}
\begin{kfQBlock}{2}{}
\kfQStem{이 시에서 말하는 이의 소망이 어떻게 점점 더 깊어지는지 서술하시오.}
\kfAnswerNote{답안}{4}
\end{kfQBlock}


\kfStartNonlit{비문학 영역 — [과학] 식물은 어떻게 높은 곳까지 물을 끌어올릴까?}


\kfSecHeader{kfAreaNonlit}{지문}{[과학] 식물은 어떻게 높은 곳까지 물을 끌어올릴까?}
\kfPassageNote{[과학] 식물은 어떻게 높은 곳까지 물을 끌어올릴까?}{%
식물이 살아가려면 물이 꼭 필요하다. 식물은 잎에서 햇빛을 이용해 스스로 양분을 만드는데, 이때 물은 아주 중요한 재료가 된다. 그런데 여기에는 큰 문제가 하나 있다. 바로 중력이다. 중력은 물을 항상 아래로 끌어당기지만 식물은 위로 자라기 때문에, 중력을 거슬러 물을 위로 올려야만 한다. 그렇다면 키가 큰 나무는 대체 어떻게 꼭대기에 있는 나뭇잎까지 물을 보내는 것일까?

세상에서 가장 키가 큰 ‘세쿼이아’ 나무는 그 높이가 아파트 40층과 맞먹는다. 이 거대한 나무가 펌프 하나 없이 꼭대기까지 물을 끌어올리는 비결은 무엇일까? 과학자들은 그 비밀이 바로 세 가지 힘의 조화에 있다고 말한다. 식물은 뿌리압, 모세관 현상, 증산 작용이라는 힘을 이용해서 이 어려운 일을 해낸다.

첫 번째 힘은 뿌리가 물을 밀어 올리는 ‘뿌리압’이다. 식물의 뿌리 안의 액체는 주변 흙에 있는 물보다 농도가 더 진하다. 그래서 흙 속의 물이 농도를 맞추기 위해 저절로 뿌리 안으로 들어온다. 이렇게 물이 계속 들어오면서 생기는 압력, 즉 뿌리압이 물줄기를 위로 조금 밀어 올린다.

두 번째 힘은 ‘모세관 현상’이다. 물이 담긴 컵에 가느다란 빨대를 넣어보면 빨대 안쪽으로 물이 조금 올라가는 것을 볼 수 있다. 이처럼 가는 관과 같은 통로를 따라 액체가 올라가거나 내려가는 것을 모세관 현상이라고 한다. 이는 물 분자가 빨대 벽에 달라붙으려는 힘이 물 분자끼리 붙으려는 힘보다 크기 때문이다. 관이 가늘수록 물이 더 높이 올라가는데, 식물 몸속에는 눈에 보이지 않을 만큼 아주 가느다란 '물관'이라는 통로가 있다. 물은 이 좁은 통로의 벽에 착 달라붙는 성질을 이용해 위로 올라간다.

세 번째 힘은 잎에서 물을 끌어당기는 ‘증산 작용’이다. 식물은 잎에 있는 작은 구멍을 통해 물을 수증기 형태로 내보낸다. 이처럼 식물체 안의 물이 잎의 기공을 통해 수증기 상태로 날아가는 현상을 증산 작용이라고 한다. 식물 속 물 분자들은 마치 사슬처럼 서로 단단히 이어져 있어서, 잎에서 물 분자 하나가 빠져나가면, 그 힘이 연결된 물줄기 전체를 위로 당겨준다. 이것이 바로 뿌리에 있는 물을 나무 꼭대기까지 끌어올리는 가장 강력한 힘이다.

이처럼 식물은 어느 한 가지 힘만 사용하는 것이 아니다. 뿌리에서 밀어주고, 줄기에서 타고 올라가고, 잎에서 강하게 끌어당기는 세 가지 힘을 조화롭게 이용하여, 아무리 키가 큰 나무라도 꼭대기까지 물을 보낼 수 있는 것이다.
}


\kfSecHeader{kfAreaNonlit}{활동}{어휘 학습}
\noindent\textit{\small 다음 표의 '의미'를 읽고, 그에 해당하는 '어휘'를 지문에서 찾아 적으시오.}\par\medskip
\begin{kfVocab}
\kfVocabItem{1}{식물이 살아가는 데 필요한 영양소}
\kfVocabItem{2}{어떤 것을 만드는 데 필요한 것}
\kfVocabItem{3}{모든 물체를 아래로 끌어당기는 힘}
\kfVocabItem{4}{액체를 낮은 곳에서 높은 곳으로 끌어올리는 기계}
\kfVocabItem{5}{남에게 알려지지 않은 특별한 방법}
\kfVocabItem{6}{여러 가지가 서로 잘 어울림}
\kfVocabItem{7}{뿌리가 물을 빨아들여 위로 밀어 올리는 압력}
\kfVocabItem{8}{액체가 아주 가는 관을 따라 저절로 올라가는 현상}
\kfVocabItem{9}{식물이 잎을 통해 물을 수증기로 내보내는 작용}
\kfVocabItem{10}{액체의 진하고 묽은 정도}
\kfVocabItem{11}{물질의 성질을 가지는 가장 작은 알갱이}
\kfVocabItem{12}{식물 줄기 속에 있는, 물이 이동하는 가느다란 관}
\kfVocabItem{13}{물이 끓거나 증발하여 기체로 변한 것}
\kfVocabItem{14}{고리가 여러 개 연결된 모양의 물건}
\end{kfVocab}

\kfSecHeader{kfAreaNonlit}{활동}{비문학 문장 독해}
\noindent\textit{\small 지문을 문장 단위로 정독하면서, 문장별로 주어진 물음에 답해 보자.}\par\medskip
\begin{kfSentenceUnit}{1}{식물이 살아가려면 물이 꼭 필요하다.}
\kfSentQ{1.}{누가 살아가려면 물이 필요한가?}
\begin{kfChoices}
\kfChoice 사람
\kfChoice 동물
\kfChoice 식물
\kfChoice 바위
\end{kfChoices}
\kfSentQ{2.}{식물이 살아가기 위해 꼭 필요한 것은 무엇인가?}
\begin{kfChoices}
\kfChoice 흙
\kfChoice 물
\kfChoice 바람
\kfChoice 그늘
\end{kfChoices}
\end{kfSentenceUnit}

\begin{kfSentenceUnit}{2}{식물은 잎에서 햇빛을 이용해 스스로 양분을 만드는데, 이때 물은 아주 중요한 재료가 된다.}
\kfSentQ{3.}{'이때'가 가리키는 때는 언제인가?}
\begin{kfChoices}
\kfChoice 식물이 춤을 출 때
\kfChoice 식물이 잠을 잘 때
\kfChoice 식물이 양분을 만들 때
\kfChoice 식물이 물을 마실 때
\end{kfChoices}
\kfSentQ{4.}{식물은 어디서 양분을 만드는가?}
\begin{kfChoices}
\kfChoice 뿌리
\kfChoice 줄기
\kfChoice 잎
\kfChoice 꽃
\end{kfChoices}
\kfSentQ{5.}{식물이 양분을 만들 때 이용하는 것은 무엇인가?}
\begin{kfChoices}
\kfChoice 달빛
\kfChoice 별빛
\kfChoice 전등빛
\kfChoice 햇빛
\end{kfChoices}
\end{kfSentenceUnit}

\begin{kfSentenceUnit}{3}{그런데 여기에는 큰 문제가 하나 있다.}
\kfSentQ{6.}{'여기'가 가리키는 것은 무엇인가?}
\begin{kfChoices}
\kfChoice 식물이 잎에서 양분을 만드는 일
\kfChoice 식물이 뿌리를 내리는 일
\kfChoice 식물이 꽃을 피우는 일
\kfChoice 식물이 열매를 맺는 일
\end{kfChoices}
\kfSentQ{7.}{이 문장은 앞으로 어떤 내용이 나올 것을 알려주는가?}
\begin{kfChoices}
\kfChoice 식물이 겪는 문제점
\kfChoice 식물의 성장 속도
\kfChoice 식물의 다양한 종류
\kfChoice 식물이 사는 곳
\end{kfChoices}
\end{kfSentenceUnit}

\begin{kfSentenceUnit}{4}{바로 중력이다.}
\kfSentQ{8.}{식물이 물을 사용하는 데 겪는 큰 문제가 무엇인가?}
\begin{kfChoices}
\kfChoice 바람
\kfChoice 햇빛
\kfChoice 중력
\kfChoice 온도
\end{kfChoices}
\kfSentQ{9.}{식물이 물을 위로 올릴 때 극복해야 하는 힘은 무엇인가?}
\begin{kfChoices}
\kfChoice 마찰력
\kfChoice 탄성력
\kfChoice 중력
\kfChoice 자기력
\end{kfChoices}
\end{kfSentenceUnit}

\begin{kfSentenceUnit}{5}{중력은 물을 항상 아래로 끌어당기지만 식물은 위로 자라기 때문에, 중력을 거슬러 물을 위로 올려야만 한다.}
\kfSentQ{10.}{중력은 물을 어느 방향으로 끌어당기는가?}
\begin{kfChoices}
\kfChoice 위로
\kfChoice 아래로
\kfChoice 옆으로
\kfChoice 앞쪽으로
\end{kfChoices}
\kfSentQ{11.}{식물은 어느 방향으로 자라는가?}
\begin{kfChoices}
\kfChoice 위로
\kfChoice 아래로
\kfChoice 옆으로
\kfChoice 뒤로
\end{kfChoices}
\kfSentQ{12.}{식물이 해결해야 할 문제는 무엇인가?}
\begin{kfChoices}
\kfChoice 물을 아래로 내리는 것
\kfChoice 중력을 거슬러 물을 위로 올리는 것
\kfChoice 뜨거운 햇빛을 피해 그늘로 숨는 것
\kfChoice 강한 바람에 흔들리지 않고 중심을 잡는 것
\end{kfChoices}
\kfSentQ{13.}{식물이 중력을 거슬러야만 하는 이유는 무엇인가?}
\begin{kfChoices}
\kfChoice 식물은 물을 싫어하기 때문이다.
\kfChoice 식물은 아래로 자라기 때문이다.
\kfChoice 식물은 위로 자라기 때문이다.
\kfChoice 중력이 물을 위로 밀기 때문이다.
\end{kfChoices}
\end{kfSentenceUnit}

\begin{kfSentenceUnit}{6}{그렇다면 키가 큰 나무는 대체 어떻게 꼭대기에 있는 나뭇잎까지 물을 보내는 것일까?}
\kfSentQ{14.}{물을 보내야 하는 목적지는 어디인가?}
\begin{kfChoices}
\kfChoice 뿌리 끝
\kfChoice 줄기 중간
\kfChoice 꽃잎
\kfChoice 꼭대기에 있는 나뭇잎
\end{kfChoices}
\kfSentQ{15.}{이 질문을 통해 앞으로 어떤 내용이 나올 것을 예상할 수 있는가?}
\begin{kfChoices}
\kfChoice 나무가 자라는 속도
\kfChoice 나무가 물을 꼭대기까지 보내는 방법
\kfChoice 나무의 종류와 특징
\kfChoice 나무가 겪는 질병
\end{kfChoices}
\end{kfSentenceUnit}

\begin{kfSentenceUnit}{7}{세상에서 가장 키가 큰 ‘세쿼이아’ 나무는 그 높이가 아파트 40층과 맞먹는다.}
\kfSentQ{16.}{세상에서 가장 키가 큰 나무의 이름은 무엇인가?}
\begin{kfChoices}
\kfChoice 소나무
\kfChoice 참나무
\kfChoice 세쿼이아
\kfChoice 대나무
\end{kfChoices}
\kfSentQ{17.}{세쿼이아 나무의 높이를 표현하기 위해 무엇과 비교했는가?}
\begin{kfChoices}
\kfChoice 63빌딩
\kfChoice 에펠탑
\kfChoice 아파트 40층
\kfChoice 자유의 여신상
\end{kfChoices}
\end{kfSentenceUnit}

\begin{kfSentenceUnit}{8}{이 거대한 나무가 펌프 하나 없이 꼭대기까지 물을 끌어올리는 비결은 무엇일까?}
\kfSentQ{18.}{'이 거대한 나무'가 가리키는 것은 무엇인가?}
\begin{kfChoices}
\kfChoice 참나무
\kfChoice 세쿼이아 나무
\kfChoice 소나무
\kfChoice 전나무
\end{kfChoices}
\kfSentQ{19.}{이 질문을 통해 앞으로 어떤 내용이 나올 것을 예상할 수 있는가?}
\begin{kfChoices}
\kfChoice 나무가 펌프 없이도 물을 끌어올리는 방법
\kfChoice 나무가 펌프를 사용하는 방법
\kfChoice 나무가 물을 저장하는 방법
\kfChoice 나무가 잎을 떨어뜨리는 이유
\end{kfChoices}
\end{kfSentenceUnit}

\begin{kfSentenceUnit}{9}{과학자들은 그 비밀이 바로 세 가지 힘의 조화에 있다고 말한다.}
\kfSentQ{20.}{'그 비밀'이 가리키는 것은 무엇인가?}
\begin{kfChoices}
\kfChoice 잎에서 햇빛을 효율적으로 받아 양분을 만드는 방법
\kfChoice 뿌리에서 흡수한 물로 광합성을 하여 영양분을 만드는 방법
\kfChoice 나무가 펌프 하나 없이 꼭대기까지 물을 끌어올리는 것
\kfChoice 계절에 맞추어 아름다운 꽃을 피우는 시기
\end{kfChoices}
\kfSentQ{21.}{나무가 펌프 하나 없이 꼭대기까지 물을 끌어올리는 힘은 몇 가지인가?}
\begin{kfChoices}
\kfChoice 한 가지
\kfChoice 두 가지
\kfChoice 세 가지
\kfChoice 네 가지
\end{kfChoices}
\end{kfSentenceUnit}

\begin{kfSentenceUnit}{10}{식물은 뿌리압, 모세관 현상, 증산 작용이라는 힘을 이용해서 이 어려운 일을 해낸다.}
\kfSentQ{22.}{'이 어려운 일'이 가리키는 것은 무엇인가?}
\begin{kfChoices}
\kfChoice 높은 곳에 있는 물을 자연스럽게 아래로 흘려보내는 일
\kfChoice 중력을 거슬러 물을 위로 끌어올리는 일
\kfChoice 잎에서 광합성을 통해 필요한 양분을 스스로 만드는 일
\kfChoice 벌과 나비를 유인하기 위해 화려한 꽃을 피우는 일
\end{kfChoices}
\kfSentQ{23.}{식물이 물을 위로 올릴 때 사용하는 세 가지 힘은 무엇인가?}
\begin{kfChoices}
\kfChoice 중력, 마찰력, 탄성력
\kfChoice 뿌리압, 모세관 현상, 증산 작용
\kfChoice 전기력, 자기력, 마찰력
\kfChoice 원심력, 구심력, 중력
\end{kfChoices}
\end{kfSentenceUnit}

\begin{kfSentenceUnit}{11}{첫 번째 힘은 뿌리가 물을 밀어 올리는 ‘뿌리압’이다.}
\kfSentQ{24.}{식물이 물을 올리는 첫 번째 힘은 무엇인가?}
\begin{kfChoices}
\kfChoice 모세관 현상
\kfChoice 증산 작용
\kfChoice 중력
\kfChoice 뿌리압
\end{kfChoices}
\kfSentQ{25.}{'뿌리압'이란 무엇을 하는 힘인가?}
\begin{kfChoices}
\kfChoice 뿌리가 물을 밀어 올리는 힘
\kfChoice 잎이 물을 당기는 힘
\kfChoice 줄기가 물을 잡는 힘
\kfChoice 꽃이 물을 마시는 힘
\end{kfChoices}
\end{kfSentenceUnit}

\begin{kfSentenceUnit}{12}{식물의 뿌리 안의 액체는 주변 흙에 있는 물보다 농도가 더 진하다.}
\kfSentQ{26.}{뿌리 안의 액체와 흙 속 물의 농도는 어떻게 다른가?}
\begin{kfChoices}
\kfChoice 뿌리 안이 더 묽다.
\kfChoice 흙 속이 더 진하다.
\kfChoice 농도가 같다.
\kfChoice 뿌리 안이 더 진하다.
\end{kfChoices}
\end{kfSentenceUnit}

\begin{kfSentenceUnit}{13}{그래서 흙 속의 물이 농도를 맞추기 위해 저절로 뿌리 안으로 들어온다.}
\kfSentQ{27.}{'그래서'의 구체적인 내용은 무엇인가?}
\begin{kfChoices}
\kfChoice 식물이 물을 싫어해서
\kfChoice 뿌리 안의 액체 농도가 더 진해서
\kfChoice 흙 속에 물이 없어서
\kfChoice 햇빛이 너무 강해서
\end{kfChoices}
\kfSentQ{28.}{흙 속의 물이 뿌리 안으로 들어오는 이유는 무엇인가?}
\begin{kfChoices}
\kfChoice 농도를 맞추기 위해서
\kfChoice 온도를 맞추기 위해서
\kfChoice 부피를 줄이기 위해서
\kfChoice 무게를 늘리기 위해서
\end{kfChoices}
\end{kfSentenceUnit}

\begin{kfSentenceUnit}{14}{이렇게 물이 계속 들어오면서 생기는 압력, 즉 뿌리압이 물줄기를 위로 조금 밀어 올린다.}
\kfSentQ{29.}{'이렇게'가 가리키는 것은 무엇인가?}
\begin{kfChoices}
\kfChoice 뜨거운 태양열에 의해 흙 속의 물이 수증기로 변해 날아가는 과정
\kfChoice 흙 속의 물이 농도를 맞추기 위해 저절로 뿌리 안으로 들어오는 과정
\kfChoice 물이 너무 많아서 뿌리가 숨을 쉬지 못해 썩어가는 과정
\kfChoice 겨울을 나기 위해 불필요한 잎을 스스로 떨어뜨리는 과정
\end{kfChoices}
\kfSentQ{30.}{물줄기가 위로 올라가는 원인은 무엇인가?}
\begin{kfChoices}
\kfChoice 강한 바람이 불어 나뭇가지가 심하게 흔들려서
\kfChoice 물이 계속 들어오면서 생기는 압력 때문에
\kfChoice 지구가 물체를 중심으로 끌어당기는 중력이 작용해서
\kfChoice 따뜻한 햇빛이 잎을 비추어 온도가 올라가서
\end{kfChoices}
\kfSentQ{31.}{'뿌리압'은 어떻게 만들어지는가?}
\begin{kfChoices}
\kfChoice 잎의 기공을 통해 물이 수증기로 변해 날아가면서
\kfChoice 나무가 성장함에 따라 줄기가 점점 굵고 튼튼해지면서
\kfChoice 물이 뿌리 안으로 계속 들어오면서
\kfChoice 잎이 햇빛을 듬뿍 받아 광합성을 활발하게 하면서
\end{kfChoices}
\end{kfSentenceUnit}

\begin{kfSentenceUnit}{15}{두 번째 힘은 ‘모세관 현상’이다.}
\kfSentQ{32.}{식물이 물을 올리는 두 번째 힘은 무엇인가?}
\begin{kfChoices}
\kfChoice 중력
\kfChoice 뿌리압
\kfChoice 증산 작용
\kfChoice 모세관 현상
\end{kfChoices}
\end{kfSentenceUnit}

\begin{kfSentenceUnit}{16}{물이 담긴 컵에 가느다란 빨대를 넣어보면 빨대 안쪽으로 물이 조금 올라가는 것을 볼 수 있다.}
\kfSentQ{33.}{이 문장은 무엇에 대한 예시인가?}
\begin{kfChoices}
\kfChoice 중력
\kfChoice 뿌리압
\kfChoice 모세관 현상
\kfChoice 증산 작용
\end{kfChoices}
\kfSentQ{34.}{빨대를 물에 넣으면 어떤 현상이 일어나는가?}
\begin{kfChoices}
\kfChoice 물이 빨대를 밀어낸다.
\kfChoice 빨대 안쪽으로 물이 조금 올라간다.
\kfChoice 빨대 안쪽 물이 내려간다.
\kfChoice 물이 끓는다.
\end{kfChoices}
\kfSentQ{35.}{모세관 현상을 설명하기 위해 어떤 실험을 예로 들었는가?}
\begin{kfChoices}
\kfChoice 물컵에 잉크를 떨어뜨리는 실험
\kfChoice 물컵에 가느다란 빨대를 넣는 실험
\kfChoice 물컵을 가열하는 실험
\kfChoice 물컵을 얼리는 실험
\end{kfChoices}
\end{kfSentenceUnit}

\begin{kfSentenceUnit}{17}{이처럼 가는 관과 같은 통로를 따라 액체가 올라가거나 내려가는 것을 모세관 현상이라고 한다.}
\kfSentQ{36.}{'이처럼'이 가리키는 것은 무엇인가?}
\begin{kfChoices}
\kfChoice 뜨거운 물에 의해 플라스틱 빨대가 녹는 현상
\kfChoice 컵에 담긴 물이 공기 중으로 자연스럽게 증발하는 현상
\kfChoice 빨대 안쪽으로 물이 조금 올라가는 현상
\kfChoice 실수로 컵을 건드려 물이 바닥으로 쏟아지는 현상
\end{kfChoices}
\kfSentQ{37.}{'모세관 현상'이 무엇인가?}
\begin{kfChoices}
\kfChoice 굵은 관에서 물이 멈추는 현상
\kfChoice 가는 관을 따라 액체가 이동하는 현상
\kfChoice 액체가 기체로 변하는 현상
\kfChoice 고체가 액체로 변하는 현상
\end{kfChoices}
\kfSentQ{38.}{모세관 현상이 일어나는 장소는 어디인가?}
\begin{kfChoices}
\kfChoice 넓은 바다
\kfChoice 굵은 수도관
\kfChoice 가는 관과 같은 통로
\kfChoice 메마른 땅
\end{kfChoices}
\end{kfSentenceUnit}

\begin{kfSentenceUnit}{18}{이는 물 분자가 빨대 벽에 달라붙으려는 힘이 물 분자끼리 붙으려는 힘보다 크기 때문이다.}
\kfSentQ{39.}{'이는'이 가리키는 것은 무엇인가?}
\begin{kfChoices}
\kfChoice 뿌리압
\kfChoice 증산 작용
\kfChoice 모세관 현상
\kfChoice 중력
\end{kfChoices}
\kfSentQ{40.}{모세관 현상이 일어나는 이유는 무엇인가?}
\begin{kfChoices}
\kfChoice 지구가 물체를 중심으로 끌어당기는 중력이 작용하기 때문에
\kfChoice 물 분자들 서로가 강하게 뭉치려는 응집력이 더 강하기 때문에
\kfChoice 물 분자가 관 벽에 붙으려는 힘이 더 크기 때문에
\kfChoice 외부에서 부는 바람이 관 속의 물을 위로 밀어 올리기 때문에
\end{kfChoices}
\end{kfSentenceUnit}

\begin{kfSentenceUnit}{19}{관이 가늘수록 물이 더 높이 올라가는데, 식물 몸속에는 눈에 보이지 않을 만큼 아주 가느다란 '물관'이라는 통로가 있다.}
\kfSentQ{41.}{모세관 현상에서 물을 더 높이 올라가게 하는 조건은 무엇인가?}
\begin{kfChoices}
\kfChoice 관이 굵을수록
\kfChoice 관이 가늘수록
\kfChoice 관이 짧을수록
\kfChoice 관이 길수록
\end{kfChoices}
\kfSentQ{42.}{관의 굵기와 물이 올라가는 높이의 관계는 어떠한가?}
\begin{kfChoices}
\kfChoice 비례한다.
\kfChoice 반비례한다.
\kfChoice 관계가 없다.
\kfChoice 굵을수록 높이 올라간다.
\end{kfChoices}
\kfSentQ{43.}{식물 몸속에 있는 통로의 이름은 무엇인가?}
\begin{kfChoices}
\kfChoice 식도
\kfChoice 기도
\kfChoice 혈관
\kfChoice 물관
\end{kfChoices}
\kfSentQ{44.}{물관의 특징은 무엇인가?}
\begin{kfChoices}
\kfChoice 눈에 보일 만큼 굵다.
\kfChoice 눈에 보이지 않을 만큼 아주 가늘다.
\kfChoice 길이가 매우 짧다.
\kfChoice 색깔이 빨간색이다.
\end{kfChoices}
\end{kfSentenceUnit}

\begin{kfSentenceUnit}{20}{물은 이 좁은 통로의 벽에 착 달라붙는 성질을 이용해 위로 올라간다.}
\kfSentQ{45.}{'이 좁은 통로'가 가리키는 것은 무엇인가?}
\begin{kfChoices}
\kfChoice 빨대
\kfChoice 물관
\kfChoice 뿌리
\kfChoice 잎맥
\end{kfChoices}
\kfSentQ{46.}{물이 위로 올라가기 위해 이용하는 성질은 무엇인가?}
\begin{kfChoices}
\kfChoice 좁은 통로의 벽에 착 달라붙는 성질
\kfChoice 자석의 같은 극끼리 서로 밀어내려는 성질
\kfChoice 중력에 의해 모든 물체가 아래로 떨어지려는 성질
\kfChoice 거울처럼 표면에서 빛을 반사하여 반짝이는 성질
\end{kfChoices}
\kfSentQ{47.}{물관 속 물이 위로 올라갈 수 있는 이유는 무엇인가?}
\begin{kfChoices}
\kfChoice 물이 좁은 통로의 벽에 달라붙는 성질 때문이다.
\kfChoice 물이 다른 액체보다 훨씬 무겁고 밀도가 높기 때문이다.
\kfChoice 물의 온도가 매우 낮아서 차가운 성질을 띠기 때문이다.
\kfChoice 물이 불순물 없이 맑고 투명하여 빛이 통과하기 때문이다.
\end{kfChoices}
\end{kfSentenceUnit}

\begin{kfSentenceUnit}{21}{세 번째 힘은 잎에서 물을 끌어당기는 ‘증산 작용’이다.}
\kfSentQ{48.}{식물이 물을 올리는 세 번째 힘은 무엇인가?}
\begin{kfChoices}
\kfChoice 중력
\kfChoice 뿌리압
\kfChoice 증산 작용
\kfChoice 모세관 현상
\end{kfChoices}
\kfSentQ{49.}{'증산 작용'은 어디에서 무엇을 하는 힘인가?}
\begin{kfChoices}
\kfChoice 뿌리에서 물을 미는 힘
\kfChoice 줄기에서 물을 잡는 힘
\kfChoice 잎에서 물을 끌어당기는 힘
\kfChoice 꽃에서 향기를 내는 힘
\end{kfChoices}
\end{kfSentenceUnit}

\begin{kfSentenceUnit}{22}{식물은 잎에 있는 작은 구멍을 통해 물을 수증기 형태로 내보낸다.}
\kfSentQ{50.}{식물이 물을 내보내는 장소는 어디인가?}
\begin{kfChoices}
\kfChoice 잎에 있는 작은 구멍
\kfChoice 뿌리 끝
\kfChoice 줄기 표면
\kfChoice 꽃술
\end{kfChoices}
\kfSentQ{51.}{물이 나갈 때의 형태는 무엇인가?}
\begin{kfChoices}
\kfChoice 얼음
\kfChoice 빗물
\kfChoice 수증기
\kfChoice 이슬
\end{kfChoices}
\kfSentQ{52.}{식물이 하는 행동은 무엇인가?}
\begin{kfChoices}
\kfChoice 물을 뿌리로 다시 흡수한다.
\kfChoice 물을 잎에 저장해 둔다.
\kfChoice 물을 수증기 형태로 내보낸다.
\kfChoice 물을 꽃으로 보낸다.
\end{kfChoices}
\end{kfSentenceUnit}

\begin{kfSentenceUnit}{23}{이처럼 식물체 안의 물이 잎의 기공을 통해 수증기 상태로 날아가는 현상을 증산 작용이라고 한다.}
\kfSentQ{53.}{'이처럼'이 가리키는 것은 무엇인가?}
\begin{kfChoices}
\kfChoice 뿌리가 물을 빨아들이는 것
\kfChoice 줄기에서 물이 이동하는 것
\kfChoice 잎이 햇빛을 받는 것
\kfChoice 잎의 작은 구멍으로 물이 나가는 것
\end{kfChoices}
\kfSentQ{54.}{앞 문장의 '작은 구멍'을 이 문장에서는 무엇이라고 부르는가?}
\begin{kfChoices}
\kfChoice 잎맥
\kfChoice 기공
\kfChoice 물관
\kfChoice 체관
\end{kfChoices}
\kfSentQ{55.}{증산 작용에서 물은 어떤 상태로 변하는가?}
\begin{kfChoices}
\kfChoice 고체 상태
\kfChoice 액체 상태
\kfChoice 기체(수증기) 상태
\kfChoice 젤리 상태
\end{kfChoices}
\kfSentQ{56.}{'증산 작용'이 무엇인가?}
\begin{kfChoices}
\kfChoice 뿌리에서 물을 흡수하는 현상
\kfChoice 줄기에서 물이 이동하는 현상
\kfChoice 잎에서 물이 수증기로 날아가는 현상
\kfChoice 꽃이 피는 현상
\end{kfChoices}
\end{kfSentenceUnit}

\begin{kfSentenceUnit}{24}{식물 속 물 분자들은 마치 사슬처럼 서로 단단히 이어져 있어서, 잎에서 물 분자 하나가 빠져나가면, 그 힘이 연결된 물줄기 전체를 위로 당겨준다.}
\kfSentQ{57.}{'그 힘'이 가리키는 것은 무엇인가?}
\begin{kfChoices}
\kfChoice 뿌리가 미는 힘
\kfChoice 물 분자가 서로 당기는 힘
\kfChoice 중력이 당기는 힘
\kfChoice 물 분자 하나가 빠져나가는 힘
\end{kfChoices}
\kfSentQ{58.}{물 분자들이 어떻게 연결되어 있는가?}
\begin{kfChoices}
\kfChoice 서로 멀리 떨어져 있다.
\kfChoice 서로 밀어내고 있다.
\kfChoice 서로 단단히 이어져 있다.
\kfChoice 아무런 관계가 없다.
\end{kfChoices}
\kfSentQ{59.}{물 분자 하나가 빠져나가면 어떤 일이 일어나는가?}
\begin{kfChoices}
\kfChoice 그 분자만 사라지고 끝난다.
\kfChoice 뿌리에서 물이 멈춘다.
\kfChoice 연결된 물줄기 전체가 위로 당겨진다.
\kfChoice 줄기가 얇아진다.
\end{kfChoices}
\kfSentQ{60.}{식물 속 물 분자들의 상태를 무엇에 비유했는가?}
\begin{kfChoices}
\kfChoice 사슬
\kfChoice 구슬
\kfChoice 실
\kfChoice 고무줄
\end{kfChoices}
\kfSentQ{61.}{물 분자들을 사슬에 비유한 이유는 무엇인가?}
\begin{kfChoices}
\kfChoice 사슬처럼 무겁기 때문에
\kfChoice 쉽게 끊어지기 때문에
\kfChoice 서로 연결되어 전체가 함께 움직이므로
\kfChoice 복잡하게 얽혀 있어서
\end{kfChoices}
\end{kfSentenceUnit}

\begin{kfSentenceUnit}{25}{이것이 바로 뿌리에 있는 물을 나무 꼭대기까지 끌어올리는 가장 강력한 힘이다.}
\kfSentQ{62.}{'이것'이 가리키는 것은 무엇인가?}
\begin{kfChoices}
\kfChoice 뿌리압
\kfChoice 증산 작용
\kfChoice 모세관 현상
\kfChoice 중력
\end{kfChoices}
\kfSentQ{63.}{증산 작용은 물을 어디서 어디까지 끌어올리는가?}
\begin{kfChoices}
\kfChoice 잎에서 뿌리까지
\kfChoice 뿌리에서 줄기 중간까지
\kfChoice 줄기에서 잎까지
\kfChoice 뿌리에서 나무 꼭대기까지
\end{kfChoices}
\kfSentQ{64.}{식물이 물을 끌어올리는 세 가지 힘 중 가장 강력한 힘은 무엇인가?}
\begin{kfChoices}
\kfChoice 중력
\kfChoice 뿌리압
\kfChoice 모세관 현상
\kfChoice 증산 작용
\end{kfChoices}
\kfSentQ{65.}{증산 작용이 가장 강력한 힘인 이유는 무엇인가?}
\begin{kfChoices}
\kfChoice 잎의 면적이 넓어 많은 양의 물을 한꺼번에 내보내기 때문에
\kfChoice 물 분자들이 서로 단단히 연결되어 전체가 당겨지기 때문에
\kfChoice 뜨거운 태양열이 잎에 직접 닿아 증발을 촉진하기 때문에
\kfChoice 뿌리가 땅속 깊은 곳까지 뻗어 있어 물을 잘 흡수하기 때문에
\end{kfChoices}
\end{kfSentenceUnit}

\begin{kfSentenceUnit}{26}{이처럼 식물은 어느 한 가지 힘만 사용하는 것이 아니다.}
\kfSentQ{66.}{'이처럼'이 가리키는 것은 무엇인가?}
\begin{kfChoices}
\kfChoice 뿌리압 하나만 작용하는 것
\kfChoice 물이 증발하는 현상
\kfChoice 설명한 세 가지 힘의 작용
\kfChoice 나무가 자라는 과정
\end{kfChoices}
\kfSentQ{67.}{식물이 물을 올리는 방법의 특징은 무엇인가?}
\begin{kfChoices}
\kfChoice 오직 뿌리압만 이용한다.
\kfChoice 중력의 힘만 이용한다.
\kfChoice 여러 가지 힘을 함께 사용한다.
\kfChoice 밤에만 물을 올린다.
\end{kfChoices}
\end{kfSentenceUnit}

\begin{kfSentenceUnit}{27}{뿌리에서 밀어주고, 줄기에서 타고 올라가고, 잎에서 강하게 끌어당기는 세 가지 힘을 조화롭게 이용하여, 아무리 키가 큰 나무라도 꼭대기까지 물을 보낼 수 있는 것이다.}
\kfSentQ{68.}{뿌리에서 하는 일은 무엇인가?}
\begin{kfChoices}
\kfChoice 물을 위로 밀어 올린다.
\kfChoice 물을 타고 올라간다.
\kfChoice 물을 강하게 끌어당긴다.
\kfChoice 물을 저장한다.
\end{kfChoices}
\kfSentQ{69.}{줄기에서 하는 일은 무엇인가?}
\begin{kfChoices}
\kfChoice 물을 위로 밀어 올린다.
\kfChoice 물이 관을 타고 올라간다.
\kfChoice 물을 강하게 끌어당긴다.
\kfChoice 물을 증발시킨다.
\end{kfChoices}
\kfSentQ{70.}{잎에서 하는 일은 무엇인가?}
\begin{kfChoices}
\kfChoice 물을 위로 밀어 올린다.
\kfChoice 물이 관을 타고 올라간다.
\kfChoice 물을 강하게 끌어당긴다.
\kfChoice 물을 흡수한다.
\end{kfChoices}
\end{kfSentenceUnit}

\kfSecHeader{kfAreaNonlit}{활동}{구조도}
\noindent\textit{\small 글의 흐름을 살펴보고, 구조도의 빈칸을 알맞게 채워 보자.}\par\medskip
\par\medskip\needspace{6\baselineskip}
\begin{tcolorbox}[enhanced, breakable=false,
  colback=kfPaper, colframe=kfRule, boxrule=0.4pt, arc=2pt,
  left=10pt, right=10pt, top=8pt, bottom=8pt,
  title={\footnotesize\bfseries\color{kfMute} 글의 구조도},
  coltitle=kfMute, colbacktitle=kfPrimaryLight!40,
  attach boxed title to top left={yshift=-1.5mm, xshift=6pt},
  boxed title style={colframe=kfRule, boxrule=0.3pt, arc=1pt}]
\setstretch{1.4}\footnotesize
지침: 글의 전체적인 논리 흐름을 파악하며 빈칸에 알맞은 내용을 채우시오. \\

\noindent \textbullet\ 1문단: 질문 - 나무는 어떻게 물을 위로 올릴까? \\
\noindent \quad\textbullet\ 문제 상황: ( ① )은 물을 아래로 끌어당김 \\
\noindent \quad\textbullet\ 핵심 질문: 키 큰 나무가 어떻게 꼭대기까지 ( ② )을 보낼까? \\
\noindent \textbullet\ 2문단: 세 가지 힘의 조화 \\
\noindent \quad\textbullet\ 나무가 물을 끌어올리는 비결: ( ③ ) 가지 힘의 조화 \\
\noindent \quad\textbullet\ 세 가지 힘의 이름: 뿌리압, ( ④ ) 현상, 증산 작용 \\
\noindent \textbullet\ 3-5문단: 세 가지 힘의 원리 \\
\noindent \quad\textbullet\ 힘 1: ( ⑤ ) (밀어 올리는 힘) \\
\noindent \quad\quad\textbullet\ 원리: 뿌리 안과 밖의 ( ⑥ ) 차이로 물이 들어와 압력 발생 \\
\noindent \quad\textbullet\ 힘 2: 모세관 현상 (타고 올라가는 힘) \\
\noindent \quad\quad\textbullet\ 원리: 가는 ( ⑦ ) 벽에 물 분자가 달라붙어 올라감 \\
\noindent \quad\textbullet\ 힘 3: ( ⑧ ) (끌어당기는 힘) \\
\noindent \quad\quad\textbullet\ 원리: ( ⑨ )에서 물이 수증기로 빠져나가면서 연결된 물줄기를 당김 \\
\noindent \quad\quad\textbullet\ 특징: 가장 ( ⑩ ) 힘 \\
\noindent \textbullet\ 6문단: 결론 \\
\noindent \quad\textbullet\ 최종 정리: 세 가지 힘을 ( ⑪ )롭게 이용하여 꼭대기까지 물을 보냄 \\
\end{tcolorbox}\par\medskip
\kfSecHeader{kfAreaNonlit}{문제}{}

\begin{kfQBlock}{1}{}
\kfQStem{윗글의 중심 내용으로 적절하지 \uline{않은} 것은?}
\begin{kfChoices}
\kfChoice 식물이 중력을 이기고 물을 위로 끌어올리는 원리
\kfChoice 세상에서 가장 키가 큰 나무들의 종류와 분류 기준
\kfChoice 잎의 기공을 통해 물이 빠져나가며 줄기 속 물을 끌어올리는 메커니즘
\kfChoice 식물이 물을 이동시키는 과학적 작용들의 협력 관계
\kfChoice 뿌리의 삼투압·줄기의 모세관 현상·잎의 증산 작용이 함께 작동하는 과정
\end{kfChoices}
\end{kfQBlock}
\begin{kfQBlock}{2}{}
\kfQStem{윗글의 내용과 일치하지 \uline{않는} 것은?}
\begin{kfChoices}
\kfChoice 식물은 양분을 만들 때 물을 재료로 사용한다.
\kfChoice 뿌리압은 뿌리 안과 밖의 농도 차이 때문에 발생한다.
\kfChoice 물관이 굵을수록 물이 더 높이 올라간다.
\kfChoice 증산 작용은 식물 속 물 분자들이 사슬처럼 이어진 것을 이용한다.
\kfChoice 식물은 세 가지 힘을 조화롭게 사용하여 물을 운반한다.
\end{kfChoices}
\end{kfQBlock}
\begin{kfQBlock}{3}{}
\kfQStem{흙 속의 물이 식물의 뿌리 안으로 저절로 들어오는 이유는 무엇인가?}
\begin{kfChoices}
\kfChoice 뿌리가 펌프처럼 물을 빨아들이기 때문에
\kfChoice 흙보다 뿌리 안의 액체 농도가 더 진하기 때문에
\kfChoice 식물이 광합성을 하기 위해 물이 필요하기 때문에
\kfChoice 뿌리에 있는 구멍이 흙 속의 물보다 크기 때문에
\kfChoice 중력이 흙 속의 물을 뿌리 쪽으로 밀어주기 때문에
\end{kfChoices}
\end{kfQBlock}
\begin{kfQBlock}{4}{}
\kfQStem{윗글을 통해 알 수 있는 내용이 \uline{아닌} 것은?}
\begin{kfChoices}
\kfChoice 식물이 물을 위로 올리는 데는 중력이 방해가 된다.
\kfChoice 세상에서 가장 키가 큰 나무는 세쿼이아이다.
\kfChoice 물 분자들은 서로 달라붙으려는 힘을 가지고 있다.
\kfChoice 식물은 밤보다 낮에 증산 작용을 더 활발하게 한다.
\kfChoice 뿌리압만으로는 키 큰 나무 꼭대기까지 물을 보낼 수 없다.
\end{kfChoices}
\end{kfQBlock}
\begin{kfQBlock}{5}{}
\kfQStem{<보기>의 현상을 일으키는 가장 핵심적인 식물의 작용은 무엇인가?}
\begin{kfEvidence}햇빛이 잘 드는 창가에 둔 화분의 나뭇잎에 비닐봉지를 씌우고 묶어두었더니, 잠시 후 봉지 안쪽에 물방울이 맺혔다.\end{kfEvidence}
\begin{kfChoices}
\kfChoice 뿌리압
\kfChoice 모세관 현상
\kfChoice 증산 작용
\kfChoice 농도 맞추기
\kfChoice 중력 거스르기
\end{kfChoices}
\end{kfQBlock}
\begin{kfQBlock}{1}{}
\kfQStem{식물이 살아가기 위해 꼭 필요한 것은 무엇인가요?}
\kfAnswerLines{1}
\end{kfQBlock}
\begin{kfQBlock}{2}{}
\kfQStem{지구에 있는 모든 물체를 아래로 끌어당기는 힘은 무엇인가?}
\kfAnswerLines{1}
\end{kfQBlock}
\begin{kfQBlock}{3}{}
\kfQStem{세상에서 가장 키가 큰 나무로 소개된 나무의 이름은 무엇인가?}
\kfAnswerLines{1}
\end{kfQBlock}
\begin{kfQBlock}{4}{}
\kfQStem{흙 속의 물이 뿌리 안으로 저절로 들어오면서 물줄기를 위로 밀어 올리는 힘은 무엇인가?}
\kfAnswerLines{1}
\end{kfQBlock}
\begin{kfQBlock}{5}{}
\kfQStem{액체가 아주 가는 관을 따라 저절로 올라가는 현상을 무엇이라고 하는가?}
\kfAnswerLines{1}
\end{kfQBlock}
\begin{kfQBlock}{1}{}
\kfQStem{키 큰 나무가 물을 꼭대기까지 보낼 수 있는 세 가지 힘의 이름을 쓰고, 각각의 힘이 주로 작용하는 나무의 부위를 서술하시오.}
\kfAnswerNote{답안}{4}
\end{kfQBlock}



\clearpage

\kfChapterCover{4}{러셀1 (챕터4)}{Chapter 4}{이번 챕터의 학습 내용을 시작합니다.}

\kfStartGrammar{문법 영역 — 움직임과 상태를 나타내는 말 (동사, 형용사)}


\kfSecHeader{kfAreaGrammar}{개념}{움직임과 상태를 나타내는 말 (동사, 형용사)}

우리말 문장에서 주어의 움직임이나 상태를 설명해 주는 말들이 있다. 이 말들은 문장의 끝부분에 와서 '어찌하다' 또는 '어떠하다'라는 뜻을 나타낸다. 이 단어들은 문장에서 쓰일 때 상황에 맞게 끝 부분의 모양이 바뀐다. 이렇게 모양이 바뀌는 말들을 '용언'이라고 부른다. 용언에는 크게 두 가지 종류가 있는데, 바로 동사와 형용사이다.

먼저 동사는 사람이나 사물의 움직임이나 행동을 나타내는 말이다. '어찌하다'에 해당하는 말이라고 생각하면 쉽다. 예를 들어 '아이가 밥을 먹다', '새가 하늘을 날다', '동생이 그림을 그리다' 같은 문장에서 '먹다', '날다', '그리다'가 바로 동사이다. '가다', '오다', '공부하다', '자다'처럼 우리가 하는 대부분의 행동이 동사에 속한다.

다음으로 형용사는 사람이나 사물의 상태나 성질을 나타내는 말이다. '어떠하다'에 해당하는 말이다. 예를 들어 '꽃이 예쁘다', '하늘이 높다', '바다가 푸르다' 같은 문장에서 '예쁘다', '높다', '푸르다'가 형용사이다. '착하다', '달다', '맛있다', '즐겁다'처럼 무언가의 모습, 느낌, 맛, 기분 등을 설명하는 말들이 형용사이다.

동사와 형용사는 둘 다 모양이 바뀌지만, 결정적인 차이점이 있다. 바로 현재의 일을 나타내는 '-는다' 또는 '-ㄴ다'를 붙여보면 알 수 있다. 동사는 '먹는다', '난다', '그린다', '간다'처럼 '-는다'나 '-ㄴ다'를 붙여도 자연스럽다. 하지만 형용사는 '예쁜다', '높는다', '푸른다', '착한다'처럼 '-는다'나 '-ㄴ다'를 붙이면 어색하다.

이처럼 동사와 형용사는 문장에서 하는 일이 다르다. 동사는 주어의 움직임이나 행동을 생생하게 보여준다. 형용사는 주어의 상태나 성질을 드러내 준다. 이 두 가지 말은 모두 문장 맨 뒤에서 주어를 설명하는 역할을 한다.
\par\medskip\begin{itemize}[leftmargin=18pt, itemsep=2pt, topsep=2pt, label=\textbullet]
\item 동사 — '먹다·날다·달리다'(움직임), '먹는다·먹었다·먹어서'(활용)
\item 형용사 — '예쁘다·높다·푸르다'(상태)
\item 구분(현재형 검사) — '먹는다(o)' / '예쁜다(x)'
\end{itemize}

\kfSecHeader{kfAreaGrammar}{활동}{내용 확인}
\noindent\textit{\small 아래 표의 빈칸에 알맞은 말을 채워 문법 내용을 확인해 보자.}\par\medskip
* 움직임과 상태를 나타내는 말, 용언 *

아래 표의 빈칸을 채우며 동사와 형용사의 개념과 차이점을 알아보자. 

\par\medskip\needspace{4\baselineskip}
\noindent\begin{adjustbox}{max width=\linewidth, center}
\renewcommand{\arraystretch}{1.45}
\arrayrulecolor{kfRule}
\begin{tabular}{|>{\raggedright\arraybackslash}p{\dimexpr 0.1856\linewidth-12pt\relax}|>{\raggedright\arraybackslash}p{\dimexpr 0.5052\linewidth-12pt\relax}|>{\raggedright\arraybackslash}p{\dimexpr 0.3093\linewidth-12pt\relax}|}
\hline
\rowcolor{kfPrimaryLight!50}\textbf{\color{kfPrimary} 구 분} & \textbf{\color{kfPrimary} 설 명} & \textbf{\color{kfPrimary} 예 시} \\\hline
[ ① \_\_\_\_\_\_ ] & 사람이나 사물의 움직임이나 행동을 나타내는 말 ('어찌하다') & 먹다, 날다, 그리다, 가다 \\\hline
[ ② \_\_\_\_\_\_ ] & 사람이나 사물의 상태나 성질을 나타내는 말 ('어떠하다') & 예쁘다, 높다, 푸르다, 착하다 \\\hline
공통점 & 문장에서 쓰일 때 상황에 맞게 끝부분의 모양이 바뀐다. (용언) & 먹고, 먹으니, 먹어서작고, 작으니, 작아서 \\\hline
구별 방법 & 현재의 일을 나타내는 [ ③ \_\_\_\_\_\_ ] (이)나 '-ㄴ다'를 붙여 본다. & 동사: 먹는다 (자연스러움) 형용사: 예쁜다 (어색함) \\\hline
\end{tabular}
\arrayrulecolor{black}
\end{adjustbox}\par\medskip

\kfSecHeader{kfAreaGrammar}{활동}{분석 훈련}
\noindent\textit{\small  동사, 형용사 찾기 훈련! 문장에 쓰인 동사와 형용사를 모두 찾아 ‘-다’로 끝나는 원래 모습(기본형)으로 써 보자.}\par\medskip
\begin{enumerate}[leftmargin=22pt, itemsep=6pt, topsep=4pt, label=\textbf{\arabic*.}]
\item 귀여운 아기가 잠을 잔다.
\item 나는 재미있는 영화를 봤다.
\item 빠른 기차가 달린다.
\item 붉은 해가 빛난다.
\item 아이들이 즐겁게 논다.
\item 우리는 따뜻한 방에 모였다.
\item 슬픈 음악을 들으니 내 마음도 슬프다.
\item 나는 매운 음식을 먹는다.
\item 착한 아이가 집으로 온다.
\item 친구가 밝게 웃는다.
\end{enumerate}

\kfSecHeader{kfAreaGrammar}{문제}{}

\begin{kfQBlock}{1}{}
\kfQStem{다음 단어들 중 성격이 \uline{다른} 하나는 무엇인가?}
\begin{kfChoices}
\kfChoice 슬프다
\kfChoice 자다
\kfChoice 놀다
\kfChoice 읽다
\kfChoice 보다
\end{kfChoices}
\end{kfQBlock}
\begin{kfQBlock}{2}{}
\kfQStem{다음 문장의 굵게 표시된 단어에 대한 설명으로 옳은 것은?}
\begin{kfChoices}
\kfChoice 다른 말을 꾸미는 부사다.
\kfChoice 주인공의 상태를 설명하는 형용사다.
\kfChoice 문장에서 홀로 쓰이는 감탄사다.
\kfChoice 주인공의 움직임을 나타내는 동사다.
\kfChoice 주인공 역할을 하는 명사다.
\end{kfChoices}
\end{kfQBlock}
\begin{kfQBlock}{3}{}
\kfQStem{다음 중 ‘움직임’을 나타내는 말이 포함된 문장은?}
\begin{kfChoices}
\kfChoice 아기가 잠을 잔다.
\kfChoice 날씨가 좋다.
\kfChoice 바다가 넓다.
\kfChoice 나는 학생이다.
\kfChoice 산이 높다.
\end{kfChoices}
\end{kfQBlock}
\begin{kfQBlock}{4}{}
\kfQStem{다음 중 성질이나 ‘상태’를 나타내는 말이 포함된 문장은?}
\begin{kfChoices}
\kfChoice 바다가 아주 넓다.
\kfChoice 동생이 밥을 먹는다.
\kfChoice 새가 노래를 부른다.
\kfChoice 나는 공을 던진다.
\kfChoice 기차가 역을 떠났다.
\end{kfChoices}
\end{kfQBlock}
\begin{kfQBlock}{5}{}
\kfQStem{다음 단어들 중 품사의 종류가 나머지와 \uline{다른} 하나는?}
\begin{kfChoices}
\kfChoice 높다
\kfChoice 작다
\kfChoice 길다
\kfChoice 달리다
\kfChoice 빠르다
\end{kfChoices}
\end{kfQBlock}
\begin{kfQBlock}{6}{}
\kfQStem{다음 밑줄 친 단어 중 품사가 \uline{다른} 하나는?}
\begin{kfChoices}
\kfChoice 새가 \uline{난다}.
\kfChoice 나는 \uline{슬프다}.
\kfChoice 물이 \uline{흐른다}.
\kfChoice 해가 \uline{솟는다}.
\kfChoice 아기가 \uline{웃는다}.
\end{kfChoices}
\end{kfQBlock}
\begin{kfQBlock}{7}{}
\kfQStem{다음 중 현재의 일을 나타내는 '-는다'나 '-ㄴ다'를 붙였을 때 어색한 말은?}
\begin{kfChoices}
\kfChoice 가다
\kfChoice 예쁘다
\kfChoice 자다
\kfChoice 오다
\kfChoice 보다
\end{kfChoices}
\end{kfQBlock}
\begin{kfQBlock}{8}{}
\kfQStem{다음 중 현재의 일을 나타내는 '-는다'나 '-ㄴ다'를 붙였을 때 자연스러운 말은?}
\begin{kfChoices}
\kfChoice 높다
\kfChoice 날다
\kfChoice 푸르다
\kfChoice 착하다
\kfChoice 맛있다
\end{kfChoices}
\end{kfQBlock}
\begin{kfQBlock}{1}{}
\kfQStem{문장에서 쓰일 때 상황에 맞게 끝부분의 모양이 바뀌는 말을 무엇이라고 하는가?}
\kfAnswerLines{1}
\end{kfQBlock}
\begin{kfQBlock}{2}{}
\kfQStem{사람이나 사물의 '움직임'이나 '행동'을 나타내는 품사는 무엇인가?}
\kfAnswerLines{1}
\end{kfQBlock}
\begin{kfQBlock}{3}{}
\kfQStem{사람이나 사물의 '상태'나 '성질'을 나타내는 품사는 무엇인가?}
\kfAnswerLines{1}
\end{kfQBlock}
\begin{kfQBlock}{4}{}
\kfQStem{동사는 '어떠하다', '어찌하다' 중 주로 어떤 물음에 대한 답이 되는가?}
\kfAnswerLines{1}
\end{kfQBlock}
\begin{kfQBlock}{5}{}
\kfQStem{형용사는 '어떠하다', '어찌하다' 중 주로 어떤 물음에 대한 답이 되는가?}
\kfAnswerLines{1}
\end{kfQBlock}
\begin{kfQBlock}{6}{}
\kfQStem{동사와 형용사를 구별할 때, 현재의 일을 나타내는 어떤 말을 붙여보는가?}
\kfAnswerLines{1}
\end{kfQBlock}
\begin{kfQBlock}{7}{}
\kfQStem{'예쁘다', '착하다', '시원하다'는 모두 어떤 품사에 속하는가?}
\kfAnswerLines{1}
\end{kfQBlock}
\begin{kfQBlock}{8}{}
\kfQStem{'달리다', '먹다', '공부하다'는 모두 어떤 품사에 속하는가?}
\kfAnswerLines{1}
\end{kfQBlock}
\begin{kfQBlock}{9}{}
\kfQStem{'자다', '웃다', '만들다'는 모두 어떤 품사에 속하는가?}
\kfAnswerLines{1}
\end{kfQBlock}
\begin{kfQBlock}{10}{}
\kfQStem{'높다', '푸르다', '길다'는 모두 어떤 품사에 속하는가?}
\kfAnswerLines{1}
\end{kfQBlock}
\begin{kfQBlock}{1}{}
\kfQStem{‘기쁘다'는 동사와 형용사 중 무엇일까? 지문에서 배운 구별 방법을 사용하여 그 이유를 간단히 설명하시오.}
\kfAnswerNote{답안}{4}
\end{kfQBlock}


\kfStartConcept{개념 영역 — 조선 시대의 노래, 시조}


\kfSecHeader{kfAreaConcept}{개념}{조선 시대의 노래, 시조}

시조는 고려 말에 발생하여 조선 시대를 거쳐 오늘날까지 이어지는 우리나라의 대표적인 '정형시(정해진 형식이 있는 시)'이다. 우리 고유의 문학 갈래로, 짧은 형식 속에 다양한 주제를 담아냈다.

시조의 가장 큰 형식적 특징은 '3장(章) 6구(句)'의 정해진 틀을 가진다는 점이다. 첫째 줄을 '초장', 둘째 줄을 '중장', 셋째 줄을 '종장'이라고 부르며(3장), 각 장은 다시 2개의 '구(句)'로 나눌 수 있어 총 6구를 이룬다. 전체 글자 수는 45자 내외로 매우 간결한 형태이다.

시조는 일정한 리듬감을 가진다. 각 장은 "이 몸이 / 죽고 죽어 / 일백 번 / 고쳐 죽어"처럼 네 마디로 끊어 읽는 '4음보'의 율격을 기본으로 한다. 이때 각 음보(마디)는 3글자 혹은 4글자로 이루어지는 '3.4조(또는 4.4조)'의 글자 수(음수율)를 갖는 것이 일반적이다.

시조가 정형시로서 지켜야 할 가장 엄격한 규칙은 종장에 있다. 종장의 첫 마디(첫 음보)는 반드시 3글자(3음절)로 고정해야 한다. 또한, 종장 전체의 글자 수는 '3/5/4/3'의 틀을 지키는 것을 정석으로 한다.

시조는 형태와 구성에 따라 평시조, 연시조, 사설시조 등으로 나뉜다.

평시조는 위에서 설명한 3장 6구, 4음보, 3.4조 음수율, 종장 규칙(3/5/4/3) 등을 엄격하게 지키는 가장 기본적이고 단일한 형태의 시조이다.

연시조는 이러한 평시조가 한 편으로 끝나지 않고, 하나의 큰 주제 아래 여러 수(首)가 연결되어 이루어진 시를 말한다. 맹사성의 「강호사시가」가 대표적인 예로, '자연 속 사계절'이라는 하나의 주제로 봄, 여름, 가을, 겨울 네 수의 시조가 묶여 있다. 조선 시대 사대부들은 주로 평시조나 연시조를 통해 자연의 아름다움이나 임금에 대한 충성심 등을 노래했다.

반면 사설시조는 조선 후기에 주로 평민들에 의해 창작되었으며, 형식이 파괴된 것이 특징이다. 초장과 중장이 4음보의 율격이나 글자 수 제한(3.4조)에서 벗어나 매우 길어진다. 평민들은 사설시조를 통해 일상의 솔직한 감정이나 사회에 대한 비판을 재미있고 해학적으로 담아냈다.
\par\medskip\begin{itemize}[leftmargin=18pt, itemsep=2pt, topsep=2pt, label=\textbullet]
\item 평시조 — 정몽주 「단심가」(3장 6구, 종장 첫 음보 3자)
\item 연시조 — 맹사성 「강호사시가」(사계절을 4수로 연결)
\item 사설시조 — 조선 후기 평민층, 중장이 길게 늘어진 형태
\end{itemize}

\kfSecHeader{kfAreaConcept}{문제}{}

\begin{kfQBlock}{1}{}
\kfQStem{지문의 내용으로 보아 ‘시조’에 대한 설명으로 적절하지 \uline{않은} 것은 무엇인가?}
\begin{kfChoices}
\kfChoice 운율은 있으나 정해진 형식이 없는 자유시의 한 종류이다.
\kfChoice 초장·중장·종장의 3장 형식을 갖춘다.
\kfChoice 한 장은 보통 4음보의 율격을 가진다.
\kfChoice 우리 문학의 대표적인 정형시 갈래로 평시조·연시조·사설시조 등으로 나뉜다.
\kfChoice 고려 말에 발생하여 오늘날까지 이어져 온 정형시이다.
\end{kfChoices}
\end{kfQBlock}
\begin{kfQBlock}{2}{}
\kfQStem{평시조의 형식적 특징으로 알맞지 \uline{않은} 것은 무엇인가?}
\begin{kfChoices}
\kfChoice 중장이 제한 없이 매우 길어진다.
\kfChoice 전체 글자 수는 45자 내외이다.
\kfChoice 각 장은 4음보의 율격을 가진다.
\kfChoice 3.4조의 음수율을 갖는 것이 일반적이다.
\kfChoice 3장 6구의 구조를 가진다.
\end{kfChoices}
\end{kfQBlock}
\begin{kfQBlock}{3}{}
\kfQStem{시조가 ‘정형시’로서 지켜야 할 규칙으로 적절하지 \uline{않은} 것은 무엇인가?}
\begin{kfChoices}
\kfChoice 종장의 첫 마디는 반드시 3글자로 고정되어야 한다.
\kfChoice 자연의 아름다움만을 주제로 노래해야 한다는 형식 규정이 있다.
\kfChoice 초장·중장·종장의 3장 6구 구조를 지켜야 한다.
\kfChoice 3·4조 또는 4·4조의 음수율을 기본으로 삼아야 한다.
\kfChoice 한 장은 보통 네 마디씩 4음보로 끊어 읽어야 한다.
\end{kfChoices}
\end{kfQBlock}
\begin{kfQBlock}{4}{}
\kfQStem{<보기>는 정몽주의 「단심가」이다. 이 시조의 율격을 바르게 설명한 것은 무엇인가?}
\begin{kfEvidence}이 몸이 / 죽고 죽어 / 일백 번 / 고쳐 죽어\end{kfEvidence}
\begin{kfChoices}
\kfChoice 3음보
\kfChoice 4음보
\kfChoice 5음보
\kfChoice 7.5조
\kfChoice 3.3.2조
\end{kfChoices}
\end{kfQBlock}
\begin{kfQBlock}{5}{}
\kfQStem{시조의 종장에 대한 설명으로 적절하지 \uline{않은} 것은 무엇인가?}
\begin{kfChoices}
\kfChoice ‘아아’와 같은 감탄사로 시작하는 것이 정해진 규칙이다.
\kfChoice 보통 3·5·4·3의 음수율을 정석으로 한다.
\kfChoice 시조의 마무리 역할을 하며 주제를 압축하여 드러낸다.
\kfChoice 작품의 정서를 갈무리하고 깊은 여운을 남기는 부분이다.
\kfChoice 첫 마디는 3음절로 고정해야 한다.
\end{kfChoices}
\end{kfQBlock}
\begin{kfQBlock}{6}{}
\kfQStem{맹사성의 「강호사시가」에 대한 설명으로 적절하지 \uline{않은} 것은 무엇인가?}
\begin{kfChoices}
\kfChoice ‘자연 속 사계절’을 주제로 묶인 연시조이다.
\kfChoice 조선 후기 평민들이 창작한 사설시조에 해당한다.
\kfChoice 자연을 즐기는 화자의 태도를 사계절의 흐름에 따라 노래한다.
\kfChoice 시조의 기본 형식인 3장 6구·4음보를 지키고 있다.
\kfChoice 평시조 4수가 모여 한 작품을 이루는 연시조이다.
\end{kfChoices}
\end{kfQBlock}
\begin{kfQBlock}{7}{}
\kfQStem{<보기>는 조선 후기에 창작된 작자 미상의 시조이다. 이 시조의 갈래는 무엇인가?}
\begin{kfEvidence}개를 여라믄이나 기르되 요 개같이 얄미우랴 미운 임 오며는 꼬리를 홰홰 치며 칩뛰락 나리뛰락 반겨서 내닫고 고운 임 오며는 뒷발을 버둥버둥 므르락 나으락 캉캉 짖어서 도라가게 한다 쉰밥이 그릇그릇 난들 너 머길 줄이 이시랴\end{kfEvidence}
\begin{kfChoices}
\kfChoice 평시조
\kfChoice 연시조
\kfChoice 사설시조
\kfChoice 고려가요
\kfChoice 현대시
\end{kfChoices}
\end{kfQBlock}
\begin{kfQBlock}{8}{}
\kfQStem{사설시조에 대한 설명으로 알맞지 \uline{않은} 것은 무엇인가?}
\begin{kfChoices}
\kfChoice 3장 6구, 45자 내외의 엄격한 형식을 지켰다.
\kfChoice 주로 평민들에 의해 창작되었다.
\kfChoice 중장이 4음보 율격에서 벗어나 매우 길어진다.
\kfChoice 일상의 솔직한 감정이나 사회 비판을 담았다.
\kfChoice 조선 후기에 주로 창작되었다.
\end{kfChoices}
\end{kfQBlock}
\begin{kfQBlock}{1}{}
\kfQStem{고려 말에 발생하여 조선 시대를 거쳐 이어진, 정해진 형식이 있는 우리나라의 대표적인 시를 무엇이라고 하는가?}
\kfAnswerLines{1}
\end{kfQBlock}
\begin{kfQBlock}{2}{}
\kfQStem{시조는 '초장', '중장', '종장'으로 구성된다. 이를 합쳐 몇 장 구성이라고 하는가?}
\kfAnswerLines{1}
\end{kfQBlock}
\begin{kfQBlock}{3}{}
\kfQStem{시조는 한 장을 네 마디로 끊어 읽는 규칙적인 리듬을 가진다. 이 리듬의 단위를 무엇이라고 하는가?}
\kfAnswerLines{1}
\end{kfQBlock}
\begin{kfQBlock}{4}{}
\kfQStem{시조 각 음보의 글자 수가 3글자 혹은 4글자로 반복되는 것을 무엇이라고 하는가?}
\kfAnswerLines{1}
\end{kfQBlock}
\begin{kfQBlock}{5}{}
\kfQStem{시조의 3장 중, 형식이 가장 엄격하게 정해져 있는 장은 어디인가?}
\kfAnswerLines{1}
\end{kfQBlock}
\begin{kfQBlock}{6}{}
\kfQStem{시조의 종장에서 첫 번째 음보는 반드시 몇 글자로 고정되어야 하는가?}
\kfAnswerLines{1}
\end{kfQBlock}
\begin{kfQBlock}{7}{}
\kfQStem{하나의 큰 주제 아래 여러 수의 평시조가 연결되어 이루어진 시조의 형식을 무엇이라고 하는가?}
\kfAnswerLines{1}
\end{kfQBlock}
\begin{kfQBlock}{1}{}
\kfQStem{'평시조'와 '사설시조'는 형식과 내용 면에서 뚜렷한 차이를 보인다. 두 시조의 형식적 차이점을 <조건>에 맞게 서술하시오.}
\kfAnswerNote{답안}{4}
\end{kfQBlock}


\kfStartLit{문학 영역 — 작품}


\kfSecHeader{kfAreaLiterature}{지문}{작품}
\kfPassageNote{작품}{%
\textbackslash{}<춘사 (봄 노래)\textbackslash{}> \\[14pt]
자연에 봄이 오니 깊은 흥이 저절로 난다 \\[4pt]
시냇가에서 막걸리 마시며 고기를 안주 삼아 \\[4pt]
이 몸이 여유롭게 지내는 것도 임금님 덕분이다 \\[14pt]
\textbackslash{}<하사 (여름 노래)\textbackslash{}> \\[14pt]
자연에 여름이 오니 초가집에 할 일이 없다 \\[4pt]
믿음직한 강물은 시원한 바람을 보내주네 \\[4pt]
이 몸이 시원하게 지내는 것도 임금님 덕분이다 \\[14pt]
 \textbackslash{}<추사 (가을 노래)\textbackslash{}> \\[14pt]
자연에 가을이 오니 물고기마다 살이 올랐다 \\[4pt]
작은 배에 그물을 싣고 물 따라 흘러가며 \\[4pt]
이 몸이 재미있게 지내는 것도 임금님 덕분이다 \\[14pt]
\textbackslash{}<동사 (겨울 노래)\textbackslash{}> \\[14pt]
자연에 겨울이 오니 눈의 깊이가 한 자 넘게 쌓였다 \\[4pt]
삿갓을 기울여 쓰고 도롱이를 옷 삼아 입으니 \\[4pt]
이 몸이 춥지 않게 지내는 것도 임금님 덕분이다
}


\kfSecHeader{kfAreaLiterature}{활동}{문학 문장 독해}
\noindent\textit{\small 작품 속 문장을 자세히 읽고, 문장별로 주어진 물음에 답해 보자.}\par\medskip
\begin{kfSentenceUnit}{1}{"자연에 봄이 오니 깊은 흥이 저절로 난다."}
\kfSentQ{1.}{다음 문장에서 서술어를 기준으로 각각의 주어는 무엇인가? (1) 서술어 '오니'의 주어는 무엇인가?}
\begin{kfChoices}
\kfChoice 자연에
\kfChoice 봄이
\kfChoice 흥이
\kfChoice 저절로
\end{kfChoices}
\kfSentQ{1.}{(2) 서술어 '떠다니네'의 주어는 무엇인가?}
\begin{kfChoices}
\kfChoice 자연이
\kfChoice 봄이
\kfChoice 깊은 흥이
\kfChoice 저절로
\end{kfChoices}
\end{kfSentenceUnit}

\begin{kfSentenceUnit}{2}{"자연에 여름이 오니 초가집에 할 일이 없다."}
\kfSentQ{2.}{다음 문장에서 서술어를 기준으로 각각의 주어는 무엇인가? (1) 서술어 '오니'의 주어는 무엇인가?}
\begin{kfChoices}
\kfChoice 자연에
\kfChoice 여름이
\kfChoice 초가집에
\kfChoice 할 일이
\end{kfChoices}
\kfSentQ{2.}{(2) 서술어 '없다'의 주어는 무엇인가?}
\begin{kfChoices}
\kfChoice 여름이
\kfChoice 초가집이
\kfChoice 할 일이
\kfChoice 자연이
\end{kfChoices}
\end{kfSentenceUnit}

\begin{kfSentenceUnit}{3}{"이 몸이 춥지 않게 지내는 것도 임금님 덕분이다."}
\kfSentQ{3.}{화자가 자연을 즐기고 임금을 그리워하는 대상은 무엇인가?}
\begin{kfChoices}
\kfChoice 임금님
\kfChoice 말하는 이 (화자)
\kfChoice 삿갓
\kfChoice 도롱이
\end{kfChoices}
\end{kfSentenceUnit}

\kfSecHeader{kfAreaLiterature}{활동}{해설 학습}
\noindent\textit{\small 작품 해설을 단원별로 읽고, 각 단원에 딸린 물음에 답해 보자.}\par\medskip
\par\noindent\textbf{\color{kfPrimary} 해설 단원 1}\par\smallskip
\kfPassageNote{해설 1}{%
「강호사시가」는 맹사성이 쓴 우리나라 최초의 연시조이다. 연시조란 여러 수의 시조를 연결해서 하나의 큰 주제를 노래하는 형식을 말한다. 이 작품은 총 4수로 이루어져 있으며, 봄·여름·가을·겨울 사계절의 아름다움을 차례대로 노래했다.

'강호'는 '강과 호수'라는 뜻으로, 자연을 의미한다. 즉 이 시는 자연 속에서 사계절을 보내며 느끼는 즐거움과 만족감을 표현한 작품이다. 각 계절마다 그 계절에 어울리는 풍경과 놀이를 제시하면서, 자연과 함께하는 삶의 아름다움을 보여준다.

특히 각 수마다 "자연에 \textasciitilde{}이 오니"로 시작하고 "임금님 덕분이다"로 끝나는 같은 구조를 반복한다. 이런 반복은 시에 운율감을 주면서도 주제 의식을 강조하는 역할을 한다. 자연을 즐기면서도 임금의 은혜를 잊지 않는 조선시대 선비들의 마음가짐을 잘 보여주는 대표적인 작품이다.
}

\begin{kfQBlock}{1}{문제}
\kfQStem{향가는 몇 구체, 8구체, 10구체로 나뉘는가?}
\kfAnswerLines{1}
\end{kfQBlock}

\begin{kfQBlock}{2}{문제}
\kfQStem{중장(제2장)에서 '갈매기'는 화자와 함께 무엇을 즐기는 대상인가?}
\kfAnswerLines{1}
\end{kfQBlock}

\begin{kfQBlock}{3}{문제}
\kfQStem{각 계절 노래(수)는 어떤 말로 끝나는가?}
\kfAnswerLines{1}
\end{kfQBlock}

\begin{kfQBlock}{undefined}{문제}
\kfQStem{종장의 마지막 구 '임금 은혜를 이제 더욱 아노이다'에서 화자는 자연 속에 있으면서도 누구의 은혜를 잊지 않고 있는가?}
\kfAnswerLines{1}
\end{kfQBlock}

\begin{kfQBlock}{5}{문제}
\kfQStem{이러한 마음가짐은 어느 시대 선비들의 특징인가?}
\kfAnswerLines{1}
\end{kfQBlock}

\par\noindent\textbf{\color{kfPrimary} 해설 단원 2}\par\smallskip
\kfPassageNote{해설 2}{%
이 시의 내용은 사계절의 흐름에 따라 전개된다. 봄에는 시냇가에서 막걸리를 마시며 물고기를 안주로 삼는 흥겨운 모습을, 여름에는 초가집에서 강바람을 맞으며 더위를 식히는 여유로운 모습을 그린다.

가을에는 살이 오른 물고기들을 보며 작은 배에 그물을 싣고 물 따라 떠다니는 한가로운 모습을, 겨울에는 깊이 쌓인 눈 속에서도 삿갓과 도롱이 차림으로 추위를 즐기는 모습을 보여준다. 이는 단순히 계절의 변화를 묘사한 것이 아니라, 각 계절마다 자연과 함께하는 특별한 즐거움이 있음을 말해준다.

특히 주목할 점은 화자가 실제로 농사를 짓거나 생계를 위해 고기를 잡는 것이 아니라는 점이다. 그물을 "물 따라 흘러가며" 던져두는 모습에서 알 수 있듯이, 이는 풍류를 즐기기 위한 놀이이다. 이처럼 화자는 자연 속에서 여유롭고 한가로운 삶을 즐기고 있으며, 이러한 평안한 삶이 모두 임금의 은혜 덕분이라고 생각하고 있다.
}

\begin{kfQBlock}{1}{문제}
\kfQStem{봄에 시냇가에서 막걸리를 마실 때 안주로 삼은 것은 무엇인가?}
\kfAnswerLines{1}
\end{kfQBlock}

\begin{kfQBlock}{2}{문제}
\kfQStem{여름에 말하는 이는 어디에서 강바람을 맞으며 더위를 식혔는가?}
\kfAnswerLines{1}
\end{kfQBlock}

\begin{kfQBlock}{3}{문제}
\kfQStem{가을에 말하는 이는 작은 배에 무엇을 싣고 물 따라 떠다녔는가?}
\kfAnswerLines{1}
\end{kfQBlock}

\begin{kfQBlock}{4}{문제}
\kfQStem{겨울에 말하는 이는 눈 속에서 무엇과 무엇을 착용했는가?}
\kfAnswerLines{1}
\end{kfQBlock}

\begin{kfQBlock}{5}{문제}
\kfQStem{가을에 그물을 배에 싣고 떠다닌 것은 생계를 위한 일이었는가, 놀이였는가?}
\kfAnswerLines{1}
\end{kfQBlock}

\begin{kfQBlock}{6}{문제}
\kfQStem{해설문에 따르면, 화자는 자연 속에서 한가롭게 무엇을 즐기고 있는가?}
\kfAnswerLines{1}
\end{kfQBlock}

\begin{kfQBlock}{7}{문제}
\kfQStem{화자는 땀을 식히기 위해 무슨 바람을 쐬었는가?}
\kfAnswerLines{1}
\end{kfQBlock}

\par\noindent\textbf{\color{kfPrimary} 해설 단원 3}\par\smallskip
\kfPassageNote{해설 3}{%
이 시에 나오는 여러 시어들은 각각 특별한 의미를 담고 있다. '막걸리'는 화자의 소박하고 서민적인 생활을 보여준다. 비싼 술이 아닌 막걸리를 마시는 것에서 화자가 검소한 생활을 하고 있음을 알 수。

'초가집'은 화자가 사는 집으로, 기와집이 아닌 볏짚으로 지붕을 덮은 소박한 집을 의미한다. '작은 배'와 '그물'은 화자가 강에서 고기잡이를 하는 모습을 보여주지만, 실제로는 생계를 위한 것이 아니라 자연을 즐기기 위한 풍류 활동임을 나타낸다.

'삿갓'과 '도롱이'는 겨울철 화자의 차림새를 나타내는데, 이 역시 소박하고 검소한 생활을 상징한다. 하지만 이런 소박한 차림으로도 추위를 느끼지 않는다고 하여, 물질적인 풍요보다는 정신적인 만족을 중시하는 화자의 가치관을 보여준다. 이는 가난해도 마음이 편안하고 도를 즐기며 사는 '안빈낙도(安貧樂道)'의 자세와 자신의 분수를 알고 만족하며 사는 '안분지족(安分知足)'의 삶의 태도를 잘 나타낸다.

무엇보다 중요한 것은 '임금의 은혜'라는 표현이다. 화자는 자연에서 누리는 모든 즐거움의 근원을 임금의 은혜로 돌리고 있다. 이는 자연과 정치 현실이 따로 분리된 것이 아니라, 임금의 다스림이라는 하나의 질서 안에 있다고 보는 조선시대 선비들의 세계관을 보여준다.
}

\begin{kfQBlock}{1}{문제}
\kfQStem{화자가 나귀를 타고 들어간 곳은 깊은 산속의 무엇인가?}
\kfAnswerLines{1}
\end{kfQBlock}

\begin{kfQBlock}{2}{문제}
\kfQStem{'바람', '명월'은 화자가 즐기는 무엇을 상징하는가?}
\kfAnswerLines{1}
\end{kfQBlock}

\begin{kfQBlock}{3}{문제}
\kfQStem{가난해도 마음이 편안하고 도를 즐기며 사는 태도를 한자성어로 무엇이라고 하는가?}
\kfAnswerLines{1}
\end{kfQBlock}

\begin{kfQBlock}{4}{문제}
\kfQStem{화자는 자연의 흥취에 젖어 나귀를 타고 산천을 구경하고 있는데, 이는 무엇을 돌리고 있는 것과 같은가?}
\kfAnswerLines{1}
\end{kfQBlock}

\begin{kfQBlock}{5}{문제}
\kfQStem{자신의 분수를 알고 만족하며 사는 삶의 태도를 한자성어로 무엇이라고 하는가?}
\kfAnswerLines{1}
\end{kfQBlock}

\kfSecHeader{kfAreaLiterature}{문제}{}

\begin{kfQBlock}{1}{}
\kfQStem{이 시의 각 계절 노래마다 공통적으로 반복되는 내용으로 적절하지 \uline{않은} 것은 무엇인가?}
\begin{kfChoices}
\kfChoice 친구들과 어울려 큰 잔치를 벌이며 즐기는 모습
\kfChoice ‘이 몸이 한가하옴도 역군은이샷다’와 같은 종장의 반복
\kfChoice 강호(자연)에서 한가롭고 즐겁게 지내는 화자의 모습
\kfChoice 자연을 즐기는 일과 임금에 대한 충의를 함께 드러내는 태도
\kfChoice 자연 속에서의 한가로움이 임금의 은혜 덕분이라는 생각
\end{kfChoices}
\end{kfQBlock}
\begin{kfQBlock}{2}{}
\kfQStem{이 시에 나타난 말하는 이의 삶의 태도로 적절하지 \uline{않은} 것은 무엇인가?}
\begin{kfChoices}
\kfChoice 자연을 개발하여 더 편리한 삶을 만들려는 적극적 개척의 태도
\kfChoice 자연의 사계절을 즐기며 한가로움을 누리려는 태도
\kfChoice 자연을 즐기는 가운데 임금의 은혜를 잊지 않으려는 충의의 태도
\kfChoice 욕심 없이 자족하며 자연과 더불어 사는 안빈낙도의 태도
\kfChoice 자연 속에서 소박하게 사는 삶에 만족하는 태도
\end{kfChoices}
\end{kfQBlock}
\begin{kfQBlock}{3}{}
\kfQStem{<추사(가을 노래)>에서 “작은 배에 그물을 싣고 물 따라 흘러가며”라는 구절을 통해 알 수 있는 것은 무엇인가?}
\begin{kfChoices}
\kfChoice 가족의 생계를 위해 부지런히 고기를 잡는 어부의 모습
\kfChoice 그물을 다루는 솜씨가 서툴러 어쩔 줄 모르는 모습
\kfChoice 고기잡이보다 가을 강의 풍경을 즐기는 풍류의 모습
\kfChoice 빨리 고기를 잡아 가족에게 가져다주려는 다급한 모습
\kfChoice 큰 배를 살 돈이 없어 작은 배를 탈 수밖에 없는 모습
\end{kfChoices}
\end{kfQBlock}
\begin{kfQBlock}{4}{}
\kfQStem{말하는 이가 자연 속에서 입고 사용하는 물건으로, 그의 소박한 삶을 보여주는 것은 무엇인가?}
\begin{kfChoices}
\kfChoice 비단옷과 옥으로 만든 갓
\kfChoice 금으로 만든 낚싯대
\kfChoice 아주 크고 화려한 배
\kfChoice 삿갓과 도롱이
\kfChoice 가죽으로 만든 신발
\end{kfChoices}
\end{kfQBlock}
\begin{kfQBlock}{5}{}
\kfQStem{<하사(여름 노래)>에서 ‘믿음직한 강물’이라는 표현에 담긴 의미로 적절하지 \uline{않은} 것은 무엇인가?}
\begin{kfChoices}
\kfChoice 시원한 바람을 보내 주어 더위를 식혀 주는 강물에 대한 신뢰
\kfChoice 강물에 대한 믿음을 가져야 복을 받는다는 미신적 사고
\kfChoice 강이 주는 한가로운 정취 속에서 자연과 하나 됨을 느끼는 마음
\kfChoice 강물을 의인화하여 자연을 다정한 벗처럼 받아들이는 태도
\kfChoice 변함없이 흐르며 화자에게 안식을 주는 자연에 대한 친근감
\end{kfChoices}
\end{kfQBlock}
\begin{kfQBlock}{6}{}
\kfQStem{이 시에 대한 설명으로 알맞지 \uline{않은} 것은 무엇인가?}
\begin{kfChoices}
\kfChoice 사계절의 흐름에 따라 시를 전개하고 있다.
\kfChoice 각 계절마다 자연을 즐기는 모습이 나타난다.
\kfChoice 자연 속에서의 즐거움을 임금의 은혜와 연결 짓고 있다.
\kfChoice 자연을 벗어나 다시 벼슬길로 돌아가고 싶은 마음을 드러낸다.
\kfChoice 비슷한 문장 구조를 반복하여 운율을 형성하고 있다.
\end{kfChoices}
\end{kfQBlock}
\begin{kfQBlock}{7}{}
\kfQStem{<추사(가을 노래)>에서 "물고기마다 살이 올랐다"는 표현은 가을의 어떤 모습을 보여주는가?}
\begin{kfChoices}
\kfChoice 쓸쓸함
\kfChoice 풍요로움
\kfChoice 화려함
\kfChoice 서늘함
\kfChoice 깨끗함
\end{kfChoices}
\end{kfQBlock}
\begin{kfQBlock}{8}{}
\kfQStem{이 시가 우리나라 시의 역사에서 중요한 작품으로 평가받는 이유로 적절하지 \uline{않은} 것은 무엇인가?}
\begin{kfChoices}
\kfChoice 임금을 강하게 비판하는 정치 풍자시의 시초이기 때문이다.
\kfChoice 자연 친화적 삶의 태도를 보여 주는 강호 시조의 출발점이 되기 때문이다.
\kfChoice 시조의 형식미와 사대부의 충의 사상을 함께 잘 드러내기 때문이다.
\kfChoice 사계절의 흐름에 따라 통일된 주제를 노래하는 작품 구성을 보여 주기 때문이다.
\kfChoice 여러 수의 시조가 하나의 제목 아래 묶인 최초의 연시조이기 때문이다.
\end{kfChoices}
\end{kfQBlock}
\begin{kfQBlock}{1}{}
\kfQStem{이 시에서 말하는 이가 소박하지만 만족스러운 삶을 살고 있음을 보여주는 모습 두 가지를 찾아 서술하시오.}
\kfAnswerNote{답안}{4}
\end{kfQBlock}
\begin{kfQBlock}{2}{}
\kfQStem{이 시의 각 계절 노래는 비슷한 구조로 되어 있다. 그 공통적인 구조를 세 단계로 나누어 서술하시오.}
\kfAnswerNote{답안}{4}
\end{kfQBlock}


\kfStartNonlit{비문학 영역 — [기술] 가상과 현실의 만남, 가상현실과 증강현실}


\kfSecHeader{kfAreaNonlit}{지문}{[기술] 가상과 현실의 만남, 가상현실과 증강현실}
\kfPassageNote{[기술] 가상과 현실의 만남, 가상현실과 증강현실}{%
영화나 게임 속 주인공처럼 새로운 세상을 탐험하거나, 내 방에 공룡이 나타나는 모습을 상상해 본 적이 있는가? 이러한 상상을 현실로 만드는 기술이 바로 가상현실(Virtual Reality, 이하 VR)과 증강현실(Augmented Reality, 이하 AR)이다. 두 기술은 실감 나는 체험을 제공한다는 점에서는 같지만, 작동 방식에는 분명한 차이가 있다. 가상현실과 증강현실은 무엇이며, 어떻게 다를까?

먼저 가상현실은 우리를 완전히 새로운 디지털 세상으로 안내하는 기술이다. 사용자는 머리에 쓰는 특수한 안경 형태의 기기를 착용한다. 이 기기는 현실 세계의 시야를 완전히 가리고, 눈앞에 컴퓨터 그래픽으로 만들어진 100퍼센트 가상의 공간을 보여 준다. 시각과 청각을 자극하여 사용자는 마치 자신이 정말 다른 세계에 와 있는 듯한 높은 몰입감을 느끼게 된다. 예를 들어 VR 기기를 쓰면 깊은 바닷속을 탐험하거나, 직접 가 보지 못한 관광 명소를 여행하는 듯한 생생한 체험을 할 수 있다.

반면 증강현실은 우리가 보고 있는 현실 세계 위에 가상의 정보를 덧붙이는 기술이다. 주로 스마트폰이나 태블릿 PC의 카메라를 통해 작동한다. 카메라로 실제 주변 환경을 비추면, 화면에 가상의 이미지나 정보가 실제 모습과 겹쳐서 나타나는 방식이다. 현실 세계가 기본 바탕이 되기 때문에 사용자는 현실감을 유지하면서도 유용한 정보를 얻거나 재미를 느낄 수 있다. 스마트폰으로 길거리를 비추며 포켓몬을 잡는 게임이나, 가구 앱을 이용해 우리 집에 가상의 소파를 미리 놓아 보는 것이 증강현실의 대표적인 예이다.

이처럼 가상현실은 현실과 단절된 새로운 공간을 창조하고, 증강현실은 현실 세계를 확장하여 더 풍부한 경험을 제공한다는 점에서 근본적으로 다르다. 이 두 기술은 게임을 넘어 교육, 의료, 산업 등 다양한 분야에서 우리 삶을 더욱 편리하고 안전하게 만들고 있다. 앞으로 가상현실과 증강현실 기술은 더욱 발전하여 현실과 가상의 경계를 넘나드는 새로운 시대를 열어갈 것이다.
}


\kfSecHeader{kfAreaNonlit}{활동}{어휘 학습}
\noindent\textit{\small 다음 표의 '의미'를 읽고, 그에 해당하는 '어휘'를 지문에서 찾아 적으시오.}\par\medskip
\begin{kfVocab}
\kfVocabItem{1}{실제로 겪는 듯한 느낌}
\kfVocabItem{2}{컴퓨터 기술로 만든 가짜 현실}
\kfVocabItem{3}{현실 세계에 가상의 정보를 덧붙여 보여주는 기술}
\kfVocabItem{4}{기계 등이 움직이는 방법}
\kfVocabItem{5}{컴퓨터 기술을 이용한, 숫자로 표현되는 방식}
\kfVocabItem{6}{안경이나 옷 등을 몸에 걸치거나 쓰는 것}
\kfVocabItem{7}{컴퓨터를 이용해 만든 그림이나 영상}
\kfVocabItem{8}{어떤 일에 깊이 빠져드는 느낌}
\kfVocabItem{9}{경치가 좋거나 유명해서 사람들이 많이 찾는 곳}
\kfVocabItem{10}{실제 세계와의 연결이 완전히 끊어짐}
\kfVocabItem{11}{원래 있던 것을 더 넓고 풍부하게 만듦}
\kfVocabItem{12}{물건을 만들거나 자원을 개발하는 등의 산업 활동}
\kfVocabItem{13}{서로 나뉘어 있는 것들 사이의 구분}
\end{kfVocab}

\kfSecHeader{kfAreaNonlit}{활동}{비문학 문장 독해}
\noindent\textit{\small 지문을 문장 단위로 정독하면서, 문장별로 주어진 물음에 답해 보자.}\par\medskip
\begin{kfSentenceUnit}{1}{영화나 게임 속 주인공처럼 새로운 세상을 탐험하거나, 내 방에 공룡이 나타나는 모습을 상상해 본 적이 있는가?}
\kfSentQ{1.}{이 문장에서 제시된 두 가지 상상의 예는 무엇인가?}
\begin{kfChoices}
\kfChoice 영화 속 주인공이 되어 악당과 싸우는 것
\kfChoice 새로운 세상을 탐험하거나 방에 공룡이 나타나는 것
\kfChoice 하늘을 나는 자동차를 타고 우주로 떠나는 것
\kfChoice 투명 인간이 되어 사람들 몰래 장난치는 것
\end{kfChoices}
\kfSentQ{2.}{이 질문을 통해 앞으로 어떤 내용이 나올 것을 예상할 수 있는가?}
\begin{kfChoices}
\kfChoice 상상을 현실로 만드는 기술에 대한 설명
\kfChoice 영화와 게임 산업의 발전 과정에 대한 소개
\kfChoice 공룡이 멸종한 과학적인 이유에 대한 분석
\kfChoice 미래 사회의 직업 변화에 대한 구체적 전망
\end{kfChoices}
\end{kfSentenceUnit}

\begin{kfSentenceUnit}{2}{이러한 상상을 현실로 만드는 기술이 바로 가상현실(Virtual Reality, 이하 VR)과 증강현실(Augmented Reality, 이하 AR)이다.}
\kfSentQ{3.}{'이러한 상상'이 가리키는 것은 무엇인가?}
\begin{kfChoices}
\kfChoice 새로운 세상 탐험이나 방 안의 공룡 등장
\kfChoice 복권에 당첨되어 부자가 되는 행복한 꿈
\kfChoice 시험에서 백 점을 맞아 칭찬받는 상황
\kfChoice 친구들과 싸우지 않고 사이좋게 지내는 것
\end{kfChoices}
\kfSentQ{4.}{상상을 현실로 만드는 두 가지 기술은 무엇인가?}
\begin{kfChoices}
\kfChoice 인공지능(AI)과 로봇 공학
\kfChoice 가상현실(VR)과 증강현실(AR)
\kfChoice 3D 프린팅과 드론 기술
\kfChoice 사물인터넷(IoT)과 빅데이터
\end{kfChoices}
\end{kfSentenceUnit}

\begin{kfSentenceUnit}{3}{두 기술은 실감 나는 체험을 제공한다는 점에서는 같지만, 작동 방식에는 분명한 차이가 있다.}
\kfSentQ{5.}{'두 기술'이 가리키는 것은 무엇인가?}
\begin{kfChoices}
\kfChoice 스마트폰과 컴퓨터
\kfChoice 가상현실과 증강현실
\kfChoice 텔레비전과 라디오
\kfChoice 사진기와 캠코더
\end{kfChoices}
\kfSentQ{6.}{가상현실(VR)과 증강현실(AR)의 공통점은 무엇인가?}
\begin{kfChoices}
\kfChoice 사용자의 눈을 완전히 가린다는 점
\kfChoice 실감 나는 체험을 제공한다는 점
\kfChoice 현실 세계를 있는 그대로 보여준다는 점
\kfChoice 반드시 스마트폰이 필요하다는 점
\end{kfChoices}
\kfSentQ{7.}{가상현실(VR)과 증강현실(AR)의 차이점은 무엇인가?}
\begin{kfChoices}
\kfChoice 기술이 개발된 시기와 장소
\kfChoice 기기의 가격과 판매량
\kfChoice 기술의 작동 방식과 원리
\kfChoice 사용하는 전기의 양과 배터리 수명
\end{kfChoices}
\kfSentQ{8.}{이 문장을 통해 앞으로 어떤 내용이 나올 것을 예상할 수 있는가?}
\begin{kfChoices}
\kfChoice VR과 AR의 정의와 차이점
\kfChoice VR 게임의 역사와 종류
\kfChoice AR 기술의 단점과 한계
\kfChoice 최신 스마트폰의 기능 소개
\end{kfChoices}
\end{kfSentenceUnit}

\begin{kfSentenceUnit}{4}{가상현실과 증강현실은 무엇이며, 어떻게 다를까?}
\kfSentQ{9.}{이 글 전체가 답하려는 핵심 질문 두 가지는 무엇인가?}
\begin{kfChoices}
\kfChoice 누가 만들었으며, 얼마인가?
\kfChoice 언제 시작됐으며, 왜 망했나?
\kfChoice 무엇이며, 어떻게 다를까?
\kfChoice 어디서 팔며, 누가 사는가?
\end{kfChoices}
\end{kfSentenceUnit}

\begin{kfSentenceUnit}{5}{먼저 가상현실은 우리를 완전히 새로운 디지털 세상으로 안내하는 기술이다.}
\kfSentQ{10.}{가상현실은 우리를 어떤 세상으로 안내하는가?}
\begin{kfChoices}
\kfChoice 과거의 역사적 사건 현장
\kfChoice 완전히 새로운 디지털 세상
\kfChoice 조금 변형된 현실 세계
\kfChoice 미래의 최첨단 도시
\end{kfChoices}
\end{kfSentenceUnit}

\begin{kfSentenceUnit}{6}{사용자는 머리에 쓰는 특수한 안경 형태의 기기를 착용한다.}
\kfSentQ{11.}{가상현실을 체험하기 위해 주로 어떤 기기를 착용하는가?}
\begin{kfChoices}
\kfChoice 머리에 쓰는 특수 안경
\kfChoice 손목에 차는 스마트 시계
\kfChoice 귀에 꽂는 무선 이어폰
\kfChoice 손에 쥐는 게임 조종기
\end{kfChoices}
\end{kfSentenceUnit}

\begin{kfSentenceUnit}{7}{이 기기는 현실 세계의 시야를 완전히 가리고, 눈앞에 컴퓨터 그래픽으로 만들어진 100퍼센트 가상의 공간을 보여 준다.}
\kfSentQ{12.}{'이 기기'가 가리키는 것은 무엇인가?}
\begin{kfChoices}
\kfChoice 가상현실 체험용 안경 기기
\kfChoice 증강현실 구현용 스마트폰
\kfChoice 소리를 전달하는 무선 헤드셋
\kfChoice 화면을 조작하는 게임 컨트롤러
\end{kfChoices}
\kfSentQ{13.}{가상현실 기기가 보여주는 공간의 특징 두 가지는 무엇인가?}
\begin{kfChoices}
\kfChoice 현실을 가리고 가상 공간을 보여줌
\kfChoice 현실 위에 가상의 정보를 덧붙임
\kfChoice 현실의 소리를 증폭시켜 들려줌
\kfChoice 가상의 냄새를 맡게 해줌
\end{kfChoices}
\kfSentQ{14.}{가상의 공간은 무엇으로 만들어진 것인가?}
\begin{kfChoices}
\kfChoice 붓으로 그린 그림
\kfChoice 카메라로 찍은 사진
\kfChoice 컴퓨터 그래픽 이미지
\kfChoice 실제 사물의 홀로그램
\end{kfChoices}
\end{kfSentenceUnit}

\begin{kfSentenceUnit}{8}{시각과 청각을 자극하여 사용자는 마치 자신이 정말 다른 세계에 와 있는 듯한 높은 몰입감을 느끼게 된다.}
\kfSentQ{15.}{가상현실이 자극하는 감각은 무엇인가?}
\begin{kfChoices}
\kfChoice 맛과 냄새를 느끼는 미각과 후각
\kfChoice 보고 듣는 것을 담당하는 시각과 청각
\kfChoice 피부로 느끼는 온도와 촉각
\kfChoice 몸의 균형을 잡는 평형 감각
\end{kfChoices}
\kfSentQ{16.}{사용자가 느끼는 감정은 무엇인가?}
\begin{kfChoices}
\kfChoice 마치 진짜 같은 높은 몰입감
\kfChoice 현실과 동떨어진 심한 이질감
\kfChoice 장시간 사용에 따른 어지러움
\kfChoice 기계 작동에 대한 어려움
\end{kfChoices}
\kfSentQ{17.}{사용자는 어떤 느낌을 받는가?}
\begin{kfChoices}
\kfChoice 마치 다른 세계에 와 있는 듯한 느낌
\kfChoice 영화관에서 영화를 보고 있는 느낌
\kfChoice 꿈속에서 하늘을 날고 있는 느낌
\kfChoice 친구와 영상 통화를 하고 있는 느낌
\end{kfChoices}
\kfSentQ{18.}{가상현실이 높은 몰입감을 주는 이유는 무엇인가?}
\begin{kfChoices}
\kfChoice 시각과 청각을 동시에 자극하기 때문
\kfChoice 실제 물건을 만질 수 있기 때문
\kfChoice 맛있는 음식을 먹을 수 있기 때문
\kfChoice 다른 사람과 대화할 수 있기 때문
\end{kfChoices}
\end{kfSentenceUnit}

\begin{kfSentenceUnit}{9}{예를 들어 VR 기기를 쓰면 깊은 바닷속을 탐험하거나, 직접 가 보지 못한 관광 명소를 여행하는 듯한 생생한 체험을 할 수。}
\kfSentQ{19.}{가상현실을 통해 할 수 있는 체험의 예시 두 가지는 무엇인가?}
\begin{kfChoices}
\kfChoice 깊은 바닷속 탐험이나 가상 여행
\kfChoice 친구들과 함께하는 축구 경기
\kfChoice 집에서 혼자 하는 요리 실습
\kfChoice 학교에서 배우는 교과서 공부
\end{kfChoices}
\end{kfSentenceUnit}

\begin{kfSentenceUnit}{10}{반면 증강현실은 우리가 보고 있는 현실 세계 위에 가상의 정보를 덧붙이는 기술이다.}
\kfSentQ{20.}{'반면'이 나타내는 관계는 무엇인가?}
\begin{kfChoices}
\kfChoice 앞의 내용과 반대되거나 대조됨
\kfChoice 앞의 내용에 대한 원인을 설명함
\kfChoice 앞의 내용에 대한 결과를 설명함
\kfChoice 앞의 내용에 대한 예시를 듦
\end{kfChoices}
\kfSentQ{21.}{증강현실은 어디에 무엇을 덧붙이는 기술인가?}
\begin{kfChoices}
\kfChoice 가상 공간에 현실의 물건을 덧붙이는 기술
\kfChoice 현실 세계에 가상의 정보를 덧붙이는 기술
\kfChoice 현실 세계에 가상의 소리를 제거하는 기술
\kfChoice 가상 공간에 현실의 날씨를 반영하는 기술
\end{kfChoices}
\kfSentQ{22.}{가상현실과 증강현실의 차이점을 이 문장을 통해 설명하면?}
\begin{kfChoices}
\kfChoice 가상현실은 현실을 확장하고, 증강현실은 가상을 창조한다
\kfChoice 가상현실은 새로운 세상을 만들고, 증강현실은 정보를 덧붙인다
\kfChoice 가상현실은 스마트폰을 쓰고, 증강현실은 안경을 쓴다
\kfChoice 가상현실은 과거를 보여주고, 증강현실은 미래를 보여준다
\end{kfChoices}
\end{kfSentenceUnit}

\begin{kfSentenceUnit}{11}{주로 스마트폰이나 태블릿 PC의 카메라를 통해 작동한다.}
\kfSentQ{23.}{증강현실은 주로 어떤 기기를 통해 작동하는가?}
\begin{kfChoices}
\kfChoice 스마트폰이나 태블릿 PC의 카메라
\kfChoice 머리에 착용하는 특수 안경 기기
\kfChoice 손목에 차는 스마트 워치와 밴드
\kfChoice 대형 스크린이 있는 컴퓨터 모니터
\end{kfChoices}
\kfSentQ{24.}{VR과 AR이 사용하는 기기의 차이점은 무엇인가?}
\begin{kfChoices}
\kfChoice VR은 특수 안경, AR은 스마트폰 카메라
\kfChoice VR은 스마트폰, AR은 특수 안경 기기
\kfChoice VR은 컴퓨터, AR은 텔레비전 화면
\kfChoice VR은 게임기, AR은 라디오 수신기
\end{kfChoices}
\end{kfSentenceUnit}

\begin{kfSentenceUnit}{12}{카메라로 실제 주변 환경을 비추면, 화면에 가상의 이미지나 정보가 실제 모습과 겹쳐서 나타나는 방식이다.}
\kfSentQ{25.}{AR이 작동하는 첫 번째 단계는 무엇인가?}
\begin{kfChoices}
\kfChoice 가상의 이미지를 먼저 선택하는 것
\kfChoice 카메라로 실제 주변 환경을 비추는 것
\kfChoice 스마트폰의 전원을 끄고 다시 켜는 것
\kfChoice 특수 안경을 머리에 착용하는 것
\end{kfChoices}
\kfSentQ{26.}{AR이 작동하는 두 번째 단계는 무엇인가?}
\begin{kfChoices}
\kfChoice 현실의 모습을 완전히 지우는 것
\kfChoice 가상의 이미지가 실제와 겹쳐 보이는 것
\kfChoice 사용자의 위치를 지도에 표시하는 것
\kfChoice 새로운 가상 공간으로 이동하는 것
\end{kfChoices}
\kfSentQ{27.}{증강현실은 어떤 방식으로 가상 정보를 보여주는가?}
\begin{kfChoices}
\kfChoice 실제 모습과 가상 정보가 겹쳐서 나타남
\kfChoice 실제 모습은 사라지고 가상 정보만 나타남
\kfChoice 가상 정보는 소리로만 들려주고 화면은 꺼짐
\kfChoice 실제 모습 옆에 가상 정보를 따로 보여줌
\end{kfChoices}
\end{kfSentenceUnit}

\begin{kfSentenceUnit}{13}{현실 세계가 기본 바탕이 되기 때문에 사용자는 현실감을 유지하면서도 유용한 정보를 얻거나 재미를 느낄 수 있다.}
\kfSentQ{28.}{AR에서 기본 바탕이 되는 것은 무엇인가?}
\begin{kfChoices}
\kfChoice 컴퓨터가 만든 가상 공간
\kfChoice 우리가 보고 있는 현실 세계
\kfChoice 캄캄한 어둠 속의 빈 공간
\kfChoice 만화 영화 속의 그림 배경
\end{kfChoices}
\kfSentQ{29.}{AR 사용자가 얻을 수 있는 두 가지 효과는 무엇인가?}
\begin{kfChoices}
\kfChoice 유용한 정보를 얻거나 재미를 느낌
\kfChoice 깊은 수면을 취하거나 휴식을 취함
\kfChoice 시력이 좋아지거나 청력이 발달함
\kfChoice 친구를 많이 사귀거나 인기가 많아짐
\end{kfChoices}
\kfSentQ{30.}{AR의 장점은 무엇인가?}
\begin{kfChoices}
\kfChoice 현실과 완전히 분리되어 몰입감을 줌
\kfChoice 현실감을 유지하며 정보와 재미를 줌
\kfChoice 복잡한 장비 없이 맨눈으로 볼 수 있음
\kfChoice 전기가 없어도 언제든 사용할 수 있음
\end{kfChoices}
\kfSentQ{31.}{AR이 현실감을 유지할 수 있는 이유는 무엇인가?}
\begin{kfChoices}
\kfChoice 현실 세계를 완전히 가리기 때문이다
\kfChoice 현실 세계가 기본 바탕이 되기 때문이다
\kfChoice 가상의 소리를 들려주기 때문이다
\kfChoice 특수 안경을 착용하기 때문이다
\end{kfChoices}
\end{kfSentenceUnit}

\begin{kfSentenceUnit}{14}{스마트폰으로 길거리를 비추며 포켓몬을 잡는 게임이나, 가구 앱을 이용해 우리 집에 가상의 소파를 미리 놓아 보는 것이 증강현실의 대표적인 예이다.}
\kfSentQ{32.}{증강현실의 대표적인 예시 두 가지는 무엇인가?}
\begin{kfChoices}
\kfChoice 심해 탐험과 우주 여행
\kfChoice 포켓몬 게임과 가구 배치 앱
\kfChoice 드론 조종과 로봇 코딩
\kfChoice 영화 감상과 음악 감상
\end{kfChoices}
\end{kfSentenceUnit}

\begin{kfSentenceUnit}{15}{이처럼 가상현실은 현실과 단절된 새로운 공간을 창조하고, 증강현실은 현실 세계를 확장하여 더 풍부한 경험을 제공한다는 점에서 근본적으로 다르다.}
\kfSentQ{33.}{VR의 특징은 무엇인가?}
\begin{kfChoices}
\kfChoice 현실을 확장하여 풍부한 경험 제공
\kfChoice 현실과 단절된 새로운 공간 창조
\kfChoice 현실 위에 가상 정보 덧입히기
\kfChoice 스마트폰 카메라로 주변 비추기
\end{kfChoices}
\kfSentQ{34.}{AR의 특징은 무엇인가?}
\begin{kfChoices}
\kfChoice 현실과 단절된 새로운 공간 창조
\kfChoice 현실 세계를 확장하여 경험 제공
\kfChoice 시야를 완전히 가리는 몰입감
\kfChoice 가상의 공간으로 이동하는 느낌
\end{kfChoices}
\kfSentQ{35.}{가상현실과 증강현실의 근본적인 차이점은 무엇인가?}
\begin{kfChoices}
\kfChoice 사용되는 전기의 양과 배터리 시간
\kfChoice 기기의 가격과 디자인의 차이
\kfChoice 현실과의 단절 여부와 공간 창조 방식
\kfChoice 개발된 국가와 제조 회사의 차이
\end{kfChoices}
\end{kfSentenceUnit}

\begin{kfSentenceUnit}{16}{이 두 기술은 게임을 넘어 교육, 의료, 산업 등 다양한 분야에서 우리 삶을 더욱 편리하고 안전하게 만들고 있다.}
\kfSentQ{36.}{'이 두 기술'이 가리키는 것은 무엇인가?}
\begin{kfChoices}
\kfChoice 인공지능과 빅데이터
\kfChoice 자율주행과 드론
\kfChoice 가상현실과 증강현실
\kfChoice 3D 프린터와 로봇
\end{kfChoices}
\kfSentQ{37.}{두 기술이 활용되는 분야의 예시 세 가지는 무엇인가?}
\begin{kfChoices}
\kfChoice 교육, 의료, 산업 분야
\kfChoice 농업, 어업, 임업 분야
\kfChoice 정치, 경제, 사회 분야
\kfChoice 요리, 청소, 세탁 분야
\end{kfChoices}
\end{kfSentenceUnit}

\begin{kfSentenceUnit}{17}{앞으로 가상현실과 증강현실 기술은 더욱 발전하여 현실과 가상의 경계를 넘나드는 새로운 시대를 열어갈 것이다.}
\kfSentQ{38.}{이 문장은 VR과 AR의 무엇에 대해 말하고 있는가?}
\begin{kfChoices}
\kfChoice 과거의 역사
\kfChoice 현재의 문제점
\kfChoice 미래의 전망
\kfChoice 기술적 한계
\end{kfChoices}
\kfSentQ{39.}{이 글이 예측하는 두 기술의 미래 모습은 어떠한가?}
\begin{kfChoices}
\kfChoice 기술적 한계로 인해 사라질 것이다
\kfChoice 특정 전문가들만 사용하게 될 것이다
\kfChoice 현실과 가상의 경계를 넘나들 것이다
\kfChoice 게임 분야에서만 제한적으로 쓰일 것이다
\end{kfChoices}
\end{kfSentenceUnit}

\kfSecHeader{kfAreaNonlit}{활동}{구조도}
\noindent\textit{\small 글의 흐름을 살펴보고, 구조도의 빈칸을 알맞게 채워 보자.}\par\medskip
\par\medskip\needspace{6\baselineskip}
\begin{tcolorbox}[enhanced, breakable=false,
  colback=kfPaper, colframe=kfRule, boxrule=0.4pt, arc=2pt,
  left=10pt, right=10pt, top=8pt, bottom=8pt,
  title={\footnotesize\bfseries\color{kfMute} 글의 구조도},
  coltitle=kfMute, colbacktitle=kfPrimaryLight!40,
  attach boxed title to top left={yshift=-1.5mm, xshift=6pt},
  boxed title style={colframe=kfRule, boxrule=0.3pt, arc=1pt}]
\setstretch{1.4}\footnotesize
지침: 글의 전체적인 논리 흐름을 파악하며 빈칸에 알맞은 내용을 채우시오. \\

\noindent \textbullet\ 1문단: 주제 소개 \\
\noindent \quad\textbullet\ 공통점: 상상을 현실로 만드는 ( ① ) 체험 기술 \\
\noindent \quad\textbullet\ 차이점: ( ② ) 방식의 차이 \\
\noindent \quad\textbullet\ 질문: 두 기술은 무엇이며 어떻게 다른가? \\
\noindent \textbullet\ 2문단: ( ③ ) (VR) \\
\noindent \quad\textbullet\ 개념: 완전히 ( ④ )된 새로운 디지털 세상 \\
\noindent \quad\textbullet\ 체험 방식: 머리에 쓰는 기기로 시야를 완전히 ( ⑤ ) \\
\noindent \quad\textbullet\ 특징: 높은 ( ⑥ ) \\
\noindent \quad\textbullet\ 예시: 바닷속 탐험, 가상 여행 \\
\noindent \textbullet\ 3문단: ( ⑦ ) (AR) \\
\noindent \quad\textbullet\ 개념: ( ⑧ ) 세계 위에 가상의 정보를 ( ⑨ ) 기술 \\
\noindent \quad\textbullet\ 체험 방식: 스마트폰 ( ⑩ )로 실제 환경을 비춤 \\
\noindent \quad\textbullet\ 특징: ( ⑪ )을 유지하며 정보 획득 \\
\noindent \quad\textbullet\ 예시: 포켓몬 잡는 게임, 가구 배치 앱 \\
\noindent \textbullet\ 4문단: 가상 현실과 증강 현실의 차이점과 활용 및 전망 \\
\noindent \quad\textbullet\ 핵심 차이점 \\
\noindent \quad\quad\textbullet\ 가상현실: 현실과 ( ⑫ )된 새로운 공간 창조 \\
\noindent \quad\quad\textbullet\ 증강현실: 현실 세계를 ( ⑬ ) \\
\noindent \quad\textbullet\ 미래 전망: 게임을 넘어 교육, 의료 등 다양한 분야로 발전 \\
\end{tcolorbox}\par\medskip
\kfSecHeader{kfAreaNonlit}{문제}{}

\begin{kfQBlock}{1}{}
\kfQStem{윗글의 내용과 일치하지 \uline{않는} 것은?}
\begin{kfChoices}
\kfChoice 가상현실과 증강현실은 모두 실감 나는 체험을 제공한다.
\kfChoice 가상현실은 주로 스마트폰 카메라를 통해 체험할 수 있다.
\kfChoice 증강현실은 현실 세계를 바탕으로 가상 정보를 보여준다.
\kfChoice 가상현실은 사용자에게 높은 몰입감을 주는 것을 목표로 한다.
\kfChoice 증강현실은 현실감을 유지하면서 추가적인 정보를 얻을 수 있다.
\end{kfChoices}
\end{kfQBlock}
\begin{kfQBlock}{2}{}
\kfQStem{가상현실(VR)과 증강현실(AR)의 차이점에 대한 설명으로 적절하지 \uline{않은} 것은?}
\begin{kfChoices}
\kfChoice VR은 현실과 단절된 가상 공간을, AR은 현실이 확장된 공간을 만든다.
\kfChoice VR은 최신 기술이고 AR은 그보다 훨씬 오래된 기술이다.
\kfChoice AR은 스마트폰 같은 기기로 현실 위에 가상 정보를 덧입힌다.
\kfChoice VR은 몰입감 중심의 체험을, AR은 현실 위에 정보를 더하는 체험을 제공한다.
\kfChoice VR은 머리에 쓰는 기기를 통해 가상 공간을 체험하게 한다.
\end{kfChoices}
\end{kfQBlock}
\begin{kfQBlock}{3}{}
\kfQStem{<보기>의 상황에 사용된 기술은 무엇인가?}
\begin{kfEvidence}나는 박물관에서 공룡 뼈 화석을 스마트폰 카메라로 비추었다. 그랬더니 화면 속 뼈 위에 살이 붙고 피부가 생기면서, 살아있는 공룡이 눈앞에서 움직이는 것처럼 보였다.\end{kfEvidence}
\begin{kfChoices}
\kfChoice 가상현실(VR)
\kfChoice 증강현실(AR)
\kfChoice 인공지능(AI)
\kfChoice 동기화 기술
\kfChoice 홀로그램
\end{kfChoices}
\end{kfQBlock}
\begin{kfQBlock}{4}{}
\kfQStem{다음 중 ‘가상현실(VR)’ 기술의 예시로 적절하지 \uline{않은} 것은?}
\begin{kfChoices}
\kfChoice 스마트폰 카메라로 길거리를 비추며 포켓몬을 잡는 게임
\kfChoice HMD를 쓰고 우주 공간을 떠다니는 시뮬레이션을 체험하는 것
\kfChoice VR 기기를 착용하고 가상 박물관을 둘러보는 것
\kfChoice VR 안경을 쓰고 공포 게임 속 세상을 직접 경험하는 것
\kfChoice 특수한 안경을 쓰고 깊은 바닷속을 탐험하는 체험
\end{kfChoices}
\end{kfQBlock}
\begin{kfQBlock}{5}{}
\kfQStem{윗글을 통해 알 수 있는 내용으로 적절하지 \uline{않은} 것은?}
\begin{kfChoices}
\kfChoice 가상현실과 증강현실은 작동 원리가 같은 기술이어서 굳이 구분할 필요가 없다.
\kfChoice 가상현실은 현실 시야를 가린 가상 공간에서 높은 몰입감을 느끼게 한다.
\kfChoice 증강현실은 현실 위에 가상 정보를 덧붙여 정보를 풍부하게 한다.
\kfChoice 두 기술은 게임뿐 아니라 교육·산업 등 다양한 분야에 활용된다.
\kfChoice 가상현실과 증강현실은 현실 세계를 활용하는 방식이 서로 다르다.
\end{kfChoices}
\end{kfQBlock}
\begin{kfQBlock}{1}{}
\kfQStem{우리를 완전히 새로운 디지털 세상으로 안내하며, 머리에 쓰는 기기를 통해 높은 몰입감을 주는 기술은 무엇인가?}
\kfAnswerLines{1}
\end{kfQBlock}
\begin{kfQBlock}{2}{}
\kfQStem{현실 세계 위에 가상의 정보를 덧붙여 보여주며, 주로 스마트폰 카메라를 통해 작동하는 기술은 무엇인가?}
\kfAnswerLines{1}
\end{kfQBlock}
\begin{kfQBlock}{3}{}
\kfQStem{가상현실과 증강현실이 공통적으로 제공하는 것은 무엇인가?}
\kfAnswerLines{1}
\end{kfQBlock}
\begin{kfQBlock}{4}{}
\kfQStem{증강현실 기술이 현실감을 유지할 수 있는 이유는 무엇이 기본 바탕이 되기 때문인가?}
\kfAnswerLines{1}
\end{kfQBlock}
\begin{kfQBlock}{5}{}
\kfQStem{가구 앱을 이용해 우리 집에 가상의 소파를 미리 놓아 보는 것은 어떤 기술의 예시인가?}
\kfAnswerLines{1}
\end{kfQBlock}
\begin{kfQBlock}{1}{}
\kfQStem{'가상현실'과 '증강현실'의 차이점을 다음 조건에 맞게 서술하시오.}
\kfAnswerNote{답안}{4}
\end{kfQBlock}



\end{document}
